 
% new file names from date 20171119: 
% block2_171119.tex 

\renewcommand{\thisfile}{block2\_190328}  
  
\cleardoublepage
\mysection{Triangulated Categories and Functors} 
\label{sec:triangulated}  

\AYcopyright

In this section we introduce triangulated categories and triangulated 
functors. We prove that for a DG ring $A$ and an abelian category  
$\cat{M}$, the homotopy category $\dcat{K}(A, \cat{M})$
is triangulated. We also prove that a DG functor 
$F : \dcat{C}(A, \cat{M}) \to \dcat{C}(B, \cat{N})$
induces a triangulated functor 
$F : \dcat{K}(A, \cat{M}) \to \dcat{K}(B, \cat{N})$.

Recall that by Convention \ref{conv:2490}, there is a nonzero 
commutative base ring $\K$, which is implicit most of the time. All rings are 
central $\K$-rings, and all ring homomorphisms are over $\K$. All linear 
categories are $\K$-linear, and all 
linear (i.e.\ additive) functors between them are $\K$-linear. 



\mysubsection{Triangulated Categories}

In this subsection we define triangulated categories, and make some remarks 
regarding them. 
  
Recall that a functor is called an isomorphism of categories if it is bijective 
on sets of objects and on sets of morphisms. 

\begin{dfn} \label{dfn:1155} 
\index{Translation functor! of an additive category}
\index{T-additive! category}
Let $\cat{K}$ be an additive category. A {\em translation functor} on $\cat{K}$ 
is an additive automorphism $\opn{T}$ of $\cat{K}$.
The pair $(\cat{K}, \opn{T})$ is called a {\em T-additive category}. 
\end{dfn}

\begin{rem}
Some texts give a more relaxed definition: the translation functor $\opn{T}$ is 
only required to be an 
additive auto-equivalence of $\cat{K}$. The resulting theory is more 
complicated (it is $2$-categorical, but most texts try to suppress this fact). 
\end{rem}

Later in the book (see Definition \ref{dfn:2120}) we will write 
$M[k] := \opn{T}^k(M)$, the $k$-th translation of an object $M$. 

\begin{dfn} \label{dfn:1159} 
\index{T-additive! functor}
\index{Translation isomorphism! of an additive functor}
Suppose $(\cat{K}, \opn{T}_{\cat{K}})$ and 
$(\cat{L}, \opn{T}_{\cat{L}})$ are T-additive categories.
A {\em T-additive functor} between them is a pair 
$(F, \tau)$, consisting of an additive functor $F : \cat{K} \to \cat{L}$, 
together with an isomorphism
\[ \tau :  F \circ \opn{T}_{\cat{K}}  \iso 
\opn{T}_{\cat{L}} \circ \, F   \]
of functors $\cat{K} \to \cat{L}$,
called a {\em translation isomorphism}. 
\end{dfn}

\begin{dfn} \label{dfn:1840}
Let $(\cat{K}_i, \opn{T}_i)$ be T-additive categories, for $i = 0, 1, 2$, and 
let 
\[ (F_i, \tau_i) : (\cat{K}_{i - 1}, \opn{T}_{i - 1}) \to 
(\cat{K}_i, \opn{T}_i) \]
be T-additive functors. The composition 
\[ (F, \tau) = (F_2, \tau_2) \circ (F_1, \tau_1) \]
is the T-additive functor 
$(\cat{K}_0, \opn{T}_0) \to (\cat{K}_2, \opn{T}_2)$
defined as follows: the functor is
$F := F_2 \circ F_1$, and the translation isomorphism 
$\tau : F \circ T_0 \iso T_2 \circ F$ 
is 
$\tau := \tau_2 \circ F_2(\tau_1)$. 
\end{dfn}

\begin{dfn} \label{dfn:1160} 
Suppose $(\cat{K}, \opn{T}_{\cat{K}})$ and 
$(\cat{L}, \opn{T}_{\cat{L}})$ are T-additive categories,
and 
\[ (F, \tau), (G, \nu) : (\cat{K}, \opn{T}_{\cat{K}}) \to 
(\cat{L}, \opn{T}_{\cat{L}}) \] 
are T-additive functors. A {\em morphism of T-additive functors}
\index{T-additive! morphism of {\indash} functors}
\[ \eta : (F, \tau) \to (G, \nu) \] 
is a morphism of functors $\eta : F \to G$, such that for every object 
$M \in \cat{K}$ this diagram in $\cat{L}$ is commutative:
\[ \UseTips  \xymatrix @C=8ex @R=6ex {
F(\opn{T}_{\cat{K}}(M))
\ar[r]^{\tau_M}
\ar[d]_{\eta_{\lms \opn{T}_{\cat{K}}(M)}}
&
\opn{T}_{\cat{L}}(F(M))
\ar[d]^{\opn{T}_{\cat{L}}(\eta_M)}
\\
G(\opn{T}_{\cat{K}}(M))
\ar[r]^{\nu_M}
&
\opn{T}_{\cat{L}}(G(M)) \ . 
} \]
\end{dfn}

\begin{dfn} \label{dfn:1156} 
Let $(\cat{K}, \opn{T})$ be a T-additive category. 
A {\em triangle} 
\index{Triangle! in a T-additive category}
in $(\cat{K}, \opn{T})$ is a diagram 
\[ L \xar{\al} M \xar{\be} N \xar{\ga} \opn{T}(L)  \]
in $\cat{K}$.
\end{dfn}

Sometimes, when the names of the morphisms in the triangle above are not 
important, we may use the more compact notation 
\begin{equation} \label{eqn:4910}
L \xar{} M \xar{} N \xar{ \, \triangle \, } .
\end{equation}

\begin{dfn} \label{dfn:1157} 
Let $(\cat{K}, \opn{T})$ be a T-additive category. Suppose 
\[  L \xar{\al} M \xar{\be} N \xar{\ga} \opn{T}(L) \]
and 
\[ L' \xar{\al'} M' \xar{\be'} N' \xar{\ga'} \opn{T}(L') \]
are triangles in  $(\cat{K}, \opn{T})$. 
A {\em morphism of triangles}
\index{Triangle! morphism of {\indash}s}
between them is a commutative diagram 
\[ \UseTips  \xymatrix @C=8ex @R=6ex { 
L 
\ar[r]^{\al}
\ar[d]_{\phi}
&
M
\ar[r]^{\be}
\ar[d]_{\psi}
& 
N
\ar[r]^{\ga}
\ar[d]_{\chi}
&
\opn{T}(L)
\ar[d]_{\opn{T}(\phi)}
\\
L' 
\ar[r]^{\al'}
&
M'
\ar[r]^{\be'}
& 
N'
\ar[r]^{\ga'}
&
\opn{T}(L') 
} \]
in $\cat{K}$. 

The morphism of triangles $(\phi, \psi, \chi)$ is called an isomorphism if 
$\phi, \psi$ and $\chi$ are all isomorphisms. 
\end{dfn}

\begin{rem} \label{rem:4690}
Why ``triangle''? This is because sometimes a triangle 
\[ L \xar{\al} M \xar{\be} N \xar{\ga} \opn{T}(L) \]
is written as a diagram
\[ \UseTips  \xymatrix @C=4ex @R=4ex {
& 
N 
\ar[dl]_{\ga}
\\
L
\ar[rr]_{\al}
& &
M
\ar[lu]_{\be}
} \]
See \cite[page 20]{RD}. 
\end{rem}

\begin{dfn} \label{dfn:1158} 
A {\em triangulated category}
\index{Triangulated category}
is a T-additive category
$(\cat{K}, \opn{T})$, \lb equipped with a set of triangles called  {\em 
distinguished triangles}. 
\index{Triangle! distinguished}
The following axioms have to be satisfied:
\begin{enumerate}[leftmargin=30pt]
\item[(TR1)] 
\begin{itemize}
\rmitem{a} Every triangle that is isomorphic to a distinguished
triangle is also a distinguished triangle. 

\rmitem{b} For every morphism $\al : L \to M$ in $\cat{K}$ there is a
distinguished triangle 
\[ L \xar{\al} M \xar{} N \xar{} \opn{T}(L) . \] 

\rmitem{c}  For every object $M$ the triangle 
\[ M \xar{\opn{id}_M} M \to 0 \to \opn{T}(M) \] 
is distinguished.
\end{itemize}

\item[(TR2)] A triangle 
\[ L \xar{\al} M \xar{\be} N \xar{\ga} \opn{T}(L) \]
is distinguished iff the triangle 
\[ M \xar{\be} N \xar{\ga} \opn{T}(L) \xar{- \opn{T}(\al)} \opn{T}(M) \] 
is distinguished.

\item[(TR3)]  Suppose 
\[ \UseTips  \xymatrix @C=8ex @R=6ex { 
L 
\ar[r]^{\al}
\ar[d]_{\phi}
&
M
\ar[r]^{\be}
\ar[d]_{\psi}
& 
N
\ar[r]^{\ga}
&
\opn{T}(L)
\\
L' 
\ar[r]^{\al'}
&
M'
\ar[r]^{\be'}
& 
N'
\ar[r]^{\ga'}
&
\opn{T}(L') 
} \]
is a commutative diagram in $\cat{K}$ in which the rows are distinguished 
triangles. 
Then there exists a morphism $\chi : N \to N'$ such that the diagram
\[ \UseTips  \xymatrix @C=8ex @R=6ex { 
L 
\ar[r]^{\al}
\ar[d]_{\phi}
&
M
\ar[r]^{\be}
\ar[d]_{\psi}
& 
N
\ar[r]^{\ga}
\ar[d]_{\chi}
&
\opn{T}(L)
\ar[d]_{\opn{T}(\phi)}
\\
L' 
\ar[r]^{\al'}
&
M'
\ar[r]^{\be'}
& 
N'
\ar[r]^{\ga'}
&
\opn{T}(L') \ .
} \]
is a morphism of triangles.

\item[(TR4)] Suppose we are given these three distinguished triangles:
\[ L \xar{\al} M \xar{\ga} P \xar{} \opn{T}(L) , \]
\[ M \xar{\be} N \xar{\ep} R \xar{} \opn{T}(M) , \]
\[ L \xar{\be \circ \al} N \xar{\de} Q \xar{} \opn{T}(L) . \]
Then there is a distinguished triangle 
\[  P \xar{\phi} Q \xar{\psi} R \xar{\rho} \opn{T}(P)  \]
making diagram (\ref{eqn:5022}) commutative. 
\begin{figure}[!htb]
\centering
\begin{equation} \label{eqn:5022}
\UseTips  \xymatrix @C=8ex @R=6ex { 
L 
\ar[r]^{\al}
\ar[d]_{\opn{id}}
&
M
\ar[r]^{\ga}
\ar[d]_{\be}
& 
P
\ar[r]^{}
\ar[d]_{\phi}
&
\opn{T}(L)
\ar[d]_{\opn{id}}
\\
L
\ar[r]^{\be \circ \al}
\ar[d]_{\al}
&
N
\ar[r]^{\de}
\ar[d]_{\opn{id}}
& 
Q
\ar[r]^{}
\ar[d]_{\psi}
&
\opn{T}(L)
\ar[d]_{\opn{T}(\al)}
\\
M
\ar[r]^{\be}
\ar[d]_{\ga}
&
N
\ar[r]^{\ep}
\ar[d]_{\de}
& 
R
\ar[r]^{}
\ar[d]_{\opn{id}}
&
\opn{T}(M)
\ar[d]_{\opn{T}(\ga)}
\\
P
\ar[r]^{\phi}
&
Q
\ar[r]^{\psi}
& 
R
\ar[r]^{\rho}
&
\opn{T}(P)
} 
\end{equation}
\end{figure}
\end{enumerate}
\end{dfn}

Here are a few remarks on triangulated categories. See Remark \ref{exer:3015} 
regarding triangulated categories arising from algebraic topology. 

\begin{rem} \label{rem:1280}
The numbering of the axioms we use is taken from \cite{RD}, and it agrees with 
the numbering in \cite{We}. The numbering in \cite{KaSc1} and \cite{KaSc2}  
is different.

In the situation that we care about, namely 
$\cat{K} = \dcat{K}(A, \cat{M})$,
the distinguished triangles will be those triangles that are isomorphic, in 
$\dcat{K}(A, \cat{M})$, to the standard triangles in $\dcat{C}(A, \cat{M})$
from Definition \ref{dfn:1140}. 
See Definition \ref{dfn:140} below for the precise statement. 

The object $N$ in item (b) of axiom (TR1) is referred to as a {\em cone} of the 
morphism $\al$. We should think of the cone as something combining 
``the cokernel'' and ``the kernel'' of $\al$. 
Axiom (TR3) guarantees a sort of weak functoriality of the cone. 
See Subsection \ref{subsec:props-tr-cats}, and in particular 
Remark \ref{rem:4160}, for more properties of the cone.

Axiom (TR2) says that if we ``turn'' a distinguished triangle (cf.\ Remark 
\ref{rem:4690}) we remain with a distinguished triangle.
\end{rem}

\begin{rem} \label{rem:1281}
The axiom (TR4) is called the {\em octahedral axiom}.
\index{Octahedral axiom}
It is supposed to replace the isomorphism  
\[ (N / L) / (M / L) \cong N / M \]
for objects $L \sub M \sub N$ in an abelian category $\cat{M}$. 
The octahedral axiom is needed for the theory of {\em t-structures}: it is 
used, in \cite{BBD}, to show that the heart of a t-structure is an abelian 
category (see Remark \ref{rem:3600}). This axiom is also needed to form 
Verdier quotients of triangulated categories (see Remark \ref{rem:4220}).
The book \cite{Ne1} has detailed discussions of the octahedral axiom. 

A T-additive category $(\cat{K}, \opn{T})$ that only satisfies axioms 
(TR1)-(TR3) is called a {\em pretriangulated category}. 
(The reader should not confuse ``pretriangulated category'', as used here, with 
the ``pretriangulated DG category'' from \cite{BoKa}; see Remark 
\ref{rem:1475}.) It is not known whether the octahedral axiom is a consequence 
of the other axioms; there was a recent paper by A. Maccioca 
(arxiv:1506.00887) claiming that, but it had a fatal error in it.

In our book the octahedral axiom does not play any role. For this reason we had 
excluded it from an earlier version of the book, in which we had discussed 
pretriangulated categories only. Our decision to include this axiom in the 
current version of the book, and thus to talk about triangulated categories 
(rather than about pretriangulated ones) is just to be more in line with the 
mainstream usage. With the exception of a longer proof of Theorem 
\ref{thm:1255} -- stating that $\dcat{K}(A, \cat{M})$ is a triangulated 
category -- there is virtually no change in the content of the book, and almost 
all definitions and results are valid for pretriangulated categories.  
\end{rem}

\begin{rem} \label{rem:2305}
The structure of triangulated categories is poorly understood, and there is no
classification of triangulated functors between them. This is true even 
for the derived category $\dcat{D}(A)$ of a ring $A$. For instance, a famous 
open question of J. Rickard is whether every triangulated auto-equivalence $F$ 
of $\dcat{D}(A)$ is isomorphic to $P \ot^{\mrm{L}}_{A} (-)$ for some tilting 
complex $P$. More on this issue in Section \ref{sec:perf-tilt-NC}. 

In this book the role of the triangulated structure on derived categories is 
secondary. It is mostly used for induction on the amplitude of the cohomology 
of complexes; cf.\ Section \ref{sec:adj-equ-cohdim}. 

Furthermore, one of the most important functors studied here -- the squaring 
operation $\opn{Sq}_{B / A}$ from Subsections \ref{subsec:squaring} and
\ref{subsec:RNCDC-uniq} -- is a functor from $\dcat{D}(B)$ to itself, which is 
not triangulated, nor even linear: it is a {\em quadratic functor}. 
\end{rem}




\mysubsection{Triangulated and Cohomological Functors} 
\label{subsec:tr-coh-funcs}

Suppose $\cat{K}$ and $\cat{L}$ are T-additive categories, with 
translation functors
$\opn{T}_{\cat{K}}$ and $\opn{T}_{\cat{L}}$ respectively. 
A T-additive functor $F : \cat{K} \to \cat{L}$ was defined in Definition 
\ref{dfn:1159}. A morphism of T-additive functors $\eta : F \to G$ was 
introduced in Definition \ref{dfn:1160}.

\begin{dfn} \label{dfn:1276}
Let  $\cat{K}$ and $\cat{L}$  be triangulated categories. 
\begin{enumerate}
\item A {\em triangulated functor} 
\index{Triangulated functor}
from $\cat{K}$ to $\cat{L}$
is a T-additive functor 
$(F, \tau) : \cat{K} \to \cat{L}$ 
that satisfies this condition: for every distinguished triangle 
\[ L \xar{\al} M \xar{\be} N \xar{\ga} \opn{T}_{\cat{K}}(L) \]
in $\cat{K}$, the triangle 
\[ F(L) \xar{F(\al)} F(M)  \xar{F(\be)} F(N)  
\xar{ \ \tau_L \circ \, F(\ga) \ } \opn{T}_{\cat{L}}(F(L)) \]
is a distinguished triangle in $\cat{L}$.

\item Suppose $(G, \nu) : \cat{K} \to \cat{L}$ is another triangulated functor.
A {\em morphism of triangulated functors}
\index{Triangulated functor! morphism of {\indash}s}
$\eta : (F, \tau) \to (G , \nu)$ is 
a morphism of T-additive functors, as in Definition \ref{dfn:1160}.
\end{enumerate}
\end{dfn}

Sometimes we keep the translation isomorphism $\tau$ implicit, and refer to $F$ 
as a triangulated functor.

\begin{dfn}  \label{dfn:1450}
Let $\cat{K}$ be a triangulated category. A 
{\em full triangulated subcategory}% 
\index{Triangulated category! full triangulated subcategory of}
of $\cat{K}$ is a subcategory 
$\cat{L} \subseteq \cat{K}$ satisfying these conditions:
\begin{enumerate}
\rmitem{a} $\cat{L}$ is a full additive subcategory (see Definition 
\ref{dfn:1025}). 

\rmitem{b}
$\cat{L}$ is closed under translations, i.e.\ $L \in \cat{L}$ iff
$\opn{T}(L) \in \cat{L}$.

\rmitem{c} $\cat{L}$ is closed under distinguished triangles, i.e.\ if 
$L' \to L \to L'' \xar{\, \triangle\, }$ 
is a  distinguished triangle in $\cat{K}$ s.t.\ $L', L \in \cat{L}$, then 
also $L'' \in \cat{L}$.
\end{enumerate}
\end{dfn}

Observe that $\cat{L}$ itself is triangulated, and the inclusion 
$\cat{L} \to \cat{K}$ is a triangulated functor. 

\begin{prop} \label{prop:2305}
For $i = 0, 1, 2$ let $(\cat{K}_i, \opn{T}_i)$ be triangulated categories, and 
let 
$(F_i, \tau_i) : (\cat{K}_{i - 1}, \opn{T}_{i - 1}) \to 
(\cat{K}_i, \opn{T}_i)$ 
be triangulated functors. Define the T-additive functor 
$(F, \tau) := (F_2, \tau_2) \circ (F_1, \tau_1)$
as in Definition \tup{\ref{dfn:1840}}. Then 
$(F, \tau) : (\cat{K}_0, \opn{T}_0) \to (\cat{K}_2, \opn{T}_2)$
is a triangulated functor. 
\end{prop}

\begin{exer} \label{exer:2300}
Prove Proposition \ref{prop:2305}. 
\end{exer}

\begin{dfn} \label{dfn:1510}
Let $\cat{K}$ be a triangulated category, and let $\cat{M}$ be an 
abelian category. 
\begin{enumerate}
\item A {\em cohomological functor}%
\index{Cohomological functor}
$F : \cat{K} \to \cat{M}$ is an
additive functor, such that for every distinguished triangle 
\[ L \xar{\al} M \xar{\be} N \xar{\ga} \opn{T}(L) \] 
in $\cat{K}$, the sequence 
\[ F(L) \xar{F(\al)} F(M) \xar{F(\be)} F(N) \] 
is exact in $\cat{M}$.

\item A {\em contravariant cohomological functor}%
\index{Cohomological functor! contravariant}
$F : \cat{K} \to \cat{M}$ is 
a contravariant  additive functor, such that for every distinguished triangle 
\[ L \xar{\al} M \xar{\be} N \xar{\ga} \opn{T}(L) \] 
in $\cat{K}$, the sequence 
\[ F(N) \xar{F(\be)} F(M) \xar{F(\al)} F(L) \] 
is exact in $\cat{M}$.
\end{enumerate}
\end{dfn}

\begin{prop} \label{prop:1281}
Let  $F : \cat{K} \to \cat{M}$ be a cohomological functor, and let 
\[ L \xar{\al} M \xar{\be} N \xar{\ga} \opn{T}(L) \]
be a distinguished triangle in $\cat{K}$. Then the sequence 
\[ \begin{aligned}
&    \cdots \to F(\opn{T}^i(L)) \xar{F(\opn{T}^i(\al))} F(\opn{T}^i(M)) 
\xar{F(\opn{T}^i(\be))} F(\opn{T}^i(N)) 
\\[0.2em]
& \hspace{2ex}
\xar{F(\opn{T}^i(\ga))} F(\opn{T}^{i + 1}(L))
\xar{F(\opn{T}^{i + 1}(\al))} F(\opn{T}^{i + 1}(M)) \to \cdots 
\end{aligned} \]
in $\cat{M}$ is exact.
\end{prop}

\begin{proof}
By axiom (TR2) we have these distinguished triangles:
\[ \opn{T}^i(L) \xar{(-1)^{i} \cd \opn{T}^i(\al)} \opn{T}^i(M) 
\xar{(-1)^{i} \cd \opn{T}^i(\be)} \opn{T}^i(N) 
\xar{(-1)^{i} \cd \opn{T}^i(\ga)} \opn{T}^{i + 1}(L) , \]
\[ \opn{T}^i(M) \xar{(-1)^{i} \cd \opn{T}^i(\be)} \opn{T}^i(N) 
\xar{(-1)^{i} \cd \opn{T}^i(\ga)} \opn{T}^{i + 1}(L) 
\xar{(-1)^{i+1} \cd \opn{T}^{i + 1}(\al)} \opn{T}^{i + 1}(M) , \]
\[ \opn{T}^{i}(N) \xar{(-1)^{i} \cd \opn{T}^{i}(\ga)} 
\opn{T}^{i + 1}(L) \xar{(-1)^{i+1} \cd \opn{T}^{i + 1}(\al)} 
\opn{T}^{i + 1}(M) \xar{(-1)^{i+1} \cd 
\opn{T}^{i + 1}(\be)} \opn{T}^{i + 1}(N) . \]
Now use the definition of a cohomological functor, noting that multiplying 
morphisms in an exact sequence by $-1$ preserves exactness.
\end{proof}

\begin{prop} \label{prop:4}
Let $\cat{K}$ be a triangulated category.
For every $P \in \cat{K}$ the functor
\[ \opn{Hom}_{\cat{K}}(P, -) : \cat{K} \to \cat{Mod} \K \]
is a cohomological functor, and the functor 
\[ \opn{Hom}_{\cat{K}}(-, P) : \cat{K} \to \cat{Mod} \K \]
is a contravariant cohomological functor. 
\end{prop}

\begin{proof}
We will prove the covariant statement. The contravariant statement is 
proved similarly, and we leave this to the reader. 

Consider a distinguished triangle 
$L \xar{\al} M \xar{\be} N \xar{\ga} \opn{T}(L)$
in $\cat{K}$. We have to prove that the sequence 
\[ \opn{Hom}_{\cat{K}}(P, L) \xar{ \opn{Hom}(\opn{id}_P, \al) } 
\opn{Hom}_{\cat{K}}(P, M) \xar{ \opn{Hom}(\opn{id}_P, \be) } 
\opn{Hom}_{\cat{K}}(P, N)  \]
in $\cat{Mod} \K$ is exact. In view of Proposition \ref{prop:1280}, all we need 
to show is that for every $\psi : P \to M$ s.t.\ $\be \circ \psi = 0$, there is 
some  $\phi : P \to L$ s.t.\ $\psi = \al \circ \phi$.  In a picture, we must 
show that the diagram below (solid arrows) 
\[ \UseTips  \xymatrix @C=8ex @R=6ex { 
P
\ar[r]^{\opn{id}}
\ar@{-->}[d]_{\phi}
&
P
\ar[r]^{}
\ar[d]_{\psi}
& 
0
\ar[r]^{}
\ar[d]_{}
&
\opn{T}(P)
\ar@{-->}[d]_{\opn{T}(\phi)}
\\
L
\ar[r]^{\al}
&
M
\ar[r]^{\be}
& 
N
\ar[r]^{\ga}
&
\opn{T}(L) \ .
}  \]
can be completed (dashed arrow). 
This is true by (TR2) (= turning) and (TR3) (= extending). 
\end{proof}

\begin{exer} \label{exer:2301}
Prove the contravariant statement in the proposition above. 
\end{exer}

\begin{que} \label{eqn:1305}
Let  $\cat{K}$ and $\cat{L}$  be triangulated categories, and let 
$F : \cat{K} \to \cat{L}$ be an additive functor. Is it true that there is at 
most one isomorphism of functors 
$\tau :  F \circ \opn{T}_{\cat{K}}  \iso \opn{T}_{\cat{L}} \circ \, F$
such that the pair $(F, \tau)$ is a triangulated functor?
\end{que}




\mysubsection{Some Properties of Triangulated Categories} 
\label{subsec:props-tr-cats}

In this subsection we prove a few general results on triangulated categories. 
Recall that all triangulated categories, and all triangulated functors between 
them, are (implicitly) $\K$-linear.

\begin{prop} \label{prop:1280}
Let $\cat{K}$ be a triangulated category. If
$L \xar{\al} M \xar{\be} N \xar{\ga} \opn{T}(L)$
is a distinguished triangle
in $\cat{K}$, then $\be \circ \al = 0$. 
\end{prop}

\begin{proof}
By axioms (TR1) and (TR3) we have a commutative diagram 
\[ \UseTips  \xymatrix @C=8ex @R=6ex { 
L
\ar[r]^{\opn{id}_L}
\ar[d]_{\opn{id}_L}
&
L
\ar[r]^{}
\ar[d]_{\al}
& 
0
\ar[r]^{}
\ar[d]_{}
&
\opn{T}(L)
\ar[d]_{\opn{T}(\opn{id}_L)}
\\
L
\ar[r]^{\al}
&
M
\ar[r]^{\be}
& 
N
\ar[r]^{\ga}
&
\opn{T}(L) \ .
}  \]
We see that $\be \circ \al$ factors through $0$. 
\end{proof}

\begin{prop}  \label{prop:1283}
Let $\cat{K}$ be a triangulated category, and let 
\[ \UseTips  \xymatrix @C=8ex @R=6ex { 
L 
\ar[r]^{\al}
\ar[d]_{\phi}
&
M
\ar[r]^{\be}
\ar[d]_{\psi}
& 
N
\ar[r]^{\ga}
\ar[d]_{\chi}
&
\opn{T}(L)
\ar[d]_{\opn{T}(\phi)}
\\
L' 
\ar[r]^{\al'}
&
M'
\ar[r]^{\be'}
& 
N'
\ar[r]^{\ga'}
&
\opn{T}(L') \ .
}  \]
be a morphism of distinguished triangles. If $\phi$ and $\psi$ are isomorphisms,
then $\chi$ is also an isomorphism.
\end{prop}

\begin{proof}
Take an arbitrary $P \in \cat{K}$, and let 
$F := \opn{Hom}_{\cat{K}}(P, -)$. We get a 
commutative diagram 
\[ \UseTips  \xymatrix @C=7ex @R=6ex { 
F(L) 
\ar[r]^{F(\al)}
\ar[d]^{F(\phi)}
&
F(M)
\ar[r]^{F(\be)}
\ar[d]_{F(\psi)}
& 
F(N)
\ar[r]^{F(\ga)}
\ar[d]_{F(\chi)}
&
F(\opn{T}(L))
\ar[d]_{F(\opn{T}(\phi))}
\ar[r]^(0.45){F(\opn{T}(\al))}
&
F(\opn{T}(M)) 
\ar[d]_{F(\opn{T}(\psi))}
\\
F(L') 
\ar[r]^{F(\al')}
&
F(M')
\ar[r]^{F(\be')}
& 
F(N')
\ar[r]^{F(\ga')}
&
F(\opn{T}(L'))
\ar[r]^(0.47){F(\opn{T}(\al'))}
&
F(\opn{T}(M')) 
} \]
in $\cat{Mod} \K$. By Proposition \ref{prop:4} the rows in the diagram 
are exact sequences.
Since the other vertical arrows are isomorphisms, it follows that 
\[ F(\chi) : \opn{Hom}_{\cat{K}}(P, N) \to 
\opn{Hom}_{\cat{K}}(P, N') \]
is an isomorphism of abelian groups. By forgetting structure, we see that 
$F(\chi)$ is an isomorphism of sets. 

Consider the Yoneda functor
$\opn{Y}_{\cat{K}} : \cat{K} \to \cat{Fun}(\cat{K}^{\mrm{op}}, \cat{Set})$
from Subsection \ref{subsec:repfunc}. 
There is a morphism 
$ \opn{Y}_{\cat{K}}(\chi) : \opn{Y}_{\cat{K}}(N) \to  
\opn{Y}_{\cat{K}}(N')$
in 
$\cat{Fun}(\cat{K}^{\mrm{op}}, \cat{Set})$. 
For every object $P \in \cat{K}$, letting 
$F := \opn{Hom}_{\cat{K}}(P, -)$
as above, we have 
$\opn{Y}_{\cat{K}}(N)(P) = F(N)$ and
$\opn{Y}_{\cat{K}}(N')(P) = F(N')$.
The calculation above shows that 
\[ F(\chi) = \opn{Y}_{\cat{K}}(\chi)(P) : \opn{Y}_{\cat{K}}(N)(P) \to 
\opn{Y}_{\cat{K}}(N')(P) \]
is an isomorphism in $\cat{Set}$. 
Therefore $\opn{Y}_{\cat{K}}(\chi)$ is an isomorphism in 
$\cat{Fun}(\cat{K}^{\mrm{op}}, \cat{Set})$. 
According to Yoneda Lemma (Theorem \ref{thm:3025})
the morphism $\chi : N \to N'$ in $\cat{K}$ is an isomorphism. 
\end{proof}

\begin{exer} \label{exer:4690}
Try to prove the last proposition without the Yoneda Lemma. 
\end{exer}

Recall that if $\eta : F \to G$ is a morphism of functors
$\cat{K} \to \cat{L}$, then for an object $M \in \cat{K}$
we denote by $\eta_M : F(M) \to G(M)$ the corresponding morphism in 
$\cat{L}$. 

\begin{cor} \label{cor:3601}
Suppose 
$F, G : \cat{K} \to \cat{L}$ 
are triangulated functors between triangulated categories, and 
$\eta : F \to G$ is a morphism of triangulated functors. 
Let 
$L \xar{\al} M \xar{\be} N \xar{\ga} \opn{T}(L)$ 
be a distinguished triangle in $\cat{K}$. If $\eta_L$ and $\eta_M$
are isomorphisms, then $\eta_N$ is an isomorphism.
\end{cor}

\begin{proof}
After applying the functors $F$ and $G$, and the morphisms of functors 
$\tau_{F}$, $\tau_{G}$ and $\eta$, we obtain a 
commutative diagram 
\[  \UseTips  \xymatrix @C=8ex @R=6ex { 
F(L) 
\ar[r]^{F(\al)}
\ar[d]_{\eta_L}
&
F(M)
\ar[r]^{F(\be)}
\ar[d]_{\eta_M}
& 
F(N)
\ar[rr]^(0.43){\tau_{F, L} \circ \, F(\ga)}
\ar[d]_{\eta_N}
&
&
\opn{T}_{\cat{L}}(F(L))
\ar[d]_{\opn{T}_{\cat{L}}(\eta_L)}
\\
G(L) 
\ar[r]^{G(\al)}
&
G(M)
\ar[r]^{G(\be)}
& 
G(N)
\ar[rr]^(0.43){\tau_{G, L} \circ \, G(\ga)}
&
&
\opn{T}_{\cat{L}}(G(L)) 
} \]
in $\cat{L}$ in which the rows are distinguished triangles. 
According to Proposition \ref{prop:1283} the morphism $\eta_N$ is an 
isomorphism.
\end{proof}

\begin{cor} \label{cor:3605}   
Let $\cat{K}$ be a triangulated category, and let
$L \xar{\al} M \xar{} N \xar{\, \triangle\, }$
be a distinguished triangle in it.
The two conditions below are equivalent.  
\begin{enumerate}
\rmitem{i} $\al : L \to M$ is an isomorphism.
\rmitem{ii} $N \cong 0$. 
\end{enumerate}
\end{cor}

\begin{exer} \label{exer:3600}
Prove Corollary \ref{cor:3605}. (Hint: Use Proposition \ref{prop:1283} and 
axiom (TR1)(c).)
\end{exer}

\begin{dfn} \label{dfn:3600}
Let $\cat{K}$ be a triangulated category, and let 
$\al : L \to M$ be a morphism in $\cat{K}$. By axiom (TR1)(b) there exists a 
distinguished triangle 
\[ \tag{$\dag$} L \xar{\al} M \xar{\be} N \xar{\ga} \opn{T}(L) \]
in $\cat{K}$. The object $N$ is called a 
{\em cone of $\al$}%
\index{Cone! of morphism in triangulated category},
and the distinguished triangle ($\dag$) is called a 
{\em distinguished triangle built on $\al$}%
\index{Triangle! distinguished}. 
\end{dfn}

\begin{cor} \label{cor:3600}
In the situation of Definition \tup{\ref{dfn:3600}}, the object $N$ and the 
distinguished triangle \tup{($\dag$)} are unique up to isomorphism.
\end{cor}

\begin{proof}
Suppose 
\[ \tag{$\dag'$} L \xar{\al} M \xar{\be'} N' \xar{\ga'} \opn{T}(L) \]
is another distinguished triangle built on $\al$. By axiom (TR3) there is a 
commutative diagram 
\[ \UseTips  \xymatrix @C=8ex @R=6ex { 
L 
\ar[r]^{\al}
\ar[d]_{\opn{id}_L}
&
M
\ar[r]^{\be}
\ar[d]_{\opn{id}_M}
& 
N
\ar[r]^{\ga}
\ar[d]_{\chi}
&
\opn{T}(L)
\ar[d]_{\opn{T}(\opn{id}_L)}
\\
L 
\ar[r]^{\al}
&
M
\ar[r]^{\be'}
& 
N'
\ar[r]^{\ga'}
&
\opn{T}(L) \ .
}  \]
in $\cat{K}$. Proposition \ref{prop:1283} says that $\chi$ is an isomorphism. 
\end{proof}

\begin{rem} \label{rem:4160}
In general the isomorphism $\chi$ in the corollary above is not unique, and 
thus the cone of $\al$ is not functorial in the morphism $\al$. However, in 
some special cases $\chi$ is unique -- see Remark \ref{rem:3600} below. 

Note also that the standard cones and triangles in the DG category 
$\dcat{C}(A, \cat{M})$ are functorial, by Proposition \ref{prop:1162}.
This implies that in the triangulated categories 
$\dcat{K}(A, \cat{M})$ and $\dcat{D}(A, \cat{M})$
we can sometimes arrange to have functorial cones.
\end{rem}

\begin{prop} \label{prop:3605}
Let $\cat{K}$ be a triangulated category, and let
$L \xar{\al} M \xar{\be} N \xar{\ga} \opn{T}(L)$ 
be a distinguished triangle in it.
The two conditions below are equivalent.  
\begin{enumerate}
\rmitem{i} The morphism $\ga$ is zero.
\rmitem{ii} There exists a morphism 
$\tau : M \to L$ such that $\tau \circ \al = \opn{id}_L$. 
\end{enumerate}
\end{prop}

\begin{proof} \mbox{}

\smallskip \noindent
(i) $\Rightarrow$ (ii): By Propositions \ref{prop:4} the contravariant functor 
$\opn{Hom}_{\cat{K}}(-, L)$ is cohomological. Applying it to the given 
distinguished triangle, and using Proposition \ref{prop:1281}, we get  
an exact sequence of $\K$-modules 
\[ \opn{Hom}_{\cat{K}}(M, L) \xar{\opn{Hom}(\al, \opn{id}_L)}  
\opn{Hom}_{\cat{K}}(L, L) \xar{\opn{Hom}(\opn{T}^{-1}(\ga), \opn{id}_L)}
\opn{Hom}_{\cat{K}}(\opn{T}^{-1}(N), L) . \]
Since the homomorphism 
$\opn{Hom}_{\cat{K}}(\opn{T}^{-1}(\ga), \opn{id}_L)$ is zero, 
there exists some morphism $\tau \in \opn{Hom}_{\cat{K}}(M, L)$
such that 
$\opn{id}_L = \opn{Hom}_{\cat{K}}(\al, \opn{id}_L)(\tau) = \tau \circ \al$.

\medskip \noindent 
(ii) $\Rightarrow$ (i):
Let us examine the commutative diagram 
\[ \UseTips  \xymatrix @C=14ex @R=6ex { 
\opn{Hom}_{\cat{K}}(N, N)
\ar[d]_(0.45){\opn{Hom}(\opn{id}, \ga)}
\\
\opn{Hom}_{\cat{K}}(N, \opn{T}(L))
\ar[d]_{\opn{Hom}(\opn{id}, \opn{T}(\al))}
\ar@(r,ul)[dr]^{ \ \opn{Hom}(\opn{id}, \opn{id})}
\\
\opn{Hom}_{\cat{K}}(N, \opn{T}(M))
\ar[r]^{\opn{Hom}(\opn{id}, \opn{T}(\tau))}
&
\opn{Hom}_{\cat{K}}(N, \opn{T}(L))
}  \]
in $\cat{Mod} \K$. Because the homomorphism 
$\opn{Hom}_{\cat{K}}(\opn{id}, \opn{id})$
is bijective, it follows that 
$\opn{Hom}_{\cat{K}}(\opn{id}, \opn{T}(\al))$
is injective. But the column is an exact sequence (by Propositions \ref{prop:4} 
and \ref{prop:1281}), and therefore 
$\opn{Hom}_{\cat{K}}(\opn{id}_{N}, \ga) = 0$. We conclude that
$\ga = \opn{Hom}_{\cat{K}}(\opn{id}_{N}, \ga)(\opn{id}_{N}) = 0$.
\end{proof}

\begin{lem} \label{lem:3600}
Let $M, M'$ be objects in a triangulated category $\cat{K}$.
Consider the canonical morphisms 
\[ M  \xar{e} M \oplus M' \xar{p} M \quad \tup{and} \quad
M'  \xar{e'} M \oplus M' \xar{p'} M' . \]
Then 
\[ M \xar{e} M \oplus M' \xar{p'} M' \xar{0} \opn{T}(M) \]
is a distinguished triangle in $\cat{K}$.
\end{lem}

\begin{proof}
By axiom (TR1)(c) and axiom (TR2) there is a distinguished triangle 
$0 \xar{0} M' \xar{\opn{id}_{M'}} M' \xar{0} 0$.
By axiom (TR1)(b) there is a distinguished triangle
\begin{equation} \label{eqn:4160}
 M \xar{e} M \oplus M' \xar{\be} N \xar{\ga} \opn{T}(M) 
\end{equation}
in $\cat{K}$, for some object $N$. 
Because $p \circ e = \opn{id}_M$, Proposition \ref{prop:3605} 
says that $\ga = 0$. 

Next, since $p' \circ e = 0$, axiom (TR3) produces a morphism of triangles 
\begin{equation} \label{eqn:3600}
\UseTips  \xymatrix @C=8ex @R=6ex { 
M 
\ar[r]^(0.4){e}
\ar[d]_{0}
&
M \oplus M'
\ar[r]^(0.6){\be}
\ar[d]_{p'}
& 
N 
\ar[r]^(0.47){\ga = 0}
\ar@{-->}[d]_{\psi}
&
\opn{T}(M)
\ar[d]_{0}
\\
0
\ar[r]^{0}
&
M'
\ar[r]^{\opn{id}_{M'}}
&
M'
\ar[r]^{0}
& 
0
} 
\end{equation}
We claim that $\psi$ is an isomorphism. This is proved 
indirectly. For every object $L \in \cat{K}$ there is a commutative diagram 
\begin{equation} \label{eqn:3602}
\UseTips  \xymatrix @C=4ex @R=6ex { 
\opn{Hom}_{\cat{K}}(L, M)
\ar[r]
&
\opn{Hom}_{\cat{K}}(L, M \oplus M')
\ar[r]
\ar[d]_{\opn{Hom}(\opn{id}_L, p')}
& 
\opn{Hom}_{\cat{K}}(L, N)
\ar[r]
\ar[d]_{\opn{Hom}(\opn{id}_L, \psi)}
&
0
\\
0
\ar[r]
&
\opn{Hom}_{\cat{K}}(L, M')
\ar[r]^{\opn{id}}
& 
\opn{Hom}_{\cat{K}}(L, M')
\ar[r]
&
0
}
\end{equation}
in $\cat{Mod} \K$, that's gotten by applying the functor 
$\opn{Hom}_{\cat{K}}(L,- )$ to the left part of diagram (\ref{eqn:3600}).
The rows in (\ref{eqn:3602}) are exact sequences. 
Because 
\[ \opn{Hom}_{\cat{K}}(L, M \oplus M') \cong 
\opn{Hom}_{\cat{K}}(L, M) \oplus \opn{Hom}_{\cat{K}}(L, M') \]
we can replace (\ref{eqn:3602}) with the next commutative diagram
\begin{equation} \label{eqn:4161}
\UseTips  \xymatrix @C=4ex @R=6ex { 
0
\ar[r]
&
\opn{Hom}_{\cat{K}}(L, M')
\ar[r]
\ar[d]_{\opn{Hom}(\opn{id}_L, \opn{id}_{M'})}
& 
\opn{Hom}_{\cat{K}}(L, N)
\ar[r]
\ar[d]_{\opn{Hom}(\opn{id}_L, \psi)}
&
0
\\
0
\ar[r]
&
\opn{Hom}_{\cat{K}}(L, M')
\ar[r]^{\opn{id}}
& 
\opn{Hom}_{\cat{K}}(L, M')
\ar[r]
&
0
}
\end{equation}
with exact rows. We see that 
$\opn{Hom}(\opn{id}_L, \psi)$ is an isomorphism of $\K$-modules. Hence it is a  
bijection of sets. Using the Yoneda Lemma, like in the proof of Proposition 
\ref{prop:1283}, we conclude that $\psi$ is an isomorphism in $\cat{K}$. 

Finally, from the commutative diagram (\ref{eqn:3600}) we know that 
$\psi \circ \be = p'$. So we have this isomorphism of triangles: 
\begin{equation} \label{eqn:3605}
\UseTips  \xymatrix @C=7ex @R=6ex { 
M 
\ar[r]^(0.4){e}
\ar[d]_{\opn{id}}
&
M \oplus M'
\ar[r]^(0.64){\be}
\ar[d]_{\opn{id}_{}}
& 
N 
\ar[r]^(0.45){\ga = 0}
\ar[d]_{\psi}^{\cong}
&
\opn{T}(M)
\ar[d]_{\opn{id}}
\\
M
\ar[r]^(0.4){e}
&
M \oplus M'
\ar[r]^(0.64){p'}
&
M'
\ar[r]^(0.45){0}
& 
\opn{T}(M)
} 
\end{equation}
The first triangle is distinguished. By axiom (TR1)(a) the second triangle is 
also distinguished. 
\end{proof}

\begin{dfn} \label{dfn:3235}
Let $\cat{K}$ be a linear category, and let $M, N \in \cat{K}$.
\begin{enumerate}
\item The object $M$ is called a {\em retract}
\index{Retract}
of $N$ if there are morphisms 
$M \xar{e} N \xar{p} M$ such that $p \circ e = \opn{id}_M$. 
The morphism $e : M \to N$ is called an embedding. 

\item The object $M$ is called a {\em direct summand}
\index{Direct summand}
of $N$ if there is an 
object $M' \in \cat{K}$ and an isomorphism $M \oplus M' \cong N$. 
The corresponding morphism $e : M \to N$ is called an embedding. 
\end{enumerate}
\end{dfn}

Note that in both items above the embedding $e : M \to N$ is a monomorphism 
in $\cat{K}$. 

\begin{thm} \label{thm:4160}  
Let $\cat{K}$ be a triangulated category, and let $e : M \to N$
be a morphism in $\cat{K}$. 
The following two conditions are equivalent\tup{:}
\begin{enumerate}
\rmitem{i} There is a morphism $e' : M' \to N$ in $\cat{K}$ such that 
\[ (e, e') : M \oplus M' \to N \]
is an isomorphism. In other words, $M$ is a direct summand of $N$,
\index{Direct summand}
with embedding $e : M \to N$. 

\rmitem{ii} There is a morphism $p : N \to M$  in $\cat{K}$  such that 
$p \circ e = \opn{id}_M$. In other words, $M$ is a retract of $N$, 
\index{Retract}
with embedding $e : M \to N$. 
\end{enumerate}
\end{thm}

\begin{proof} \mbox{}

\smallskip \noindent 
(i) $\Rightarrow$ (ii): This is trivial (and only requires $\cat{K}$ to be a 
linear category). 

\medskip \noindent 
(ii) $\Rightarrow$ (i): Let 
$M \xar{e} N \xar{\be} M' \xar{\ga} \opn{T}(M)$
be a distinguished triangle in $\cat{K}$ built on $e$. By Proposition 
\ref{prop:3605} we know that $\ga = 0$. According to Lemma \ref{lem:3600},
$M \xar{\ep} M \oplus M' \xar{\pi'} M' \xar{0} \opn{T}(M)$
is also the distinguished triangle in $\cat{K}$;
here $\ep$ is the embedding and $\pi'$ is the projection. 
Turning these two distinguished triangles and using axiom (TR3) we get a 
morphism $\th$ such that the diagram 
\[ \UseTips  \xymatrix @C=8ex @R=6ex {
\opn{T}^{-1}(M') 
\ar[r]^(0.6){0}
\ar[d]_{\opn{id}}
&
M 
\ar[r]^(0.4){\ep}
\ar[d]_{\opn{id}}
&
M \oplus M'
\ar[r]^(0.58){\pi'}
\ar@{-->}[d]_{\th}
& 
M'
\ar[d]_{\opn{id}}
\\
\opn{T}^{-1}(M') 
\ar[r]^(0.6){0}
&
M
\ar[r]^(0.5){e}
& 
N 
\ar[r]^(0.45){\be}
&
M'
} \]
is commutative. By Proposition \ref{prop:1283}, $\th$ is an isomorphism. 
Finally we define $e' := \th \circ \ep'$, where 
$\ep' : M' \to M \oplus M'$ is the embedding. 
\end{proof}

\begin{dfn} \label{dfn:4915}
Let $\cat{K}$ be a triangulated category. A {\em saturated full triangulated 
subcategory} of $\cat{K}$ is a full triangulated subcategory 
$\cat{K}' \sub \cat{K}$ (Definition \ref{dfn:1450}) which is closed in $\cat{K}$
under taking direct summands and under isomorphisms. 
\end{dfn}

\begin{dfn} \label{dfn:4916}
Let $\cat{K}$ be a triangulated category and let $Z \sub \opn{Ob}(\cat{K})$ be 
a set of objects. The {\em saturated full triangulated subcategory of $\cat{K}$
generated by $Z$} is the smallest saturated full triangulated subcategory 
$\cat{K}' \sub \cat{K}$ such that $Z \sub \opn{Ob}(\cat{K}')$.
\end{dfn}

\begin{prop} \label{prop:4470}
Let $\cat{K}$ be a triangulated category, let $Z \sub \opn{Ob}(\cat{K})$ 
be a set of objects, and let $\cat{K}' \sub \cat{K}$ be the saturated full 
triangulated subcategory of $\cat{K}$ generated by $Z$. An object 
$M \in \cat{K}$ belongs to $\cat{K}'$ if 
and only if there is a sequence $M_0, \ldots, M_r$ of objects of $\cat{K}$, 
such that $M_r = M$, and for every $i \leq r$ at least one of the following 
conditions holds\tup{:}
\begin{itemize}
\rmitem{a} $M_i \in Z$. 

\rmitem{b} There is an isomorphism $M_i \cong \opn{T}^p(M_j)$ in $\cat{K}$
for some $j < i$ and $p \in \Z$. 

\rmitem{c} There is a distinguished triangle 
$M_j \to M_k \to M_i \xar{\, \triangle\, }$
in $\cat{K}$ for some $j, k < i$.

\rmitem{d} $M_i$ is a direct summand of $M_j$ in $\cat{K}$ for some $j < i$.
\end{itemize}
\end{prop}

\begin{prop} \label{prop:4603}
Let $\cat{K}$ and $\cat{L}$ be triangulated categories, 
let $F, G : \cat{K} \to \cat{L}$ be triangulated functors, and let 
$\ze : F \to G$ be a morphism of triangulated functors. 
Denote by $\cat{K}'$ the full subcategory of $\cat{K}$ on the objects
$M$ such that $\ze_M : F(M) \to G(M)$ is an isomorphism in $\cat{L}$.
Then $\cat{K}'$ is a saturated full triangulated subcategory of $\cat{K}$. 
\end{prop}

\begin{cor} \label{cor:4920}
Let $\cat{K}$ and $\cat{L}$ be triangulated categories, 
let $F : \cat{K} \to \cat{L}$ be a triangulated functor, and let
$\cat{K}'$ be the kernel of $F$, i.e.\ $\cat{K}' \sub \cat{K}$ is the 
full subcategory on the objects $M \in \cat{K}$ such that $F(M) = 0$. 
Then $\cat{K}'$ is a saturated full triangulated subcategory of $\cat{K}$.
\end{cor}

\begin{exer} \label{exer:4603}
Prove Proposition \ref{prop:4470}, Proposition \ref{prop:4603}  and Corollary 
\lb \ref{cor:4920}. 
\end{exer}

We end this subsection with two remarks. 

\begin{rem} \label{rem:4777}
The story of the concept ``saturated full triangulated 
subcategory'' is convoluted. In the text \cite{Ver0} from 1977, J.-L. Verdier 
introduced {\em \'epaisse full triangulated subcategories}
(a rather messy definition). However, J. Rickard proved (see  
\cite[Proposition 1.3]{Ric0}, published in 1989) that \'epaisse is equivalent 
to saturated, as defined above. Rickard did not give a name for this new 
and useful property. 

Verdier died in 1989. A few years later, in 1996, G. Maltsiniotis published 
Verdier's thesis \cite{Ver}. In this thesis the word ``\'epaisse'' is absent, 
and instead we find the definition of saturated categories, identical to 
Definition \ref{dfn:4915}. Maltsiniotis notes that Verdier was aware (although 
this is only implicit in his writing) of the equivalence between \'epaisse and 
saturated; indeed, these are precisely the categories that occur as kernels of 
triangulated functors (see Corollary \ref{cor:4920} above, and Remark 
\ref{rem:4220} on Verdier quotients and saturated denominator 
sets of morphisms). 
\end{rem}

\begin{rem} \label{rem:3600}
As mentioned above, in Remark \ref{rem:4160}, the cone of a morphism 
$\al : L \to M$ in a triangulated category $\cat{D}$ is usually not functorial. 
However, if $\cat{D}$ is endowed with a {\em t-structure},
\index{t-structure}
then some cones can be made functorial. Here is a quick explanation. 

The concept of t-structure was introduced by A.A. Beilinson, J. Bernstein and 
P. Deligne in the book \cite{BBD}. A t-structure on a triangulated category 
$\cat{D}$ consists of a pair $(\cat{D}^{\leq 0}, \cat{D}^{\geq 0})$ of full 
subcategories of $\cat{D}$, that satisfy a few axioms (see e.g.\ 
\cite[Chapter X]{KaSc1} or \cite[Definition 4.1]{YeZh2}). The prototypical 
example is the {\em standard t-structure} on $\cat{D} := \dcat{D}(\cat{M})$, 
the derived category of an abelian category $\cat{M}$
(see Section \ref{sec:der-cat} of the book). Here $\cat{D}^{\leq 0}$ 
is the full subcategory of $\dcat{D}(\cat{M})$ on the complexes $M$ with 
nonpositive cohomology, and 
$\cat{D}^{\geq 0}$ is the full subcategory of $\dcat{D}(\cat{M})$ 
on the complexes $M$ with nonnegative cohomology. Other t-structures on 
$\cat{D} = \dcat{D}(\cat{M})$ are often called {\em perverse t-structures}. 

The {\em heart} of the t-structure $(\cat{D}^{\leq 0}, \cat{D}^{\geq 0})$ is 
the category 
$\cat{D}^{0} := \cat{D}^{\leq 0} \cap \cat{D}^{\geq 0}$.
In the prototypical example above, the heart is the category of complexes $M$ 
with cohomology concentrated in degree $0$, so it is equivalent to $\cat{M}$. 
Something similar happens in general: the heart $\cat{D}^{0}$ is always an 
abelian category, and there is a cohomological functor 
$\mrm{H}^0 : \cat{D} \to \cat{D}^{0}$.
The short exact sequences 
$0 \to L \xar{\al} M \xar{\be} N \to 0$
in $\cat{D}^{0}$ are the distinguished triangles 
$L \xar{\al} M \xar{\be} N \xar{\, \triangle\, }$
in $\cat{D}$ such that $L, M, N \in \cat{D}^{0}$.

Here is how a t-structure $(\cat{D}^{\leq 0}, \cat{D}^{\geq 0})$ affects 
functoriality of cones. 
Axiom (TR3) guarantees that for a pair of distinguished triangles and 
morphism $\phi, \psi$ as in the solid diagram below, there is {\em at least 
one} morphism $\chi$ (the dashed arrow) making the whole diagram commutative.  
\[ \UseTips  \xymatrix @C=8ex @R=6ex { 
L 
\ar[r]^{\al}
\ar[d]_{\phi}
&
M
\ar[r]^{\be}
\ar[d]_{\psi}
& 
N
\ar[r]^{\ga}
\ar@{-->}[d]_{\chi}
&
\opn{T}(L)
\ar[d]_{\opn{T}(\phi)}
\\
L' 
\ar[r]^{\al'}
&
M'
\ar[r]^{\be'}
& 
N'
\ar[r]^{\ga'}
&
\opn{T}(L') 
} \]
According to \cite[Proposition 1.1.9]{BBD} (see also \cite[Lemma 4.10]{YeZh2}),
if $L \in \cat{D}^{\leq 0}$ and $\opn{T}(N') \in \cat{D}^{\geq 0}$, then this 
morphism $\chi$ is {\em unique}. 

The idea of a t-structure emerged in the study of algebraic topology,  
microlocal analysis and representations of Lie algebras. In that context the 
triangulated category $\cat{D}$ was 
the full subcategory of $\dcat{D}(\K_X)$ on complexes with suitable 
finiteness and constructibility conditions. The base ring $\K$ was either a 
field or the $l$-adic integers $\what{\Z}_l$. The space (or site) $X$ was 
either a topological space with a classical metric topology, or
a scheme with its \'etale topology. 
The objects of the heart where called 
{\em perverse sheaves}. This theory is explained in the original book 
\cite{BBD}, and also in \cite{Bor} and in \cite[Chapter X]{KaSc1}. 
In algebraic geometry one can also consider {\em perverse coherent sheaves}; 
see \cite{YeZh2} and \cite{AriBez}. 
\end{rem}




\mysubsection{The Homotopy Category is Triangulated} 
\label{subsec:K-is-triang}

In this subsection we consider a $\K$-linear abelian category $\cat{M}$ and a 
DG central $\K$-ring $A$, where $\K$ is the commutative base ring. These 
ingredients give rise to the $\K$-linear DG category $\dcat{C}(A, \cat{M})$ of 
DG $A$-modules in $\cat{M}$, as in Subsection \ref{subsec:DGModinM}. 

The strict category $\dcat{C}_{\mrm{str}}(A, \cat{M})$ and the 
homotopy category $\dcat{K}(A, \cat{M})$ were introduced in Definition
\ref{dfn:1132}. Recall that these $\K$-linear categories have the same objects 
as $\dcat{C}(A, \cat{M})$. The morphism $\K$-modules are 
\[ \opn{Hom}_{\dcat{C}_{\mrm{str}}(A, \cat{M})} (M_0, M_1) = 
\opn{Z}^0 \bigl( \opn{Hom}_{\dcat{C}(A, \cat{M})} (M_0, M_1) \bigr) \]
and 
\[ \opn{Hom}_{\dcat{K}(A, \cat{M})} (M_0, M_1) = 
\opn{H}^0 \bigl( \opn{Hom}_{\dcat{C}(A, \cat{M})} (M_0, M_1) \bigr) . \]
There is a full additive functor 
$\opn{P} : \dcat{C}_{\mrm{str}}(A, \cat{M}) \to \dcat{K}(A, \cat{M})$ 
which is the identity on objects, and on morphisms it sends a homomorphism to 
its homotopy class. 
Morphisms in $\dcat{K}(A, \cat{M})$ will often by decorated with a bar, like
$\bar{\phi}$. 

Consider the translation functor $\opn{T}$ from  Definition 
\ref{dfn:1275}. Since $\opn{T}$ is a DG functor from $\dcat{C}(A, \cat{M})$ to 
itself (see Corollary \ref{cor:1300}), it restricts to a linear functor
from $\dcat{C}_{\mrm{str}}(A, \cat{M})$ to itself, and it induces a linear 
functor $\bar{\opn{T}}$ from $\dcat{K}(A, \cat{M})$ to itself, such that 
$\opn{P} \circ \opn{T} = \bar{\opn{T}} \circ \opn{P}$.

\begin{prop} \label{prop:1416} 
\mbox{}
\begin{enumerate}
\item The category $\dcat{C}_{\mrm{str}}(A, \cat{M})$,
equipped with the translation functor $\opn{T}$, is a T-additive category. 

\item The category $\dcat{K}(A, \cat{M})$,
equipped with the translation functor $\bar{\opn{T}}$, is a T-additive 
category. 

\item Let $\tau : \opn{P} \circ \opn{T} \iso \bar{\opn{T}} \circ \opn{P}$
be the identity automorphism. Then the pair 
\[ (\opn{P}, \tau) : \dcat{C}_{\mrm{str}}(A, \cat{M}) \to \dcat{K}(A, \cat{M}) 
\]
is a T-additive functor.
\end{enumerate}
\end{prop}

\begin{proof}
(1) We need to prove that $\dcat{C}_{\mrm{str}}(A, \cat{M})$
is additive. Of course the zero complex is a zero object. 
Next we consider finite direct sums. 
Let $M_1, \ldots, M_r$ be a finite collection of objects in 
$\dcat{C}(A, \cat{M})$. Each $M_i$ is a DG $A$-module in $\cat{M}$, and we 
write it as $M_i = \{ M_i^j \}_{j \in \Z}$. 
In each degree $j$ the direct sum 
$M^j := \bigoplus_{i = 1}^r M^j_i$ 
exists in $\cat{M}$. Let $M := \{ M^j \}_{j \in \Z}$ be the resulting graded 
object in $\cat{M}$. The differential $\d_M : M^j \to M^{j + 1}$ exists 
by the universal property of direct sums; so we obtain a complex 
$M \in \dcat{C}(\cat{M})$. 
The DG $A$-module structure on $M$ is defined 
similarly: for $a \in A^k$, there is an induced degree $k$ morphism 
$f(a) : M \to M$ in $\dcat{C}(\cat{M})$. Thus $M$ becomes an object of
$\dcat{C}(A, \cat{M})$. But the embeddings $e_i : M_i \to M$ are strict 
morphisms, so $(M, \{ e_i \})$ is a coproduct of the collection $\{ M_i \}$ 
in $\dcat{C}_{\mrm{str}}(A, \cat{M})$.

\medskip \noindent 
(2) Now consider the category $\dcat{K}(A, \cat{M})$.
Because the functor 
$\opn{P} : \dcat{C}_{\mrm{str}}(A, \cat{M}) \to \dcat{K}(A, \cat{M})$
is additive, and is bijective on objects, part (1) above and Proposition 
\ref{prop:2} say that $\dcat{K}(A, \cat{M})$ is an additive category. 

\medskip \noindent 
(3) Clear. 
\end{proof}

From now on  we  denote by $\opn{T}$, instead of by $\bar{\opn{T}}$, the 
translation functor of $\dcat{K}(A, \cat{M})$. 


\begin{dfn} \label{dfn:140}
A triangle 
\[ L \xar{\bar{\al}} M \xar{\bar{\be}} N \xar{\bar{\ga}} \opn{T}(L)  \]
in $\dcat{K}(A, \cat{M})$ is said to be a {\em distinguished triangle}
\index{Triangle! distinguished}
if there is a standard triangle 
\[ L' \xar{\al'} M' \xar{\be'} N' \xar{\ga'} \opn{T}(L')  \]
in $\dcat{C}_{\mrm{str}}(A, \cat{M})$, as in Definition \ref{dfn:1140},
and an isomorphism of triangles 
\[ \UseTips  \xymatrix @C=8ex @R=6ex { 
L' 
\ar[r]^{\opn{P}(\al')}
\ar[d]_{ \bar{\phi} } 
&
M'
\ar[r]^{\opn{P}(\be')}
\ar[d]_{ \bar{\psi} }
& 
N'
\ar[r]^{\opn{P}(\ga')}
\ar[d]_{ \bar{\chi} }
&
\opn{T}(L')
\ar[d]_{ \opn{T}(\bar{\phi}) }
\\
L 
\ar[r]^{\bar{\al}}
&
M
\ar[r]^{\bar{\be}}
& 
N
\ar[r]^{\bar{\ga}}
&
\opn{T}(L)
\ .
}  \]
in $\dcat{K}(A, \cat{M})$.
\end{dfn}

\begin{thm} \label{thm:1255}
The T-additive category $\dcat{K}(A, \cat{M})$, with the set of 
distinguished triangles defined above, is a triangulated category%
\index{Triangulated category! $\dcat{K}(A, \cat{M})$ is a}%
\index{1-K(A,M)@$\dcat{K}(A, \cat{M})$}.
\end{thm}

The proof is after three lemmas. 

\begin{lem}  \label{lem:2}
Let $M \in \dcat{C}(A, \cat{M})$, and consider the standard cone 
$N := \opn{Cone}(\opn{id}_M)$. Then the DG module $N$ is null-homotopic, i.e.\ 
$0 \to N$ is an isomorphism in $\dcat{K}(A, \cat{M})$.
\end{lem}

\begin{proof}
We shall exhibit a homotopy $\th$ from $0_N$ to $\opn{id}_N$. 
Recall from Subsection \ref{subsec:cone} that 
$N = \opn{Cone}(\opn{id}_M) = M \oplus \opn{T}(M) = 
\sbmat{ M \\[0.2em]\opn{T}(M) }$
as graded modules, with a differential whose matrix presentation is
$\d_{N} = \sbmat{\d_M & \opn{t}_M^{-1} \\[0.2em] 0 & \d_{\opn{T}(M)}}$.
And by the definition in Subsection \ref{subsec:translation} we have
$\d_{\opn{T}(M)} = - \opn{t}_M \circ \, \d_M \circ \opn{t}_M^{-1}$. 
Define 
$\th : N \to N$ to be the degree $-1$ morphism with  matrix presentation
$\th :=  \sbmat{ 0 & \, 0 \\[0.2em] \opn{t}_M & \, 0 }$.
Then, using the formulas above for $\d_N$ and $\d_{\opn{T}(M)}$, we get
\[ \d_N \circ \th + \th \circ \d_N =  
\bmat{ \opn{id}_M & 0 \\[0.1em] 0 & \opn{id}_{\opn{T}(M)} } = \opn{id}_N. 
\qedhere \]
\end{proof}

\begin{exer} \label{exer:1255}
Here is a generalization of Lemma \ref{lem:2}. Consider a morphism 
$\phi : M_0 \to M_1$ in $\dcat{C}_{\mrm{str}}(A, \cat{M})$. Show that the three 
conditions below are equivalent:
\begin{enumerate}
\rmitem{i} $\phi$ is a homotopy equivalence.
\rmitem{ii} $\bar{\phi}$ is an isomorphism in $\dcat{K}(A, \cat{M})$.
\rmitem{iii} The DG module $\opn{Cone}(\phi)$ is null-homotopic.
\end{enumerate}
Try to do this directly, not using Proposition \ref{prop:4}(2) and Theorem 
\ref{thm:1255}.
\end{exer}

The next lemma is based on \cite[Lemma 1.4.2]{KaSc1}.

\begin{lem} \label{lem:1}
Consider a morphism $\al : L \to M$ in 
$\dcat{C}_{\mrm{str}}(A, \cat{M})$, 
the standard  triangle  
$L \xar{\al} M \xar{\be} N \xar{\ga} \opn{T}(L)$
built on $\al$, and the standard  triangle
$M \xar{\be} N \xar{\phi} P \xar{\psi}  \opn{T}(M)$
built on $\be$, all in $\dcat{C}_{\mrm{str}}(A, \cat{M})$. 
So $N = \opn{Cone}(\al)$ and $P = \opn{Cone}(\be)$.
There is a morphism 
$\rho : \opn{T}(L) \to P$ in $\dcat{C}_{\mrm{str}}(A, \cat{M})$ s.t.\ 
$\bar{\rho}$ is an isomorphism in $\dcat{K}(A, \cat{M})$, and the diagram 
\[ \UseTips \xymatrix @C=8ex @R=6ex { 
M
\ar[r]^{\bar{\be}}
\ar[d]_{\opn{id}_M}
&
N
\ar[r]^{\bar{\ga}}
\ar[d]_{\opn{id}_N}
& 
\opn{T}(L)
\ar[r]^{- \opn{T}(\bar{\al})}
\ar[d]_{\bar{\rho}}
&
\opn{T}(M)
\ar[d]_{\opn{id}_{\opn{T}(M)}}
\\
M
\ar[r]^{\bar{\be}}
&
N
\ar[r]^{\bar{\phi}}
& 
P
\ar[r]^{\bar{\psi}}
&
\opn{T}(M) 
} \]
commutes in $\dcat{K}(A, \cat{M})$.
\end{lem}

\begin{proof}
Note that 
$N = M \oplus \opn{T}(L)$ and
$P = N \oplus \opn{T}(M) = M \oplus \opn{T}(L) \oplus \opn{T}(M)$
as graded modules. Thus $P$ and $\d_P$ have the following matrix presentations:
\[ P = \bmat{ M \\[0.2em] \opn{T}(L) \\[0.2em] \opn{T}(M) } \ , \quad
\d_P =  \bmat{ \d_M & \al \circ \opn{t}_L^{-1} & \opn{t}_M^{-1} 
\\[0.3em] 0 & \d_{\opn{T}(L)} & 0 
\\[0.3em] 0 &  0 & \d_{\opn{T}(M)} } \ . \]
Define morphisms $\rho: \opn{T}(L) \to P$ and $\chi : P \to \opn{T}(L)$ in 
$\dcat{C}_{\mrm{str}}(A, \cat{M})$ by the matrix presentations
\[ \rho := \bmat{ 0 \\[0.1em] \opn{id}_{\opn{T}(L)}  
\\[0.2em] - \opn{T}(\al) } \  , \quad
\chi := \bmat{ 0 & \opn{id}_{\opn{T}(L)} & 0 } \ .  \]
Direct calculations show that
$\chi \circ \rho = \opn{id}_{\opn{T}(L)}$, \,
$\rho \circ \ga = \rho \circ \chi \circ \phi$ and 
$\psi \circ \rho = \lb - \opn{T}(\al)$.

It remains to prove that $\rho \circ \chi$ is homotopic to $\opn{id}_P$. 
Define a degree $-1$ morphism $\th : P \to P$ by the matrix
\[ \th := \bmat{ 0 & 0 & 0 \\[0em] 0 & 0 & 0 
\\[0em] \opn{t}_{M} & 0 & 0} . \]
Then a direct calculation, using the equalities 
$\opn{t}_{M} \circ \, \d_M + \d_{\opn{T}(M)} \circ \opn{t}_{M} = 0$
and 
$\opn{T}(\al) =  \opn{t}_{M} \circ \, \al \circ \opn{t}_{L}^{-1}$
gives 
$\th \circ \d_P + \d_P \circ \th = \opn{id}_P - \rho  \circ \chi$.
\end{proof}

\begin{lem} \label{lem:1265} 
Consider a  standard triangle 
$L \xar{\al} M \xar{\be} N \xar{\ga} \opn{T}(L)$
in $\dcat{C}_{\mrm{str}}(A, \cat{M})$.
For every integer $k$, the triangle 
\[ \opn{T}^k(L) \xar{ \, \opn{T}^k(\al) \, } \opn{T}^k(M) 
\xar{ \, \opn{T}^k(\be) \, } \opn{T}^k(N) 
\xar{ \, (-1)^k \cd \opn{T}^k(\ga) \, } 
\opn{T}^{k + 1}(L)  \]
is isomorphic, in $\dcat{C}_{\mrm{str}}(A, \cat{M})$, to a standard triangle.
\end{lem}

\begin{proof}
Combine Corollary \ref{cor:1300}, Corollary \ref{cor:1915} with $F = \opn{T}$,  
and Proposition \ref{prop:1300}.
\end{proof}

\begin{proof}[Proof of Theorem \tup{\ref{thm:1255}}] 
We essentially follow the proof of \cite[Proposition 1.4.4]{KaSc1}, adding some
details.

\medskip \noindent
(TR1): By definition the set of distinguished triangles in 
$\dcat{K}(A, \cat{M})$ is closed under isomorphisms. This establishes item (a). 

As for item (b): consider any morphism $\bar{\al} : L \to M$ in
$\dcat{K}(A, \cat{M})$. It is represented by a morphism $\al : L \to M$ in 
$\dcat{C}_{\mrm{str}}(A, \cat{M})$. Take the standard triangle on $\al$ 
in $\dcat{C}_{\mrm{str}}(A, \cat{M})$. Its image in $\dcat{K}(A, \cat{M})$
has the desired property. 

Finally, Lemma \ref{lem:2} shows that the triangle 
$M \xar{\opn{id}_M} M \to 0 \to \opn{T}(M)$ 
is isomorphic in $\dcat{K}(A, \cat{M})$ to the triangle
$M \xar{\opn{id}_M} M \xar{\bar{e}} \opn{Cone}(\opn{id}_M) \xar{\bar{p}} 
\opn{T}(M)$.
The latter is the image of a standard triangle, and so it is distinguished.

\medskip \noindent
(TR2): Consider the triangles
\begin{equation} \label{eqn:1255}
L \xar{\bar{\al}} M \xar{\bar{\be}} N \xar{\bar{\ga}} \opn{T}(L)
\end{equation}
and
\begin{equation} \label{eqn:1256}
M \xar{\bar{\be}} N \xar{\bar{\ga}} \opn{T}(L)\xar{- \opn{T}(\bar{\al})} 
\opn{T}(M) 
\end{equation}
in $\dcat{K}(A, \cat{M})$. If (\ref{eqn:1255}) is distinguished, then by 
Lemma \ref{lem:1} so is (\ref{eqn:1256}). 

Conversely, if (\ref{eqn:1256}) is distinguished, then by turning it $5$ times, 
and using the previous step (namely by Lemma \ref{lem:1}), we see that the 
triangle
\[ \opn{T}^2(L) \xar{\opn{T}^2(\bar{\al})} \opn{T}^2(M) 
\xar{\opn{T}^2(\bar{\be})} \opn{T}^2(N) \xar{\opn{T}^2(\bar{\ga})} 
\opn{T}^3(L) \]
is distinguished. 
According to Lemma \ref{lem:1265} (with $k = -2$), the triangle 
gotten by applying $\opn{T}^{-2}$ to this is distinguished.
But this is just the triangle (\ref{eqn:1255}).

\medskip \noindent
(TR3): 
Consider a commutative diagram in $\dcat{K}(A, \cat{M})$~:
\begin{equation} \label{eqn:1915}
\UseTips \xymatrix @C=8ex @R=6ex { 
\bar{L} 
\ar[r]^{\bar{\al}}
\ar[d]_{\bar{\phi}}
&
\bar{M}
\ar[r]^{\bar{\be}}
\ar[d]_{\bar{\psi}}
& 
\bar{N}
\ar[r]^{\bar{\ga}}
&
\opn{T}(\bar{L})
\\
\bar{L}' 
\ar[r]^{\bar{\al}'}
&
\bar{M}'
\ar[r]^{\bar{\be}'}
& 
\bar{N}'
\ar[r]^{\bar{\ga}'}
&
\opn{T}(\bar{L}') 
}
\end{equation}
where the horizontal triangles are distinguished. 
By definition the rows in (\ref{eqn:1915}) are isomorphic in 
$\dcat{K}(A, \cat{M})$
to the images under the functor $\opn{P}$ of 
standard triangles in $\dcat{C}_{\mrm{str}}(A, \cat{M})$. These are the rows in 
diagram (\ref{eqn:1916}) below. 
The vertical morphisms in  (\ref{eqn:1915}) are also induced from morphisms in 
$\dcat{C}_{\mrm{str}}(A, \cat{M})$, i.e.\ 
$\bar{\phi} = \opn{P}(\phi)$ and  $\bar{\psi} = \opn{P}(\psi)$. 
Thus (\ref{eqn:1915}) is isomorphic to the image under $\opn{P}$ of the 
following diagram:
\begin{equation} \label{eqn:1916}
\UseTips \xymatrix @C=8ex @R=6ex { 
L 
\ar[r]^{\al}
\ar[d]_{\phi}
&
M
\ar[r]^{\be}
\ar[d]_{\psi}
& 
N
\ar[r]^{\ga}
&
\opn{T}(L)
\\
L' 
\ar[r]^{\al'}
&
M'
\ar[r]^{\be'}
& 
N'
\ar[r]^{\ga'}
&
\opn{T}(L') 
} 
\end{equation}
Warning: the diagram (\ref{eqn:1916}) is only  commutative up to homotopy in 
 $\dcat{C}_{\mrm{str}}(A, \cat{M})$.

Since the rows in (\ref{eqn:1916}) are standard triangles
(see Definition \ref{dfn:1140}), the objects $N$ and 
$N'$ are cones: $N = \opn{Cone}(\al)$ and $N' = \opn{Cone}(\al')$.
The commutativity up to homotopy of this diagram means 
that there is a degree $-1$ morphism 
$\th : L \to M'$ in $\dcat{C}(A, \cat{M})$ such that
$\al' \circ \phi = \psi \circ \al + \d(\th)$.

Define the morphism 
$\chi : N = \sbmat{ M \\[0.2em] \opn{T}(L) } \to 
N' = \sbmat{ M' \\[0.2em] \opn{T}(L') }$
by the matrix presentation 
\[ \chi :=  
\bmat{ \psi & \th \circ \opn{t}_{L}^{-1} \\[0.3em] 0 & \opn{T}(\phi) } . \]
An easy calculation shows that $\chi$ is a morphism in 
$\dcat{C}_{\mrm{str}}(A, \cat{M})$, and that there are equalities
$\opn{T}(\phi) \circ \ga = \ga' \circ \chi$
and 
$\chi \circ \be = \be' \circ \psi$. 
Therefore, when we apply the functor $\opn{P}$, and conjugate by the original 
isomorphism between (\ref{eqn:1915}) and the image of (\ref{eqn:1916}), 
we obtain a commutative diagram 
\[ \UseTips \xymatrix @C=8ex @R=6ex { 
\bar{L} 
\ar[r]^{\bar{\al}}
\ar[d]_{\bar{\phi}}
&
\bar{M}
\ar[r]^{\bar{\be}}
\ar[d]_{\bar{\psi}}
& 
\bar{N}
\ar[r]^{\bar{\ga}}
\ar[d]_{\bar{\chi}}
&
\opn{T}(\bar{L})
\ar[d]_{\opn{T}(\bar{\phi})}
\\
\bar{L}' 
\ar[r]^{\bar{\al}'}
&
\bar{M}'
\ar[r]^{\bar{\be}'}
& 
\bar{N}'
\ar[r]^{\bar{\ga}'}
&
\opn{T}(\bar{L}') 
} \]
in $\dcat{K}(A, \cat{M})$, where $\bar{\chi}$ is conjugate to
$\opn{P}(\chi)$. 

\medskip \noindent
(TR4): We may assume that the three given distinguished triangles are 
standard triangles in $\dcat{C}_{\mrm{str}}(A, \cat{M})$. 
Namely, we can assume that 
$\al : L \to M$ and $\be : M \to N$ are morphisms in 
$\dcat{C}_{\mrm{str}}(A, \cat{M})$; the DG modules $P, Q, R$ are 
$P = \opn{Cone}(\al)$, 
$Q = \opn{Cone}(\be \circ \al)$ and
$R = \opn{Cone}(\be)$; and the morphisms $\ga, \de, \ep$ in 
$\dcat{C}_{\mrm{str}}(A, \cat{M})$ are
$\ga = e_{\al}$, $\de = e_{\be \circ \al}$ and 
$\ep = e_{\be}$. All this in the notation of Subsection \ref{subsec:cone}. 

In matrix notation we have
\[ P = \bmat{M \\[0.1em] \opn{T}(L)} , \quad 
Q = \bmat{N \\[0.1em] \opn{T}(L)} , \quad 
R = \bmat{N \\[0.1em] \opn{T}(M)} . \]
We define the morphisms 
$\phi : P \to Q$ and $\psi : Q \to R$ in $\dcat{C}_{\mrm{str}}(A, \cat{M})$ 
by the matrix presentations 
\[ \phi := \bmat{\be & 0 \\[0.1em] 0 & \opn{id}_{\opn{T}(L)}} , \quad 
\psi := \bmat{\opn{id}_{N} & 0 \\[0.1em] 0 & \opn{T}(\al)} . \]
(We leave to the reader to verify that $\phi$ and $\psi$ commute with the 
differentials $\d_P$, $\d_Q$ and $\d_R$; this is just linear algebra, using 
the matrix presentations of the differentials of the cones from Definition 
\ref{dfn:1172}.)
Define the morphism 
$\rho : R \to \opn{T}(Q)$ in $\dcat{C}_{\mrm{str}}(A, \cat{M})$
to be the composition of the morphisms 
$R \to \opn{T}(M) \xar{\opn{T}(\ga)} \opn{T}(Q)$.
Then the diagram (\ref{eqn:5022}) in $\dcat{C}_{\mrm{str}}(A, \cat{M})$ is 
commutative.

It remains to prove that the triangle 
\begin{equation} \label{eqn:2485}
 P \xar{\bar{\phi}} Q \xar{\bar{\psi}} R \xar{\bar{\rho}} \opn{T}(P) 
\end{equation}
in $\dcat{K}(A, \cat{M})$ is distinguished. 
Let $C := \opn{Cone}(\phi)$; so we have a standard triangle
\begin{equation} \label{eqn:2486}
 P \xar{\phi} Q \xar{e_{\phi}} C \xar{p_{\phi}} \opn{T}(P)
\end{equation}
in $\dcat{C}_{\mrm{str}}(A, \cat{M})$. We are going to prove that the triangles 
(\ref{eqn:2485}) and (\ref{eqn:2486}) are isomorphic in 
$\dcat{K}(A, \cat{M})$, by producing an isomorphism 
$\bar{\chi} : C \iso R$ in $\dcat{K}(A, \cat{M})$ that makes the diagram 
\[ \UseTips  \xymatrix @C=8ex @R=6ex { 
P
\ar[r]^{\bar{\phi}}
\ar[d]^{\opn{id}}
&
Q
\ar[r]^{\bar{e}_{\phi}}
\ar[d]^{\opn{id}}
&
C
\ar[r]^{\bar{p}_{\phi}}
\ar[d]^{\bar{\chi}}
&
\opn{T}(P)
\ar[d]^{\opn{id}}
\\
P
\ar[r]^{\bar{\phi}}
&
Q
\ar[r]^{\bar{\psi}}
&
R
\ar[r]^{\bar{\rho}}
&
\opn{T}(P)
} \]
commutative. 

Here are the matrices for the object $C$, the morphism 
$\chi : C \to R$, and another morphism 
$\om: R \to C$, both in $\dcat{C}_{\mrm{str}}(A, \cat{M})$. 
\[ C = \bmat{N \\[0.1em] \opn{T}(L) \\[0.3em] \opn{T}(M) 
\\[0.1em] \opn{T}^2(L)} , \quad 
\chi := \bmat{\opn{id}_N & 0 & 0 & 0 
\\[0.1em] 0 & \opn{T}(\al) & \opn{id}_{\opn{T}(M)} & 0 } , \quad 
\om := \bmat{\opn{id}_N & 0 \\[0.1em] 0 & 0 
\\[0.1em] 0 & \opn{id}_{\opn{T}(M)} \\[0.1em] 0 & 0} . \]
Again, we leave it to the reader to check that $\chi$ and $\om$ commute with 
the differentials. It is easy to see that 
$\om \circ \psi = e_{\phi}$, 
$\rho \circ \chi = p_{\phi}$ and 
$\chi \circ \om = \opn{id}_R$. 

Finally we must find a homotopy between $\om \circ \chi$ and $\opn{id}_C$.
Consider the degree $-1$ endomorphism $\th$ of $C$~:
\[ \th := 
\bmat{0 & 0 & 0 & 0 
\\[0em] 0 & 0 & 0 & 0 
\\[0em] 0 & 0 & 0 & 0 
\\[0em] 0 & \opn{t}_{\opn{T}(L)} & 0 & 0} . \]
Then 
$\d_C \circ \th + \th \circ \d_C = \opn{id}_C - \, \om \circ \chi$.
\end{proof}

A full subcategory of a DG category is of course a DG category itself. 
Full additive (resp.\ triangulated) subcategories of additive
(resp.\ triangulated) categories were defined in 
Definition \ref{dfn:1025} (resp.\ \ref{dfn:1450}). 

\begin{cor} \label{cor:3015}
Let $\cat{C} \sub \dcat{C}(A, \cat{M})$ be a full subcategory satisfying these 
three conditions\tup{:}
\begin{itemize}
\rmitem{a} $\cat{C}_{\mrm{str}}$ is a full additive subcategory of 
$\dcat{C}_{\mrm{str}}(A, \cat{M})$. 

\rmitem{b} $\cat{C}$ is translation invariant, i.e.\
$M \in \cat{C}$ if and only $\opn{T}(M) \in \cat{C}$.

\rmitem{c} $\cat{C}$ is closed under standard cones, i.e.\ for every morphism 
$\phi$ in $\cat{C}_{\mrm{str}}$ the object $\opn{Cone}(\phi)$ belongs to 
$\cat{C}$.
\end{itemize}
Then $\cat{K} := \opn{Ho}(\cat{C})$
is a full triangulated subcategory of 
$\dcat{K}(A, \cat{M})$. 
\end{cor}

\begin{proof}
Each condition here implies the same numbered condition in Definition 
\ref{dfn:1450}. 
\end{proof}

Here is a partial converse to the corollary. 

\begin{prop} \label{prop:3015}
Suppose $\cat{K} \sub \dcat{K}(A, \cat{M})$ is a full triangulated subcategory. 
Let $\cat{C} \sub \dcat{C}(A, \cat{M})$ be the full subcategory on the set of 
objects $\opn{Ob}(\cat{K})$. Then the DG category $\cat{C}$ satisfies 
conditions \tup{(a)-(c)} of Corollary \tup{\ref{cor:3015}}, and 
$\cat{K} = \opn{Ho}(\cat{C})$.
\end{prop}

\begin{exer} \label{exer:3015}
Prove Proposition \ref{prop:3015}. 
\end{exer}

Here are two remarks about other sorts of triangulated categories.

\begin{rem} \label{rem:4911}
Let $\cat{E}$ be an exact category in the sense of D. Quillen \cite{Qu1}. 
Assume that $\cat{E}$ is a {\em Frobenius category}, i.e.\ $\cat{E}$ has 
enough injectives and enough projectives, and an object is injective if and only 
if it is projective. The 
{\em stable category}%
\index{Stable category}
of $\cat{E}$ is the 
quotient category $\ul{\cat{E}} := \cat{E} / \cat{E}_0$, where 
$\cat{E}_0 \sub \cat{E}$ is the two-sided ideal of morphisms 
factoring through a projective object. The set of objects of $\ul{\cat{E}}$ is 
the same as that of $\cat{E}$. D. Happel \cite{Hap0} has shown that 
$\ul{\cat{E}}$ is 
a triangulated category. 

By definition, an {\em algebraic triangulated category} is a 
triangulated category $\cat{K}$ which is equivalent to the stable category 
$\ul{\cat{E}}$ of some Frobenius category $\cat{E}$.
It turns out that all the derived categories studied in this book
(see Section \ref{sec:der-cat}) are in fact 
algebraic triangulated categories; see the paper \cite{Kel0} of B. Keller. 

Here are a couple of other examples of algebraic triangulated categories. 
Let $\cat{M}$ be an additive category, and let 
$\cat{E} := \dcat{C}_{\mrm{str}}(\cat{M})$, the strict category of complexes in 
$\cat{M}$. Put on $\cat{E}$ the exact structure in which the short exact 
sequences are the degreewise split sequences. Then $\cat{E}$ becomes a 
Frobenius category; the projective-injective objects in it are the contractible 
complexes. In this case the ideal $\cat{E}_0$ coincides with the null-homotopic 
morphisms in $\cat{E} = \dcat{C}_{\mrm{str}}(\cat{M})$, and hence 
$\ul{\cat{E}} = \dcat{K}(\cat{M})$, the homotopy category of $\cat{M}$. 

For the second example, suppose that $\K$ is a field and $A$ a noetherian 
$\K$-ring which is Gorenstein, i.e.\ $A$ has finite injective dimension 
as a left and as a right module over itself. The category of finitely generated 
$A$-modules is $\dcat{M}_{\mrm{f}}(A) = \cat{Mod}_{\mrm{f}} A$.
Let $\cat{E} \sub \dcat{M}_{\mrm{f}}(A)$ 
be the full subcategory on the modules $M$ such that 
$\opn{Ext}^i_A(M, A) = 0$ for all $i > 0$. These modules are sometimes called 
{\em maximal Cohen-Macaulay modules}. Put on $\cat{E}$ the exact structure in 
which the exact sequences are the exact sequences in $\dcat{M}_{\mrm{f}}(A)$
with all terms belonging to $\cat{E}$. Then $\cat{E}$ is a Frobenius category. 
The projective-injective objects of $\cat{E}$ are the modules in it that are
projective $A$-modules. Moreover, according to R.-O. Buchweitz \cite{Buch},
the triangulated category $\ul{\cat{E}}$ in this case is equivalent to the {\em 
stable derived category} of $A$, which is the Verdier quotient 
$\dcat{D}^{\mrm{b}}_{\mrm{f}}(A) / \dcat{D}(A)_{\mrm{perf}}$.
Here $\dcat{D}^{\mrm{b}}_{\mrm{f}}(A)$ is the derived category of bounded 
complexes with finitely generated cohomologies (see subsection 
\ref{subsec:thick-subcats}), $\dcat{D}(A)_{\mrm{perf}}$ is the full 
subcategory on the perfect complexes (see Section \ref{sec:perf-tilt-NC}), and 
Verdier quotients are explained in Remark \ref{rem:4220}. 
Note that the triangulated category 
$\dcat{D}^{\mrm{b}}_{\mrm{f}}(A) / \dcat{D}(A)_{\mrm{perf}}$
is called the {\em triangulated category of singularities}  by
D.O. Orlov \cite{Orl}, and this was later shortened to the {\em singularity 
category} of $A$.
\end{rem}

\begin{rem} \label{rem:4920}
In 1961, A. Dold and D. Puppe \cite{DoPu} tried to formalize the properties 
of the {\em stable homotopy category of topological spaces}. What they got is a 
theory almost identical to the theory of triangulated categories studied here 
(Definition \ref{dfn:1158}), with the exception of the octahedral axiom (TR4). 

In more modern terms we can consider the notion of a 
{\em stable model category}%
\index{Stable category}, 
that we now define. Let $\cat{C}$ be a pointed Quillen model category, meaning 
that it has a zero object. One defines a {\em suspension functor} $\Sigma$ and a 
{\em loop functor} $\Om$ on $\cat{C}$ using pushout and pullback diagrams 
involving the zero object. Let $\cat{D}$ be the {\em homotopy category} of 
$\cat{C}$. (In the terminology of our book, $\cat{D}$ is actually the {\em 
derived category} of $\cat{C}$, because it is gotten by inverting the weak 
equivalences in $\cat{C}$; see Sections 
\ref{sec:loc-cats}-\ref{sec:der-funcs}.) If the functors 
$\Sigma$ and $\Om$ induce auto-equivalences on $\cat{D}$, quasi-inverse to each 
other, then $\cat{C}$ is called a stable model category. In this case it can be 
shown that $\cat{D}$ is a triangulated category, with translation functor 
$\opn{T} := \Sigma$.
A more general notion is that of a {\em stable infinity category}
(see \cite{Lur2} or \cite{Hi2}). 

A {\em topological triangulated category} is, by definition, a triangulated 
category $\cat{K}$ which is equivalent to a full triangulated subcategory 
of the homotopy category $\cat{D}$ of a stable model category $\cat{C}$, as 
above. For more information on these topics see \cite{Ho1}, \cite{Sch}, 
\cite{ShSc}, \cite{Lur2} and \cite{CaSt}. 
\end{rem}

% trying to avoid a lot of blank space; new section forces a page break here
\newpage



\mysubsection{From DG Functors to Triangulated Functors} 
\label{subsec:DG-to-tri-fun}

We now add a second DG ring $B$, and a second abelian category $\cat{N}$.
DG functors were introduced in Subsection \ref{subsec:DGFunc}.

Consider a DG functor 
$F  : \dcat{C}(A, \cat{M}) \to \dcat{C}(B, \cat{N})$.
From Theorem \ref{thm:1150} we know that the translation isomorphism $\tau_F$ 
is an isomorphism of DG functors
$\tau_F : F \circ \opn{T}_{A, \cat{M}} \iso 
\opn{T}_{B, \cat{N}} \circ \, F$.
Therefore, when we pass to the homotopy categories, 
and writing $\bar{F} := \opn{Ho}(F)$
and $\bar{\tau}_F := \opn{Ho}(\tau_F)$,
we get a T-additive functor 
$(\bar{F}, \bar{\tau}_F): \dcat{K}(A, \cat{M}) \to \dcat{K}(B, \cat{N})$.

If
$G : \dcat{C}(A, \cat{M}) \to \dcat{C}(B, \cat{N})$
is another DG functor, then we have another T-additive functor 
$(\bar{G}, \bar{\tau}_G) :  \dcat{K}(A, \cat{M}) \to \dcat{K}(B, \cat{N})$.
And if $\eta : F \to G$ is a strict morphism of DG functors, then there is a 
morphism of additive functors 
$\bar{\eta} := \opn{Ho}(\eta) : \bar{F} \to \bar{G}$.
This notation will be used in the next theorem. 

\begin{thm} \label{thm:1265}
Let $A$ and $B$ be DG rings, let $\cat{M}$ and $\cat{N}$ be abelian categories, 
and let 
$F : \dcat{C}(A, \cat{M}) \to \dcat{C}(B, \cat{N})$
be a DG functor. 
\index{Differential graded functor}
\begin{enumerate}
\item The T-additive functor  
\[ (\bar{F}, \bar{\tau}_F) :  \dcat{K}(A, \cat{M}) \to \dcat{K}(B, \cat{N}) \]
is a triangulated functor. 
\index{Triangulated functor}

\item Suppose 
$G : \dcat{C}(A, \cat{M}) \to \dcat{C}(B, \cat{N})$
is another DG functor, and $\eta : F \to G$ is a strict morphism of DG functors.
\index{Differential graded functor! strict morphism of {\indash}s} 
Then 
\[ \bar{\eta} : (\bar{F}, \bar{\tau}_F) \to (\bar{G}, \bar{\tau}_G) \]
is a morphism of triangulated functors. 
\index{Triangulated functor! morphism of {\indash}s}
\end{enumerate}
\end{thm}

\begin{proof} \mbox{}

\smallskip \noindent 
(1) Take a distinguished triangle 
$L \xar{\bar{\al}} M \xar{\bar{\be}} N \xar{\bar{\ga}} \opn{T}(L)$
in $\dcat{K}(A, \cat{M})$. Since we are only interested in triangles up to 
isomorphism, we can assume that this is the image under the functor $\opn{P}$ 
of a standard triangle  
$L \xar{\al} M \xar{\be} N \xar{\ga} \opn{T}(L)$
in $\dcat{C}_{\mrm{str}}(A, \cat{M})$. According to Theorem \ref{thm:1185}
and Corollary \ref{cor:1915}, 
there is a  standard triangle  
$L' \xar{\al'} M' \xar{\be'} N' \xar{\ga'} \opn{T}(L')$
in $\dcat{C}_{\mrm{str}}(B, \cat{N})$, and a commutative diagram 
\[ \UseTips  \xymatrix @C=6ex @R=6ex { 
F(L) 
\ar[r]^{F(\al)}
\ar[d]_{\phi}
&
F(M) 
\ar[r]^{F(\be)}
\ar[d]_{\psi}
& 
F(N)
\ar[rr]^(0.44){\tau_{F, L} \, \circ \, F(\ga)}
\ar[d]_{\chi}
&
&
\opn{T}( F(L))
\ar[d]_{\opn{T}(\phi)}
\\
L' 
\ar[r]^{\al'}
&
M'
\ar[r]^{\be'}
& 
N'
\ar[rr]^{\ga'}
&
&
\opn{T}(L') 
} \]
in $\dcat{C}_{\mrm{str}}(B, \cat{N})$,
in which the vertical arrows are isomorphisms. (Actually, we can take 
$L' = F(L)$, $\phi = \opn{id}_{F(L)}$, etc.) 
After applying the functor $\opn{P}$ to this diagram, we see that the condition 
in Definition \ref{dfn:1276}(1) is satisfied. 

\medskip \noindent
(2) By Definition \ref{dfn:1276}(2), all we need to prove that $\bar{\eta}$ is 
a morphism of T-add\-itive functors. Let's use these abbreviations:
$\cat{K} := \dcat{K}(A, \cat{M})$ and 
$\cat{L} := \dcat{K}(B, \cat{N})$. 
Definition \ref{dfn:1160} requires that for every object 
$M \in \cat{K}$ the diagram 
\begin{equation} \label{eqn:3215}
\UseTips  \xymatrix @C=8ex @R=6ex {
\bar{F}(\opn{T}_{\cat{K}}(M))
\ar[r]^{\bar{\tau}_{F, M}}
\ar[d]_{\bar{\eta}_{\opn{T}_{\cat{K}}(M)}}
&
\opn{T}_{\cat{L}}(\bar{F}(M))
\ar[d]^{\opn{T}_{\cat{L}}(\bar{\eta}_M)}
\\
\bar{G}(\opn{T}_{\cat{K}}(M))
\ar[r]^{\bar{\tau}_{G, M}}
&
\opn{T}_{\cat{L}}(\bar{G}(M)) \ . 
} 
\end{equation}
in $\cat{L}$ is commutative.
Going back to Definitions \ref{dfn:1171} and \ref{dfn:1150}  we see that 
$\opn{T}_{\cat{L}}(\eta_M) = \opn{t}_{G(M)} \circ \, \eta_M \circ 
\opn{t}_{F(M)}^{-1}$,
$\tau_{F, M} = \opn{t}_{F(M)} \circ \, F(\opn{t}_M^{-1})$
and 
$\tau_{G, M} = \opn{t}_{G(M)} \circ G(\opn{t}_M^{-1})$, 
as morphisms in  $\dcat{C}(B, \cat{N})$.
Thus the path in (\ref{eqn:3215}) that starts by going right is represented by 
the morphism 
\[ \begin{aligned}
&
\opn{T}_{\cat{L}}(\eta_M) \circ \tau_{F, M} =
\opn{t}_{G(M)} \circ \, \eta_M \circ \opn{t}_{F(M)}^{-1} \circ 
\opn{t}_{F(M)} \circ \, F(\opn{t}_M^{-1})
\\
& \quad 
= \opn{t}_{G(M)} \circ \, \eta_M \circ F(\opn{t}_M^{-1})  
\end{aligned} \]
in $\dcat{C}_{\mrm{str}}(B, \cat{N})$, and other path is represented by 
$\opn{t}_{G(M)} \circ \, G(\opn{t}_M^{-1}) \circ 
\eta_{\lms \opn{T}_{\cat{K}}(M)}$.
Because $\opn{t}_{G(M)}$ is an isomorphism (of degree $-1$) in 
$\dcat{C}(B, \cat{N})$, it suffices to prove that 
$\eta_M \circ F(\opn{t}_M^{-1}) =  G(\opn{t}_M^{-1}) \circ
\eta_{\opn{T}_{\cat{K}}(M)}$
in $\dcat{C}(B, \cat{N})$; namely that the diagram 
\[ \UseTips  \xymatrix @C=8ex @R=6ex {
F(\opn{T}_{\cat{K}}(M))
\ar[r]^(0.55){F(\opn{t}_M^{-1})}
\ar[d]_{\eta_{\lms \opn{T}_{\cat{K}}(M)}}
&
F(M)
\ar[d]^{\eta_M}
\\
G(\opn{T}_{\cat{K}}(M))
\ar[r]^(0.55){G(\opn{t}_M^{-1})}
&
G(M)
} \]
is commutative. This is true because $\eta :F \to G$ is a strict morphism of DG 
functors, and here it acts on the degree $1$ morphism 
$\opn{t}_M^{-1} : \opn{T}_{\cat{K}}(M) \to M$
in $\dcat{C}(B, \cat{N})$.
\end{proof}

\begin{cor} \label{cor:1265}
For every integer $k$, the pair
$\bigl( \opn{T}^k, (-1)^k \cd \opn{id}_{\opn{T}^{k + 1}} \bigr)$
is a triangulated functor from $\dcat{K}(A, \cat{M})$ to itself. 
\end{cor}

\begin{proof}
Combine Theorem \ref{thm:1265} and Proposition \ref{prop:1300}. 
\end{proof}

\begin{rem} \label{rem:1475}
In \cite{BoKa}, A.I. Bondal and M.M. Kapranov introduced the concept of 
{\em pretriangulated DG category}%
\index{Differential graded category! pretriangulated}.
This is a DG category $\cat{C}$ for which the 
homotopy category $\opn{Ho}(\cat{C})$ is canonically triangulated (the details 
of the definition are too complicated to mention here). Our DG categories 
$\dcat{C}(A, \cat{M})$ are pretriangulated in the sense of \cite{BoKa}; but 
they have a lot more structure (e.g.\ the objects have cohomologies too). 

Suppose $\cat{C}$ and $\cat{C}'$ are pretriangulated DG categories. 
In \cite{BoKa} there is a (rather complicated) definition of {\em pre-exact DG 
functor} $F : \cat{C} \to \cat{C}'$. It is stated there that if $F$ is a 
pre-exact DG functor, then 
$\opn{Ho}(F) : \opn{Ho}(\cat{C}) \to \opn{Ho}(\cat{C}')$
is a triangulated functor. This is analogous to our Theorem \ref{thm:1265}.
Presumably, Theorems \ref{thm:1150} and \ref{thm:1185} imply that every DG 
functor 
$F : \dcat{C}(A, \cat{M}) \to \dcat{C}(A', \cat{M}')$
is pre-exact in the sense of \cite{BoKa}; but we did not verify this.  
\end{rem}



\mysubsection{The Opposite Homotopy Category is Triangulated}
\label{subsec:opp-hom-triang}

Here we introduce a canonical triangulated structure on the opposite homotopy 
category $\dcat{K}(A, \cat{M})^{\mrm{op}}$. This gives us a way to talk about 
contravariant triangulated functors whose source is a full subcategory of 
$\dcat{K}(A, \cat{M})$. Our solution is precise, but it is not totally 
satisfactory -- see Remark \ref{rem:2306}. 

We already gave a thorough treatment of contravariant DG functors in 
Subsection \ref{subsec:contrvar-dg-func}. In the previous subsection we 
explained precisely how to pass from DG functors to triangulated functors. 
In this subsection we treat the contravariant case.

As before, $A$ is a central DG $\K$-ring, and $\cat{M}$ is a $\K$-linear 
abelian category. The DG category of DG $A$-modules in $\cat{M}$ is 
$\dcat{C}(A, \cat{M})$. 
In Subsection \ref{subsec:contrvar-dg-func} we introduced the flipped DG 
category 
$\dcat{C}(A^{\mrm{op}}, \cat{M}^{\mrm{op}})$. 
In Theorem \ref{thm:2495} we had a canonical isomorphism of DG categories
\begin{equation} \label{eqn:3000}
\opn{Flip} : 
\dcat{C}(A, \cat{M})^{\mrm{op}} \iso 
\dcat{C}(A^{\mrm{op}}, \cat{M}^{\mrm{op}}) . 
\end{equation}

\begin{dfn} \label{dfn:2375}
The {\em flipped category}
\index{Homotopy category! flipped}
\index{1-K(A,M)flip@$\dcat{K}(A, \cat{M})^{\mrm{flip}}$}
of $\dcat{K}(A, \cat{M})$ is the triangulated category 
\[ \dcat{K}(A, \cat{M})^{\mrm{flip}} :=
\dcat{K}(A^{\mrm{op}}, \cat{M}^{\mrm{op}}) = 
\opn{Ho} \bigl( \dcat{C}(A^{\mrm{op}}, \cat{M}^{\mrm{op}}) \bigr) . \]
See Theorem \ref{thm:1255}. Its translation functor is denoted by 
$\opn{T}^{\mrm{flip}}$. 
\end{dfn}

Since 
$\opn{Ho} \bigl( \dcat{C}(A, \cat{M})^{\mrm{op}} \bigr) = 
\dcat{K}(A, \cat{M})^{\mrm{op}}$,
the isomorphism (\ref{eqn:3000}) 
induces a $\K$-linear isomorphism of categories 
\begin{equation} \label{eqn:2995}
\ol{\opn{Flip}} := \opn{Ho}(\opn{Flip}) :  
\dcat{K}(A, \cat{M})^{\mrm{op}} \iso \dcat{K}(A, \cat{M})^{\mrm{flip}} .
\end{equation}

\begin{dfn} \label{dfn:2995}
The category $\dcat{K}(A, \cat{M})^{\mrm{op}}$
\index{Homotopy category! opposite}
\index{1-K(A,M)op@$\dcat{K}(A, \cat{M})^{\mrm{op}}$}
is given the triangulated 
category structure induced from the flipped category 
$ \dcat{K}(A, \cat{M})^{\mrm{flip}}$
under the isomorphism (\ref{eqn:2995}). 
The translation functor of $\dcat{K}(A, \cat{M})^{\mrm{op}}$ is denoted by 
$\opn{T}^{\mrm{op}}$.
\end{dfn}

Thus 
\begin{equation} \label{eqn:3001}
\opn{T}^{\mrm{op}} = \ol{\opn{Flip}}^{\, -1} \circ 
\opn{T}^{\mrm{flip}} \circ \, \ol{\opn{Flip}} .
\end{equation}
The distinguished triangles of $\dcat{K}(A, \cat{M})^{\mrm{op}}$
are the images under $\ol{\opn{Flip}}^{\, -1}$ of the distinguished triangles 
of $ \dcat{K}(A, \cat{M})^{\mrm{flip}}$.
And tautologically, 
\begin{equation} \label{eqn:3002}
(\ol{\opn{Flip}}, \opn{id}) : 
\dcat{K}(A, \cat{M})^{\mrm{op}} \to  \dcat{K}(A, \cat{M})^{\mrm{flip}}
\end{equation}
is an isomorphism of triangulated categories. 

\begin{dfn} \label{dfn:3000}
Let $\cat{K} \sub \dcat{K}(A, \cat{M})$ be a full additive subcategory, and 
assume that $\cat{K}^{\mrm{op}}$ is a triangulated subcategory of 
$\dcat{K}(A, \cat{M})^{\mrm{op}}$.
Then we give $\cat{K}^{\mrm{op}}$ the triangulated structure induced from 
$\dcat{K}(A, \cat{M})^{\mrm{op}}$.
\index{1-Kop@$\cat{K}^{\mrm{op}}$} 
\end{dfn}

To say this a bit differently, the condition on $\cat{K}$ in the definition is 
that 
%\begin{equation} \label{eqn:3003}
$\cat{K}^{\mrm{flip}} := \ol{\opn{Flip}}(\cat{K}^{\mrm{op}})
\sub  \dcat{K}(A, \cat{M})^{\mrm{flip}}$ 
%\end{equation}
is a triangulated subcategory. The triangulated structure we put on 
$\cat{K}^{\mrm{op}}$ is such that 
$(\ol{\opn{Flip}}, \opn{id}) : (\cat{K}^{\mrm{op}}, \opn{T}^{\mrm{op}})
\to (\cat{K}^{\mrm{flip}}, \opn{T}^{\mrm{flip}})$
is an isomorphism of triangulated categories. 

\begin{rem} \label{rem:3170}
It could happen (though we have no example of it) that 
$\cat{K} \sub \dcat{K}(A, \cat{M})$ is a full triangulated subcategory, 
yet $\cat{K}^{\mrm{op}} \sub \dcat{K}(A, \cat{M})^{\mrm{op}}$ 
is not a triangulated subcategory; or vice versa. 
More on this issue in Remark \ref{rem:2306}. 
\end{rem}

The remark above notwithstanding, in many important cases, such as in 
Propositions \ref{prop:3050} and
\ref{prop:3051}, both $\cat{K} \sub \dcat{K}(A, \cat{M})$
and $\cat{K}^{\mrm{op}} \sub \dcat{K}(A, \cat{M})^{\mrm{op}}$
are triangulated. 

\begin{dfn} \label{dfn:2996}
Let $\cat{L}$ be a triangulated category, and let 
$\cat{K} \sub \dcat{K}(A, \cat{M})$ be a full subcategory 
such that $\cat{K}^{\mrm{op}} \sub \dcat{K}(A, \cat{M})^{\mrm{op}}$ is 
triangulated. A {\em contravariant triangulated functor}
\index{Triangulated functor! contravariant}
from $\cat{K}$ to $\cat{L}$ is, by definition, a triangulated functor 
$(F, \tau) : \cat{K}^{\mrm{op}} \to \cat{L}$
in the sense of Definition \ref{dfn:1276}, where $\cat{K}^{\mrm{op}}$ has the 
triangulated structure from Definition \ref{dfn:3000}.
\end{dfn}

\begin{thm} \label{thm:2995}
Let $A$ and $B$ be DG rings, and let $\cat{M}$ and $\cat{N}$ be abelian 
categories. Let 
$\cat{C} \sub \dcat{C}(A, \cat{M})$ be a full subcategory, and let 
$F : \cat{C}^{\mrm{op}} \to \dcat{C}(B, \cat{N})$
be a DG functor%
\index{Differential graded functor! contravariant}.
Consider the homotopy category 
$\cat{K} := \opn{Ho}(\cat{C}) \sub 
\opn{Ho} \bigl( \dcat{C}(A, \cat{M}) \bigr) = \dcat{K}(A, \cat{M})$
and the induced additive functor 
$\bar{F} := \opn{Ho}(F) : \cat{K}^{\mrm{op}} \to \dcat{K}(B, \cat{N})$.
Suppose that 
$\cat{K}^{\mrm{op}} \sub \dcat{K}(A, \cat{M})^{\mrm{op}}$ is triangulated.
Then there is a canonical translation isomorphism 
$\bar{\tau}$ such that 
\[ (\bar{F}, \bar{\tau}) : \cat{K}^{\mrm{op}} \to \dcat{K}(B, \cat{N}) \]
is a triangulated functor. 
\index{Triangulated functor! contravariant}
\end{thm}

\begin{proof}
Define 
\[ \cat{C}^{\mrm{flip}} := \opn{Flip}(\cat{C}^{\mrm{op}}) \sub
\dcat{C}(A^{\mrm{op}}, \cat{M}^{\mrm{op}}) =
\dcat{C}(A, \cat{M})^{\mrm{flip}} . \]
This is a DG category, and 
$\opn{Flip} : \cat{C}^{\mrm{op}} \to \cat{C}^{\mrm{flip}}$
is an isomorphism of DG categories. Also
$\opn{Ho} (\cat{C}^{\mrm{flip}}) = \cat{K}^{\mrm{flip}}$.
There is a DG functor
\[ F^{\mrm{flip}} := F \circ \opn{Flip}^{-1} : \cat{C}^{\mrm{flip}} 
\to \dcat{C}(B, \cat{N}) , \]
and it induces an additive functor 
\[ \bar{F}^{\mrm{flip}} := \opn{Ho}(F^{\mrm{flip}}) : 
\cat{K}^{\mrm{flip}} \to \dcat{K}(B, \cat{N}) . \]
By Theorem \ref{thm:1265} there is a translation isomorphism 
$\bar{\tau}^{\mrm{flip}}$ such that 
\[ (\bar{F}^{\mrm{flip}}, \bar{\tau}^{\mrm{flip}}) : 
\cat{K}^{\mrm{flip}} \to \dcat{K}(B, \cat{N})  \]
is a triangulated functor.
Using formula (\ref{eqn:3001}) we obtain these equalities and isomorphisms:
\[ \begin{aligned}
& \bar{F} \circ \opn{T}^{\mrm{op}} = 
\bar{F} \circ \ol{\opn{Flip}}^{\, -1} \circ \opn{T}^{\mrm{flip}} \circ \, 
\ol{\opn{Flip}} 
= \bar{F}^{\mrm{flip}} \circ \opn{T}^{\mrm{flip}} \circ \, 
\ol{\opn{Flip}} 
\\
& \quad 
\xar{\bar{\tau}^{\mrm{flip}}}
\opn{T}_{\dcat{K}(B, \cat{N})} \circ \, \bar{F}^{\mrm{flip}} \, \circ  \, 
\ol{\opn{Flip}} 
= \opn{T}_{\dcat{K}(B, \cat{N})} \circ \, \bar{F} . 
\end{aligned} \]
We define $\bar{\tau}$ to be the composed isomorphism. By construction 
the pair $(\bar{F}, \bar{\tau})$ is a triangulated functor. 
\end{proof}

\begin{exer} \label{exer:3220}
Let $\cat{C} \sub \dcat{C}(A, \cat{M})$ be a full subcategory s.t.\ 
$\cat{K} := \opn{Ho}(\cat{C})$ is a triangulated subcategory of 
$\dcat{K}(A, \cat{M})$, and let 
$G : \cat{C} \to \dcat{C}(B, \cat{N})^{\mrm{op}}$
be a DG functor. Write $\bar{G} := \opn{Ho}(G)$. 
Show that there is a translation isomorphism $\bar{\nu}$ s.t.\ 
$(\bar{G}, \bar{\nu}) : \cat{K} \to \dcat{K}(B, \cat{N})^{\mrm{op}}$
is a triangulated functor. 
(Hint: study the proof of the theorem above.) 
\end{exer}

\begin{rem} \label{rem:2306}
Let $\cat{K}$ be a triangulated category. There is a way to introduce a 
triangulated structure on the opposite category $\cat{K}^{\mrm{op}}$. 
\index{Triangulated category! opposite}
The translation functor is 
$\opn{T}^{\mrm{op}} := \opn{Op} \circ \opn{T}^{-1} \circ \opn{Op}^{-1}$.
The distinguished triangles are defined to be 
\[ N \xar{\opn{Op}(\be)} M \xar{\opn{Op}(\al)} L 
\xar{\opn{Op}(- \opn{T}^{-1}(\ga))}  \opn{T}^{\mrm{op}}(N) , \]
where 
$L \xar{\al} M \xar{\be} N \xar{\ga} \opn{T}(L)$
goes over all distinguished triangles in $\cat{K}$.
See \cite[Remark 10.1.10(ii)]{KaSc2}. 
This approach has an advantage: it allows to talk about contravariant 
triangulated functors 
\index{Triangulated functor! contravariant}
$\cat{K} \to \cat{L}$ without any restriction on the 
category $\cat{K}$. This is in contrast to our quite restricted Definition 
\ref{dfn:2996}. The disadvantage of the approach presented in this paragraph   
is that there is no easy way to tie it with the DG theory (as opposed to our 
Theorems \ref{thm:1265} and \ref{thm:2995}). 

It would have been very pleasing if in the case 
$\cat{K} = \dcat{K}(A, \cat{M})$, the triangulated structure on 
$\dcat{K}(A, \cat{M})^{\mrm{op}}$ presented in the paragraph above would 
coincide with the triangulated structure from Definition \ref{dfn:2995}.
However, our calculations seem to indicate otherwise. 

Since in this book we are only interested in triangulated functors that are 
of DG origin (covariant, as in Theorem \ref{thm:1265}, or contravariant, as in 
Theorem \ref{thm:2995}), we decided to adopt Definition \ref{dfn:2996}.
It is a reliable, yet somewhat awkward approach. For instance, it requires us 
to perform particular calculations for the composition of contravariant 
triangulated functors, as in Lemma \ref{lem:2156}. 
\end{rem}






\cleardoublepage
\mysection{Localization of Categories} \label{sec:loc-cats}

\AYcopyright

This section is devoted to the general theory of {\em Ore 
localization of categories}. We are given a category $\cat{A}$ and a 
multiplicatively closed set of morphisms $\cat{S} \sub \cat{A}$.
The localized category $\cat{A}_{\cat{S}}$, gotten by formally inverting the 
morphisms in $\cat{S}$, always exists. The goal is to have a presentation of the
morphisms of $\cat{A}_{\cat{S}}$ as left or right fractions. 

In Section \ref{sec:der-cat} we shall apply the results of this section to 
triangulated categories. 


\mysubsection{The Formalism of Localization}
We will start with a category $\cat{A}$, without even assuming it is 
linear. Still we use the notation $\cat{A}$, because it will be suggestive to 
think about a linear category $\cat{A}$ with a single object, which is just a 
ring $A$. The reason is that our localization procedure is the same as that in 
noncommutative ring theory -- even when the category is not linear, and it has 
multiple objects. In Subsection \ref{subsec:loc-lin} we treat linear 
categories, and thus we recover the ring theoretic localization as a special 
case. 

The emphasis will be on morphisms rather than on objects. Thus it will be 
convenient to write 
$\cat{A}(M, N) := \opn{Hom}_{\cat{A}}(M, N)$
for $M, N \in \opn{Ob}(\cat{A})$. 
We sometimes use the notation $a \in \cat{A}$ for a morphism 
$a \in \cat{A}(M, N)$, leaving the objects implicit. When we write $b \circ a$ 
for $a, b \in \cat{A}$, we implicitly mean that these morphisms are 
composable. 

For heuristic purposes, we can think of $\cat{A}$ as a linear category (e.g.\ 
living inside some category of modules), with objects $M, N, \ldots$.
For any given object $M$, we then have a genuine ring 
$\cat{A}(M) := \cat{A}(M, M)$. 

\begin{dfn} \label{dfn:1405}
Let $\cat{A}$ be a category. A 
{\em multiplicatively closed set of morphisms}% 
\index{Multiplicatively closed set of morphisms}
in $\cat{A}$ is a subcategory $\cat{S} \subseteq \cat{A}$ such that 
$\opn{Ob}(\cat{S}) = \opn{Ob}(\cat{A})$.
\end{dfn}

In other words, for any pair of objects $M, N \in \cat{A}$ 
there is a subset $\cat{S}(M, N) \subseteq \cat{A}(M, N)$, 
such that $\opn{id}_M \in \cat{S}(M, M)$, and such that for any
$s \in \cat{S}(L, M)$ and $t \in \cat{S}(M, N)$, the composition
$t \circ s \in \cat{S}(L, N)$. 

Using our shorthand, we can write the definition like this: 
$\opn{id}_M \in \cat{S}$,
and $s, t \in \cat{S}$ implies $t \circ s \in \cat{S}$. 

If $A = \cat{A}$ is a single object linear category, namely a ring, then 
$S = \cat{S}$ is a multiplicatively closed set in the sense of ring theory. 

There are various notions of localization in the literature.  We restrict 
attention to two of them. Here is the first: 

\begin{dfn} \label{dfn:1406}
Let $\cat{S}$ be a multiplicatively closed set of morphisms in a category 
$\cat{A}$. A {\em localization} 
\index{Localization! of a category}
of $\cat{A}$ with respect to $\cat{S}$  
is a pair $(\cat{A}_{\cat{S}}, \opn{Q})$, consisting of a category 
$\cat{A}_{\cat{S}}$ and a functor 
$\opn{Q} : \cat{A} \to \cat{A}_{\cat{S}}$, called the localization functor,
having the following properties:
\begin{itemize}[leftmargin=32pt] 
\item[(Loc1)] There is equality 
$\opn{Ob}(\cat{A}_{\cat{S}}) = \opn{Ob}(\cat{A})$, and $\opn{Q}$ is the 
identity 
on objects.

\rmitem{Loc2}
For every $s \in \cat{S}$, the morphism $\opn{Q}(s) \in \cat{A}_{\cat{S}}$
is invertible (i.e.\ it is an isomorphism). 

\rmitem{Loc3} Suppose $\cat{B}$ is a category, and 
$F : \cat{A} \to \cat{B}$
is a functor such that $F(s)$ is invertible for every 
$s \in \cat{S}$. 
Then there is a unique functor 
$F_{\cat{S}} : \cat{A}_{\cat{S}} \to \cat{B}$
such that 
$F_{\cat{S}} \circ \opn{Q} = F$ 
as functors $\cat{A} \to \cat{B}$.
\end{itemize}
\end{dfn}

In a commutative diagram:
\[ \UseTips  \xymatrix @C=10ex @R=6ex  {
\cat{S}
\ar[r]^{\mrm{inc}}
&
\cat{A}
\ar@{-->}[r]^{F}="q" 
\ar[d]_{\opn{Q}}
& 
\cat{B}
\\
&
\cat{A}_{\cat{S}}
\ar@{-->}[ur]_{F_{\cat{S}}} ^(0.5){}="f"
} \]

\begin{exa} \label{exa:2310}
Suppose $\cat{A}$ has a single object, say $x$. Here is what the definition 
above says in this case. We have a ring 
$A := \cat{A}(x, x)$ and a multiplicatively closed set
$S := \cat{S}(x, x) \sub A$. The localization of $A$ with respect to $S$ is a 
ring $A_S := \cat{A}_{\cat{S}}(x, x)$, with a ring homomorphism 
$\opn{Q} : A \to A_S$, such that $\opn{Q}(S) \sub (A_S)^{\times}$,
the group of invertible elements of $A_S$. 
Condition (Loc3) says that any ring homomorphism
$F : A \to B$ such that $F(S) \sub B^{\times}$ factors uniquely through 
$A_S$. 
\end{exa}

\begin{thm} \label{thm:2310}
Let $\cat{S}$ be a multiplicatively closed set of morphisms in a category 
$\cat{A}$. A localization 
\index{Localization! of a category}
$(\cat{A}_{\cat{S}}, \opn{Q})$ of 
$\cat{A}$ with respect to $\cat{S}$, in the sense of Definition 
\tup{\ref{dfn:1406}}, exists, and it is unique up to a unique isomorphism. 
\end{thm}

\begin{proof}
Uniqueness: Suppose $(\cat{A}'_{\cat{S}}, \opn{Q}')$ is another localization.
By condition (Loc3), for both localizations, there are unique functors 
$G : \cat{A}_{\cat{S}} \to \cat{A}'_{\cat{S}}$
and 
$G' : \cat{A}'_{\cat{S}} \to \cat{A}_{\cat{S}}$,
that satisfy 
$\opn{Q}' = G \circ \opn{Q}$ and 
$\opn{Q} = G' \circ \opn{Q}'$. These functors must be the identity on 
objects, and are inverses to each other. 

Existence: Here we encounter the set theoretic issue alluded to in Subsection 
\ref{subsec:set-theor}. The category $\cat{A}_{\cat{S}}$ that we will construct 
might have to belong to a bigger universe than that in which $\cat{A}$ lives; 
this is because of its large morphism sets. 
We proceed as follows. Let $\cat{B}$ be the free category on the set of objects 
$\opn{Ob}(\cat{A})$, with set of generating morphisms 
$\cat{A} \sqcup \cat{S}^{\mrm{op}}$.
Thus, the morphisms in $\cat{B}$ are the finite
sequences of composable morphisms in $\cat{A} \sqcup \cat{S}^{\mrm{op}}$, and 
composition is concatenation. Inside $\cat{B}$ we have the congruence 
(two-sided 
ideal) $\cat{N} \sub \cat{B}$ generated by the relations of $\cat{A}$, 
the relations of $\cat{S}^{\mrm{op}}$, and the 
relations $\opn{Op}(s) \circ s = \opn{id}$ and 
$s \circ \opn{Op}(s) = \opn{id}$ for all $s \in \cat{S}$. 
The quotient category 
$\cat{A}_{\cat{S}} := \cat{B} / \cat{N}$ has the desired properties. 
Cf.\ \cite{GaZi} for more details. 
\end{proof}

\begin{prop} \label{prop:3030}
Let $\cat{S}$ be a multiplicatively closed set in a category $\cat{A}$, 
let $\cat{B}$ be a category, 
and let $\opn{Q} : \cat{A} \to \cat{B}$ be a functor. 
Consider the multiplicatively closed set 
$\cat{S}^{\mrm{op}} \sub \cat{A}^{\mrm{op}}$ and the functor 
$\opn{Q}^{\mrm{op}} : \cat{A}^{\mrm{op}} \to \cat{B}^{\mrm{op}}$. 
The two conditions below are equivalent\tup{:}
\begin{itemize}
\rmitem{i} The pair $(\cat{B}, \opn{Q})$ is a localization of 
$\cat{A}$ with respect to $\cat{S}$. 

\rmitem{ii} The pair $(\cat{B}^{\mrm{op}}, \opn{Q}^{\mrm{op}})$ is a  
localization of $\cat{A}^{\mrm{op}}$ with respect to $\cat{S}^{\mrm{op}}$. 
\end{itemize}
\end{prop}

\begin{exer} \label{exer:3031}
Prove Proposition \ref{prop:3030}. 
\end{exer}

Often the localization $\cat{A}_{\cat{S}}$ is of little value, because there is 
no practical way to describe the morphisms in it. This issue will be addressed 
in the next subsection. 




\mysubsection{Ore Localization}
There is a better notion of localization. 
The references here are \cite{RD}, \cite{GaZi}, \cite{We}, \cite{KaSc1}, 
\cite{Ste} and \cite{Row}. The first four references talk about localization of 
categories; and the last two talk about noncommutative rings. 
It seems that historically, this {\em noncommutative calculus of fractions} was 
discovered by Ore and Asano in ring theory, around 1930. There was progress in 
the categorical side, notably by Gabriel around 1960. 

In this subsection we mostly follow the treatment of \cite{Ste}; but we 
sometimes use diagrams instead of formulas.

\begin{dfn}  \label{dfn:1407}
Let $\cat{S}$ be a multiplicatively closed set of morphisms in a category 
$\cat{A}$. A {\em right Ore localization}
\index{Localization! Ore {\indash} of a category}
of $\cat{A}$ with respect to 
$\cat{S}$ is a pair $(\cat{A}_{\cat{S}}, \opn{Q})$, consisting of a category 
$\cat{A}_{\cat{S}}$ and a functor 
$\opn{Q} : \cat{A} \to \cat{A}_{\cat{S}}$, having the following 
properties:
\begin{itemize}[leftmargin=30pt] 
\item[(RO1)] There is equality 
$\opn{Ob}(\cat{A}_{\cat{S}}) = \opn{Ob}(\cat{A})$, and $\opn{Q}$ is the 
identity on
objects.

\item[(RO2)]  For every $s \in \cat{S}$, the morphism $\opn{Q}(s) \in
\cat{A}_{\cat{S}}$
is an isomorphism. 

\item[(RO3)]  Every morphism $q \in \cat{A}_{\cat{S}}$ can be written as
$q = \opn{Q}(a) \circ \opn{Q}(s)^{-1}$ for some 
$a \in \cat{A}$ and $s \in \cat{S}$.

\item[(RO4)]  Suppose $a, b \in \cat{A}$ satisfy 
$\opn{Q}(a) = \opn{Q}(b)$. Then $a \circ s = b \circ s$ for some 
$s \in \cat{S}$. 
\end{itemize}
\end{dfn}

The letters ``RO'' stand for ``right Ore''. 
We refer to the expression 
$q = \opn{Q}(a) \circ \opn{Q}(s)^{-1}$
as a {\em right fraction representation of $q$}. 
Here is the left sided version of this definition: 

\begin{dfn}  \label{dfn:2310}  
Let $\cat{S}$ be a multiplicatively closed set of morphisms in a category 
$\cat{A}$. A {\em left Ore localization}
\index{Localization! Ore {\indash} of a category}
of $\cat{A}$ with respect to 
$\cat{S}$ is a pair $(\cat{A}_{\cat{S}}, \opn{Q})$, consisting of a category 
$\cat{A}_{\cat{S}}$ and a functor 
$\opn{Q} : \cat{A} \to \cat{A}_{\cat{S}}$,
having the following properties:
\begin{itemize}[leftmargin=30pt] 
\item[(LO1)] There is equality 
$\opn{Ob}(\cat{A}_{\cat{S}}) = \opn{Ob}(\cat{A})$, and $\opn{Q}$ is the 
identity on
objects.

\item[(LO2)]  For every $s \in \cat{S}$, the morphism $\opn{Q}(s) \in
\cat{A}_{\cat{S}}$
is an isomorphism. 

\item[(LO3)]  Every morphism $q \in \cat{A}_{\cat{S}}$ can be written as
$q = \opn{Q}(s)^{-1} \circ \opn{Q}(a)$ for some 
$a \in \cat{A}$ and $s \in \cat{S}$.

\item[(LO4)]  Suppose $a, b \in \cat{A}$ satisfy 
$\opn{Q}(a) = \opn{Q}(b)$. Then $s \circ a = s \circ b$
for some $s \in \cat{S}$. 
\end{itemize}
\end{dfn}

As in the right case, the letters ``LO'' stand for ``left Ore''; and we refer 
to the expression 
$q = \opn{Q}(s)^{-1} \circ \opn{Q}(a)$
as a {\em left fraction representation of $q$}. 

\begin{rem} \label{rem:1845}
The results to follow in this subsection will be stated for right Ore 
localizations only. They all have ``left'' versions, with identical proofs 
(just a matter of reversing some arrows or compositions),  and therefore they 
will be omitted. 
\end{rem}

To reinforce the last remark, we give: 

\begin{prop} \label{prop:1845}
Let $\cat{S}$ be a multiplicatively closed set in a category $\cat{A}$, 
let $\cat{B}$ be a category, 
and let $\opn{Q} : \cat{A} \to \cat{B}$ be a functor. 
Consider the multiplicatively closed set 
$\cat{S}^{\mrm{op}} \sub \cat{A}^{\mrm{op}}$ and the functor 
$\opn{Q}^{\mrm{op}} : \cat{A}^{\mrm{op}} \to \cat{B}^{\mrm{op}}$. 
The two conditions below are equivalent\tup{:}
\begin{itemize}
\rmitem{i} The pair $(\cat{B}, \opn{Q})$ is a left Ore localization of 
$\cat{A}$ with respect to $\cat{S}$. 

\rmitem{ii} The pair $(\cat{B}^{\mrm{op}}, \opn{Q}^{\mrm{op}})$ is a right Ore 
localization of $\cat{A}^{\mrm{op}}$ with respect to $\cat{S}^{\mrm{op}}$. 
\end{itemize}
\end{prop}

\begin{exer} \label{exer:1845}
Prove Proposition \ref{prop:1845}. 
\end{exer}

\begin{lem} \label{lem:3}
Let $(\cat{A}_{\cat{S}}, \opn{Q})$ be a right Ore localization, 
let $a_1, a_2 \in \cat{A}$ and $s_1, s_2 \in \cat{S}$.
The following conditions are equivalent\tup{:}
\begin{enumerate}
\rmitem{i} 
$\opn{Q}(a_1) \circ \opn{Q}(s_1)^{-1} =  
\opn{Q}(a_2) \circ \opn{Q}(s_2)^{-1}$
in $\cat{A}_{\cat{S}}$. 

\rmitem{ii} There are $b_1, b_2 \in \cat{A}$ s.t.\ 
$a_1 \circ b_1 = a_2 \circ b_2$,
and $s_1 \circ b_1  = s_2 \circ b_2 \in \cat{S}$.
\end{enumerate}
\end{lem}

\begin{proof} \mbox{}

\smallskip \noindent
(ii) $\Rightarrow$ (i):
Since $\opn{Q}(s_i)$ and $\opn{Q}(s_i \circ b_i)$ are invertible, it 
follows that 
$\opn{Q}(b_i)$ are invertible. So 
\[ \begin{aligned}
& \opn{Q}(a_1) \circ \opn{Q}(s_1)^{-1} = \opn{Q}(a_1) \circ 
\opn{Q}(b_1) \circ \opn{Q}(b_1)^{-1} \circ \opn{Q}(s_1)^{-1} 
\\
& \quad = 
\opn{Q}(a_2) \circ \opn{Q}(b_2) \circ \opn{Q}(b_2)^{-1} \circ 
\opn{Q}(s_2)^{-1} =  \opn{Q}(a_2) \circ \opn{Q}(s_2)^{-1} .
\end{aligned} \]

\medskip \noindent
(i) $\Rightarrow$ (ii):
By property (RO3) there are $c \in \cat{A}$ and $u \in \cat{S}$ s.t.\
\begin{equation} \label{eqn:1425}
\opn{Q}(s_2)^{-1} \circ \opn{Q}(s_1) = \opn{Q}(c) \circ \opn{Q}(u)^{-1} . 
\end{equation}
Rewriting this equation we get 
\begin{equation} \label{eqn:1426}
\opn{Q}(s_1 \circ u) = \opn{Q}(s_2 \circ c) . 
\end{equation}
It is given that 
$\opn{Q}(a_1) =  \opn{Q}(a_2) \circ \opn{Q}(s_2)^{-1} \circ \opn{Q}(s_1)$.
Plugging (\ref{eqn:1425}) into it we obtain 
$\opn{Q}(a_1) =  \opn{Q}(a_2) \circ \opn{Q}(c) \circ \opn{Q}(u)^{-1}$.
Rearranging this equation we get 
$\opn{Q}(a_1 \circ u) = \opn{Q}(a_2 \circ c)$. 
By property (RO4) there is 
$v \in \cat{S}$ s.t.\ 
$a _1 \circ u \circ v = a_2 \circ c \circ v$.
Likewise, from equation (\ref{eqn:1426}) and property (RO4), there is 
$v' \in \cat{S}$ s.t.\ 
$s_1 \circ u \circ v' = s_2 \circ c \circ v'$. 

Again using property (RO3), there are $d \in \cat{A}$ and $w \in \cat{S}$ s.t.\
$\opn{Q}(v)^{-1} \circ \opn{Q}(v') = \opn{Q}(d) \circ \opn{Q}(w)^{-1}$.
Rearranging we get 
$\opn{Q}(v' \circ w) = \opn{Q}(v \circ d)$.
By property (RO4) there is $w' \in \cat{S}$ s.t.\ 
$v' \circ w \circ w' = v \circ d \circ w'$.
Define 
$b_1 := u \circ v \circ  d \circ w'$ 
and
$b_2 := c \circ v \circ d \circ w'$.
Then 
\[ \begin{aligned}
& s_1 \circ b_1 = s_1 \circ u \circ v \circ  d \circ w' = 
s_1 \circ u \circ v' \circ w \circ w' 
\\
& \quad = s_2 \circ c \circ v' \circ w \circ w' = s_2 \circ b_2 , 
\end{aligned} \]
and it is in $\cat{S}$. 
Also 
\[ a_1 \circ b_1 = a_1 \circ u \circ v \circ  d \circ w' =
a_2 \circ c \circ v  \circ  d \circ w' = a_2 \circ b_2 . \qedhere \]
\end{proof}


\begin{prop} \label{prop:103}
A right Ore localization $(\cat{A}_{\cat{S}}, \opn{Q})$ is a localization in 
the sense of Definition \tup{\ref{dfn:1406}}. 
\end{prop}

\begin{proof}
Say  $\cat{B}$ is a category, and 
$F : \cat{A} \to \cat{B}$
is a functor such that $F(s)$ is an isomorphism for every 
$s \in \cat{S}$. 

The uniqueness of a functor
$F_{\cat{S}} : \cat{A}_{\cat{S}} \to \cat{B}$
satisfying $F_{\cat{S}} \circ \opn{Q} = F$  is clear from property (RO3).
We have to prove existence.

Define $F_{\cat{S}}$ to be $F$ on objects.
For a morphism $q$ in $\cat{A}_{\cat{S}}$, property (RO3) says that
there is a right fraction presentation 
$q = \opn{Q}(a_1) \circ \opn{Q}(s_1)^{-1}$,
with $a_1 \in \cat{A}$ and $s_1 \in \cat{S}$.
Let $F_{\cat{S}}(q) := F(a_1) \circ F(s_1)^{-1}$ in $\cat{B}$.
We have to prove that this is well-defined, i.e.\ it does not depend on the 
presentation. So suppose that 
$q = \opn{Q}(a_2) \circ \opn{Q}(s_2)^{-1}$
is another right fraction presentation of $q$.
Let $b_1, b_2 \in \cat{A}$ be as in Lemma \ref{lem:3}. 
Since $F(s_i)$ and $F(s_i \circ b_i)$ are invertible, then so is $F(b_i)$. 
We get
$F(a_2) = F(a_1) \circ F(b_1) \circ F(b_2)^{-1}$
and 
$F(s_2) = F(s_1) \circ F(b_1) \circ F(b_2)^{-1}$.
Hence 
$F(a_2) \circ F(s_2)^{-1}  = F(a_1) \circ F(s_1)^{-1}$, 
as required. 

It remains to prove that $F_{\cat{S}}$ is a functor. Since the identity 
$\opn{id}_M$ of the object $M$ in $\cat{A}_{\cat{S}}$ can be presented 
as 
$\opn{id}_M = \opn{Q}(\opn{id}_M) \circ \opn{Q}(\opn{id}_M)^{-1}$, 
we see that $F_{\cat{S}}(\opn{id}_M) = \opn{id}_{F(M)}$. 

Next let $q_1$ and $q_2$ be morphisms in $\cat{A}_{\cat{S}}$, such that the 
composition $q_2 \circ q_1$ exists (i.e.\ the target of $q_1$ is the source of 
$q_2$). We have to show that 
$F_{\cat{S}}(q_2 \circ q_1)$ equals
$F_{\cat{S}}(q_2) \circ F_{\cat{S}}(q_1)$.
Choose presentations 
$q_i = \opn{Q}(a_i) \circ \opn{Q}(s_i)^{-1}$, 
so that 
\begin{equation} \label{eqn:1421}
F_{\cat{S}}(q_2) \circ F_{\cat{S}}(q_1) = 
F(a_2) \circ F(s_2)^{-1} \circ F(a_1) \circ F(s_1)^{-1}  . 
\end{equation}
By property (RO3) there is a right fraction presentation 
\begin{equation} \label{eqn:1428}
\opn{Q}(s_2)^{-1} \circ \opn{Q}(a_1) = \opn{Q}(b) \circ \opn{Q}(t)^{-1} 
\end{equation}
for some $b \in \cat{A}$ and $t \in \cat{S}$. 
Because
$\opn{Q}(a_1 \circ t) = \opn{Q}(s_2 \circ b)$,
by (RO4) there is some $r \in \cat{S}$ such that 
$a_1 \circ t \circ r = s_2 \circ b \circ r$.
Therefore 
$F(a_1 \circ t \circ r) = F(s_2 \circ b \circ r)$,
which implies, by canceling the invertible morphism $F(r)$ and rearranging, that
\begin{equation} \label{eqn:1420}
F(s_2)^{-1} \circ F(a_1) = F(b) \circ F(t)^{-1}
\end{equation}
in $\cat{B}$. 

Let us continue. Using equation (\ref{eqn:1428}) we have  
\[ \begin{aligned}
& q_2 \circ q_1 = \opn{Q}(a_2) \circ \opn{Q}(s_2)^{-1} \circ 
\opn{Q}(a_1) \circ \opn{Q}(s_1)^{-1} 
\\ 
& \qquad 
= \opn{Q}(a_2) \circ \opn{Q}(b) \circ \opn{Q}(t)^{-1} \circ \opn{Q}(s_1)^{-1} =
\opn{Q}(a_2 \circ b) \circ \opn{Q}(s_1 \circ t)^{-1} . 
\end{aligned} \]
Using this presentation of $q_2 \circ q_1$, and the equality (\ref{eqn:1420}), 
we obtain 
\[ \begin{aligned}
& F_{\cat{S}}(q_2 \circ q_1) = 
F(a_2 \circ b) \circ F(s_1 \circ t)^{-1} 
= F(a_2) \circ F(b) \circ F(t)^{-1} \circ F(s_1)^{-1}
\\
& \qquad 
= F(a_2) \circ F(s_2)^{-1} \circ F(a_1) \circ F(s_1)^{-1} . 
\end{aligned} \]
This is the same as (\ref{eqn:1421}). 
\end{proof}

\begin{cor} \label{cor:1405}
Let $\cat{S}$ be a multiplicatively closed set of morphisms in a category 
$\cat{A}$. Assume that these two assertions hold\tup{:}
\begin{itemize}
\item The pair $(\cat{B}, \opn{Q})$ is either a right or a left Ore 
localization of $\cat{A}$ with respect to $\cat{S}$.

\item The pair $(\cat{B}', \opn{Q}')$ is either a right or a left 
Ore localization of $\cat{A}$ with respect to $\cat{S}$.
\end{itemize}
Then there is a unique isomorphism of localizations 
$(\cat{B}, \opn{Q}) \cong (\cat{B}', \opn{Q}')$.
\end{cor}

\begin{proof}
By Proposition \ref{prop:103} (in its right or left versions, as the case may 
be), both $(\cat{B}, \opn{Q})$ and $(\cat{B}', \opn{Q}')$ 
are localizations in the sense of Definition \ref{dfn:1406}. Hence, by 
Theorem \ref{thm:2310}, there is a unique isomorphism 
$(\cat{B}, \opn{Q}) \cong (\cat{B}, \opn{Q}')$.
\end{proof}

\begin{dfn} \label{dfn:21}
Let $\cat{S}$ be multiplicatively closed set of morphisms in a  category 
$\cat{A}$. We say that $\cat{S}$ is a {\em right denominator set}
\index{Denominator set}
if it satisfies these two conditions:
\begin{enumerate}[leftmargin=30pt] 
\item[(RD1)] (Right Ore condition) Given $a \in \cat{A}$ and $s \in \cat{S}$,
there exist $b \in \cat{A}$ and $t \in \cat{S}$ such that
$a \circ t = s \circ b$.

\item[(RD2)] (Right cancellation condition) 
Given $a, b \in \cat{A}$ and  $s \in \cat{S}$ such that
$s \circ a = s \circ b$, there exists $t \in \cat{S}$ such that $a \circ t = b 
\circ t$.
\end{enumerate}
\end{dfn}

In the definition above the implicit assumption is that the sources and targets 
of the morphisms match; e.g., in (RD1) the morphisms $a$ and $s$ have the same 
target. Here are the diagrams illustrating the definition: 
\begin{equation} \label{eqn:2310}
\UseTips  \xymatrix @C=5ex @R=4ex  {
& 
K
\ar@{-->}[dl]_{t}
\ar@{-->}[dr]^{b}
\\
M 
\ar[dr]_{a}
& &
N
\ar[dl]^{s}
\\
&
L 
} 
\qquad 
\UseTips  \xymatrix @C=7ex @R=6ex {
K
\ar@{-->}[r]^{t}
&
M
\ar@(dl,ul)[d]_{a}
\ar@(dr,ur)[d]^{b}
\\
&
N
\ar[r]_{s}
&
L
} 
\end{equation}

Now the left version of Definition \ref{dfn:21}:

\begin{dfn} \label{dfn:2311}
Let $\cat{S}$ be multiplicatively closed set of morphisms in a  category 
$\cat{A}$. We say that $\cat{S}$ is a {\em left denominator set}
\index{Denominator set}
if it satisfies these two conditions:
\begin{enumerate}[leftmargin=30pt] 
\item[(LD1)] (Left Ore condition) Given $a \in \cat{A}$ and $s \in \cat{S}$,
there exist $b \in \cat{A}$ and $t \in \cat{S}$ such that
$t \circ a = b \circ s$.

\item[(LD2)] (Left cancellation condition) 
Given $a, b \in \cat{A}$ and  $s \in \cat{S}$ such that
$a \circ s = b \circ s$, there exists $t \in \cat{S}$ such that $t \circ a = t 
\circ b$.
\end{enumerate}
\end{dfn}

\begin{prop} \label{prop:2310}
Let $\cat{S}$ be multiplicatively closed set of morphisms in a  category 
$\cat{A}$. The two conditions below are equivalent\tup{:}
\begin{itemize}
\rmitem{i} $\cat{S}$ is a left denominator set in $\cat{A}$. 

\rmitem{ii} $\cat{S}^{\mrm{op}}$ is a right denominator set in 
$\cat{A}^{\mrm{op}}$. 
\end{itemize}
\end{prop}

\begin{exer} \label{exer:2312}
Prove Proposition \ref{prop:2310}. 
\end{exer}


Here is the main theorem regarding Ore localization. 

\begin{thm} \label{thm:103}
The following conditions are equivalent for a category $\cat{A}$ 
and a multiplicatively closed set of morphisms $\cat{S} \subseteq \cat{A}$. 
\begin{enumerate}
\rmitem{i} The right Ore localization
\index{Localization! Ore {\indash} of a category}
$(\cat{A}_{\cat{S}}, \opn{Q})$ exists.

\rmitem{ii} $\cat{S}$ is a right denominator set.
\index{Denominator set}
\end{enumerate}
\end{thm}

The proof of Theorem \ref{thm:103} is after some preparation. 
The hard part is proving that (ii) $\Rightarrow$ (i). 
The general idea is the same as in 
commutative localization: we consider the set of pairs of morphisms
$\cat{A} \times \cat{S}$, and define a relation $\sim$ on it, with the hope 
that this is an equivalence relation, and that the quotient set 
$\cat{A}_{\cat{S}}$ will be a category, and it will have the desired 
properties. 

Let's assume that $\cat{S}$ is a right denominator set. 
For any $M, N \in \opn{Ob}(A)$ consider the set 
\[ (\cat{A} \times \cat{S})(M, N) :=  
\coprod_{L \in \opn{Ob}(A)} 
\cat{A}(L, N) \times \cat{S}(L, M) . \]

\begin{rem} \label{rem:1405}
The set $(\cat{A} \times \cat{S})(M, N)$ could be big, namely not an element 
of the initial universe $\cat{U}$. This would require the introduction of a 
larger universe, say $\cat{V}$, in which $\cat{U}$ is an element. 
And the resulting category $\cat{A}_{\cat{S}}$ will be a 
$\cat{V}$-category. This is the same issue we had in the proof of Theorem 
\ref{thm:2310}. 

We will ignore this issue. Moreover, in many cases of interest (derived 
categories where there are DG enhancements, such as the K-injective 
enhancement), there will be an alternative presentation of $\cat{A}_{\cat{S}}$
as a $\cat{U}$-category. We will refer to this when we get to it. 
\end{rem}

An element 
$(a, s) \in (\cat{A} \times \cat{S})(M, N)$ can be pictured as a diagram
\begin{equation} \label{eqn:1408}
\UseTips  \xymatrix @C=5ex @R=4ex  {
& L 
\ar[dl]_{s}
\ar[dr]^{a}
\\
M 
& &
N
}
\end{equation}
in $\cat{A}$.
This diagram will eventually represent the right fraction 
$\opn{Q}(a) \circ \opn{Q}(s)^{-1} : M \to N$.

\begin{dfn}  \label{dfn:1408}
We define a relation $\sim$ on the set $\cat{A} \times \cat{S}$
like this:
$(a_1, s_1) \sim (a_2, s_2)$ 
if there exist $b_1, b_2 \in \cat{A}$ s.t.\ 
$a_1 \circ b_1 = a_2 \circ b_2$ 
and
$s_1 \circ b_1 = s_2 \circ b_2 \in \cat{S}$.
\end{dfn}

Note that the relation $\sim$ imposes condition (ii) of Lemma \ref{lem:3}.
Here it is in a commutative diagram, in which we have made the objects explicit:
\begin{equation} \label{eqn:1409}
\UseTips  \xymatrix @C=6ex @R=6ex  {
&
K
\ar@{-->}[dl]_{b_1}
\ar@{-->}[dr]^{b_2}
\\
L_1
\ar[d]_{s_1}
\ar[drr]_(0.75){a_1}
& &
L_2
\ar[dll]_(0.75){s_2}
\ar[d]^{a_2}
\\
M 
& &
N
} 
\end{equation}
The paths (i.e.\ morphisms) ending at $M$ are in $\cat{S}$. 

\begin{lem} \label{lem:4}
If $\cat{S}$ is a right denominator set, then the relation $\sim$ is an 
equivalence.
\end{lem}

\begin{proof}
Reflexivity: take $K := L$ and $b_i := \opn{id}_L : L \to L$. 
Symmetry is trivial. 

Now to prove transitivity. 
Suppose we are given 
$(a_1, s_1) \sim (a_2, s_2)$ and $(a_2, s_2) \sim (a_3, s_3)$.
So we have the solid part of the first diagram in (\ref{eqn:4925}),
and it is commutative. The morphisms ending at $M$ are in $\cat{S}$.
%\begin{figure}[!htb]
\begin{equation} \label{eqn:4925}
\UseTips  \xymatrix @C=6.5ex @R=6ex  {
H
\ar@{-->}[r]^{u}
&
I
\ar@{-->}[dl]_{t}
\ar@{-->}[dr]^{d}
\\
K
\ar[d]_{b_1}
\ar[dr]^{b_2}
& & 
J
\ar[dl]_{c_2}
\ar[d]^{c_3}
\\
L_1
\ar[d]_{s_1}
\ar[drr]_(0.65){a_1}
&
L_2
\ar[dl]_(0.6){s_2}
\ar[dr]^(0.6){a_2}
&
L_3
\ar[dll]^(0.65){s_3}
\ar[d]^{a_3}
\\
*+++{M} 
& &
*+++{N}
}
\qquad 
\UseTips  \xymatrix @C=5ex @R=6ex  {
H
\ar[dd]_{b_1 \circ t \circ u}
\ar[ddrr]^{c_3 \circ d \circ u}
\\
\rule{0pt}{0.5em}
\\
L_1
\ar[d]_{s_1}
\ar[drr]_(0.7){a_1}
&
&
L_3
\ar[dll]^(0.7){s_3}
\ar[d]^{a_3}
\\
*+{M}
& &
*+{N}
} 
\end{equation}
%\caption{Diagrams for the proof of Lemma \ref{lem:4}} 
%\label{fig:5}
%\end{figure}
 
By condition (RD1) applied to the morphisms $K \to M \leftarrow J$ there are 
$t \in \cat{S}$ and $d \in \cat{A}$ s.t.\ 
$(s_3 \circ c_3) \circ d = (s_1 \circ b_1) \circ t$.
Comparing the morphisms $I \to M$ in this diagram, we see that  
\[ s_2 \circ (b_2 \circ t) = s_1 \circ b_1 \circ t = s_3 \circ c_3 \circ d =
s_2 \circ (c_2 \circ d) . \]
By (RD2) there is $u \in \cat{S}$ s.t.\ 
$(b_2 \circ t) \circ u = (c_2 \circ d) \circ u$.
So all morphisms $H \to M$ are equal and belong to $\cat{S}$, and  all 
morphisms $H \to N$ are equal.
Now delete the object $L_2$ and the arrows going through it. 
Then delete the objects $I, J, K$,  but keep the paths going through them. 
We get the second diagram in (\ref{eqn:4925}). It is commutative, and all
morphisms ending at $M$ are in $\cat{S}$. This is evidence for 
$(a_1, s_1) \sim (a_3, s_3)$.
\end{proof}

\begin{proof}[Proof of Theorem \tup{\ref{thm:103}}] \mbox{}

\smallskip \noindent
Step 1. In this step we prove (i) $\Rightarrow$ (ii). 
Take $a \in \cat{A}$ and $s \in \cat{S}$ with the same target. 
Consider $q := \opn{Q}(s)^{-1} \circ \opn{Q}(a)$. By (RO3) there are 
$b \in \cat{A}$ and $t \in \cat{S}$ s.t.\ 
$q = \opn{Q}(b) \circ \opn{Q}(t)^{-1}$. 
So
$\opn{Q}(s \circ b) = \opn{Q}(a \circ t)$.
By (RO4) there is $u \in \cat{S}$ s.t.\ 
$(s \circ b) \circ u = (a \circ t) \circ u$.
We read this as 
$s \circ (b \circ u) = a \circ (t \circ u)$,
and note that $t \circ u \in \cat{S}$. So (RD1) holds.

Next $a, b \in \cat{A}$ and $s \in \cat{S}$ s.t.\ 
$s \circ a = s \circ b$.  Then 
$\opn{Q}(s \circ a) = \opn{Q}(s \circ b)$. But $\opn{Q}(s)$ is invertible, so 
$\opn{Q}(a) = \opn{Q}(b)$. By (RO4) there is $t \in \cat{S}$ s.t.\ 
$a \circ t = b \circ t$. We have proved (RD2). 


\medskip \noindent
Step 2. Now we assume that condition (ii) holds, and we define 
the morphism sets $\cat{A}_{\cat{S}}(M, N)$, composition between them, and the 
identity morphisms. 

For any $M, N \in \opn{Ob}(\cat{A})$ let
\[ \cat{A}_{\cat{S}}(M, N) :=
(\cat{A} \times \cat{S})(M, N) \, / \, \sim , \]
the quotient set modulo the relation $\sim$ from 
Definition \ref{dfn:1408}, which is an
equivalence relation by Lemma \ref{lem:4}.

We define composition like this. Given 
$q_1 \in \cat{A}_{\cat{S}}(M_0, M_1)$
and 
$q_2 \in \cat{A}_{\cat{S}}(M_1, M_2)$,
choose representatives 
$(a_i, s_i) \in (\cat{A} \times \cat{S})(M_{i-1}, M_i)$.
We use the notation 
$q_i = \ol{ (a_i, s_i) }$ to indicate this. 
By (RD1) there are $c \in \cat{A}$ and $u \in \cat{S}$ s.t.\ 
$s_2 \circ c = a_1 \circ u$. 
The composition 
$q_2 \circ q_1 \in \cat{A}_{\cat{S}}(M_0, M_2)$
is defined to be
\[ q_2 \circ q_1 := \ol{ (a_2 \circ c, s_1 \circ u) } 
\in (\cat{A} \times \cat{S})(M_{0}, M_2)) . \]
The idea behind the formula can be seen in diagram (\ref{eqn:4926}). 
%\begin{figure}[!htb]
\begin{equation} \label{eqn:4926}
\UseTips  \xymatrix @C=3.5ex @R=3.5ex  {
& & 
K
\ar[dl]_{u}
\ar[dr]^{c}
\\
&
L_1
\ar[dl]_{s_1}
\ar[dr]_(0.5){a_1}
& & 
L_2
\ar[dl]_(0.5){s_2}
\ar[dr]^{a_2}
\\
M_0 
& &
M_1
& &
M_2
}
\end{equation}
% \caption{Diagram for the proof of of Theorem \tup{\ref{thm:103}}}
%\label{fig:2}
%\end{figure}

We have to verify that this definition is independent of the representatives.
So suppose we take other representatives 
$q_i = \ol{ (a'_i, s'_i) }$, and we choose morphisms $u', c'$ to
construct the composition. This is the solid part of the diagram
(\ref{eqn:4927}), and it is a commutative diagram. We must prove that 
\[ \ol{ (a_2 \circ c, s_1 \circ u) }  = \ol{ (a'_2 \circ c', s'_1 \circ u') } . 
\]
%\begin{figure}[!htb]
\begin{equation} \label{eqn:4927}
\UseTips  \xymatrix @C=8ex @R=6ex  {
J_1
\ar@{-->}[d]_{b_1}
\ar@{-->}[dr]^(0.7){b'_1}
&
K
\ar[dl]^(0.7){u}
\ar[drr]^(0.7){c}
& & 
K'
\ar[dll]_(0.7){u'}
\ar[dr]_(0.7){c'}
& 
J_2
\ar@{-->}[dl]_(0.7){b_2}
\ar@{-->}[d]^{b'_2}
\\
L_1
\ar[d]_{s_1}
\ar[drr]_(0.65){a_1}
&
L'_1
\ar[dl]^(0.65){s'_1}
\ar[dr]^{a'_1}
& & 
L_2
\ar[dl]_{s_2}
\ar[dr]_(0.65){a_2}
& 
L'_2
\ar[dll]^(0.65){s'_2}
\ar[d]^{a'_2}
\\
M_0 
& &
*++{M_1}
& &
M_2
} 
\end{equation}
% \caption{{Diagram for the proof of of Theorem \tup{\ref{thm:103}}}} 
% \label{fig:3}
% \end{figure}

There are morphisms $b_i, b'_i$ that are evidence for 
$(a_i, s_i) \sim (a'_i, s'_i)$.
They are depicted as the dashed arrows in (\ref{eqn:4927}). 
That whole diagram is commutative. 
The morphisms $J_1 \to M_0$, $K \to M_0$, $K' \to M_0$ and 
$J_2 \to M_1$ are all in $\cat{S}$.

Choose $v_1 \in \cat{S}$ and $d_1 \in \cat{A}$ s.t.\ the first diagram in 
(\ref{eqn:4928}) is commutative. This can be done by (RD1). 
\begin{equation} \label{eqn:4928}
\UseTips \xymatrix @C=3.5ex @R=3.5ex {
\\
&
I_1
\ar[dl]_{v_1}
\ar[dr]^{d_1}
\\
J_1
\ar[dr]_{b_1}
& &
K
\ar[dl]^{u}
\\
&
L_1
} 
\quad  
\UseTips  \xymatrix @C=3.5ex @R=3.5ex {
I_2
\ar@{-->}[dr]^{\til{v}}
\\
& 
\til{I}_2
\ar@{-->}[dl]_{\til{v}_2}
\ar@{-->}[dr]^{\til{d}_2}
\\
K'
\ar[dr]_{c'}
& &
J_2
\ar[dl]^{b'_2}
\\
&
L'_2 
\ar[dr]_{s'_2}
\\
&
&
M_1
} 
\quad 
\xymatrix @C=3.5ex @R=3.5ex {
\\
&
I_2
\ar[dl]_{v_2}
\ar[dr]^{d_2}
\\
K'
\ar[dr]_{c'}
& &
J_2
\ar[dl]^{b'_2}
\\
&
L'_2
} 
\end{equation}

Consider the solid part of the middle diagram in (\ref{eqn:4928}).
Since $J_2 \to M_1$ is in $\cat{S}$, by (RD1) there are 
$\til{v}_2 \in \cat{S}$ and $\til{d}_2 \in \cat{A}$ s.t.\ the two paths
$\til{I}_2 \to M_1$ are equal.
By (RD2) there is $\til{v} \in \cat{S}$ s.t.\ the two paths 
$I_2 \to L'_2$ are equal. 
We get the commutative diagram in the middle of (\ref{eqn:4928}).
Next, defining 
$d_2 := \til{d}_2 \circ \til{v}$ and 
$v_2 := \til{v}_2 \circ \til{v} \in \cat{S}$,
we obtain the third commutative diagram in (\ref{eqn:4928}).

We now embed the first and third diagrams from (\ref{eqn:4928}) into the 
diagram (\ref{eqn:4927}). This gives us the solid diagram in (\ref{eqn:4929}),
and it is commutative. 
The morphisms $I_1 \to M_0$ belong to $\cat{S}$.  

% \begin{figure}[!htb]
\begin{equation} \label{eqn:4929}
\UseTips  \xymatrix @C=8ex @R=6ex  {
H''
\ar@{-->}[r]^{w''}
& 
H'
\ar@{-->}[r]^{w'}
& 
H
\ar@{-->}[dl]_{w}
\ar@{-->}[dr]^{e}
\\
&
I_1
\ar@{->}[dl]_{v_1}
\ar@{->}[d]^{d_1}
& & 
I_2
\ar@{->}[d]_{v_2}
\ar@{->}[dr]^{d_2}
\\
J_1
\ar[d]_{b_1}
\ar[dr]^(0.7){b'_1}
&
K
\ar[dl]^(0.7){u}
\ar[drr]^(0.7){c}
& & 
K'
\ar[dll]_(0.7){u'}
\ar[dr]_(0.7){c'}
& 
J_2
\ar[dl]_(0.7){b_2}
\ar[d]^{b'_2}
\\
L_1
\ar[d]_{s_1}
\ar[drr]_(0.65){a_1}
&
L'_1
\ar[dl]^(0.65){s'_1}
\ar[dr]^{a'_1}
& & 
L_2
\ar[dl]_{s_2}
\ar[dr]_(0.65){a_2}
& 
L'_2
\ar[dll]^(0.65){s'_2}
\ar[d]^{a'_2}
\\
M_0 
& &
*++{M_1}
& &
M_2
}
\end{equation}
% \caption{{Diagram for the proof of of Theorem \tup{\ref{thm:103}}}} 
% \label{fig:8}
% \end{figure}

Choose $w \in \cat{S}$ and $e \in \cat{A}$, starting at an object $H$, to fill 
the diagram $I_1 \to M_0 \leftarrow I_2$, using (RD1). The path 
$H \to I_1 \to M_0$ is in $\cat{S}$, and all the paths $H \to M_0$ are equal.   
But we could have failure of commutativity in the paths $H \to L'_1$ and 
$H \to L_2$. 

The two paths  $H \to L'_1$ in (\ref{eqn:4929}) satisfy 
\[ s'_1 \circ (b'_1 \circ v_1 \circ w) = s'_1 \circ (u' \circ v_2 \circ e) . \]
Therefore, by (RD2), there is $w' \in \cat{S}$ s.t.\ 
\[  (b'_1 \circ v_1 \circ w) \circ w' = (u' \circ v_2 \circ e) \circ w' . \]
Next, the two paths  $H' \to L_2$ satisfy 
\[ s_2 \circ (c \circ d_1 \circ w \circ w') = s_2 \circ (b_2 \circ d_2 \circ e 
\circ w') ; \]
this is because we can take a detour through $L'_1$. 
Therefore, again by (RD2), there is $w'' \in \cat{S}$ s.t.\ 
\[ (c \circ d_1 \circ w \circ w') \circ w'' = (b_2 \circ d_2 \circ e \circ w') 
\circ w'' . \]
Now all paths $H'' \to M_2$ in (\ref{eqn:4929}) are equal. All paths 
$H'' \to M_0$ are equal and are in $\cat{S}$. 

Erase the objects $M_1, J_1, J_2$ and all arrows touching them
from (\ref{eqn:4929}). 
Then erase $H, H'$, but keep the paths through them. 
We obtain the commutative diagram (\ref{eqn:4930}).
%\begin{figure}[!htb]
%\cmnt{ (190123)  }
\begin{equation} \label{eqn:4930}
\UseTips  \xymatrix @C=8.5ex @R=6ex  {
&
I_1
\ar[d]^{d_1}
& 
H''
\ar[l]_{w \circ w' \circ w''}
\ar[r]^{e \circ w' \circ w''}
& 
I_2
\ar[d]_{v_2}
\\
&
K
\ar[dl]^(0.7){u}
\ar[drr]^(0.7){c}
& & 
K'
\ar[dll]_(0.7){u'}
\ar[dr]_(0.7){c'}
& 
\\
L_1
\ar[d]_{s_1}
&
L'_1
\ar[dl]^(0.65){s'_1}
& & 
L_2
\ar[dr]_(0.65){a_2}
& 
L'_2
\ar[d]^{a'_2}
\\
M_0 
& &
& &
M_2
}
\end{equation}
% \caption{{Diagram for the proof of of Theorem \tup{\ref{thm:103}}}} 
% \label{fig:4}
% \end{figure}
This is evidence for 
\[  (a_2 \circ c, s_1 \circ u) \sim (a'_2 \circ c', s'_1 \circ u') . \]
The proof that composition is well-defined is done.

The identity morphism $\opn{id}_M$ in $\cat{A}_{\cat{S}}$
of an object $M$ is $\ol{(\opn{id}_M, \opn{id}_M)}$.  

\medskip \noindent
Step 3. 
We have to verify the associativity and the identity properties of 
composition in $\cat{A}_{\cat{S}}$. Namely that $\cat{A}_{\cat{S}}$ is a
category. This seems to be not too hard, given Step
2, and we leave it as an exercise.

\medskip \noindent
Step 4.  
The functor $\opn{Q} : \cat{A} \to \cat{A}_{\cat{S}}$ is defined to be 
$\opn{Q}(M) := M$ on objects, and $\opn{Q}(a) := \ol{(a, \opn{id}_M)}$ for 
$a : M \to N$ in $\cat{A}$. We have to verify this is a functor... Again, an 
exercise.

\medskip \noindent
Step 5.  
Finally we verify properties (RO1)-(RO4). (RO1) is clear. The inverse of
$\opn{Q}(s)$ is $\ol{(\opn{id}, s)}$, so (RO2) holds. 

It is not hard to see that 
$\ol{(a, s)} = \ol{(a, \opn{id})} \circ \ol{(\opn{id}, s)}$;
this is (RO3). 

If $\opn{Q}(a_1) = \opn{Q}(a_2)$, then 
$(a_1, \opn{id}_M) \sim (a_2, \opn{id}_M)$; so there
are $b_1, b_2 \in \cat{A}$  s.t.\ 
$a_1 \circ b_1 = a_2 \circ b_2$ and 
$\opn{id} \circ \, b_1 = \opn{id} \circ \, b_2 \in \cat{S}$. 
Writing $s := b_1 \in \cat{S}$, we get $a_1 \circ s = a_2 \circ s$. This proves 
(RO4). 
\end{proof}

\begin{exer} \label{exer:3030}
Finish the details in the proof of Theorem \tup{\ref{thm:103}}. 
\end{exer}

\begin{prop}[Common Denominator] \label{prop:107}
Let $\cat{A}$ be a category, let $\cat{S}$ be a right
denominator set in $\cat{A}$, and let $(\cat{A}_{\cat{S}}, \opn{Q})$ be the 
right Ore
localization.
For every two morphisms $q_1, q_2 : M \to N$ in $\cat{A}_{\cat{S}}$ there
is a common right denominator. Namely we can write 
$q_i = \opn{Q}(a_i) \circ \opn{Q}(s)^{-1}$
for suitable $a_i \in \cat{A}$ and $s \in \cat{S}$. 
\end{prop}

\begin{proof}
Choose representatives
$q_i = \opn{Q}(a'_i) \circ \opn{Q}(s'_i)^{-1}$. 
By (RD1) applied to $L_1 \to M \leftarrow L_2$, 
there are $b \in \cat{A}$ and $t \in \cat{S}$ 
such that the part of diagram (\ref{eqn:4931}) that lies above $M$ is 
commutative:
\begin{equation} \label{eqn:4931}
\UseTips  \xymatrix @C=8ex @R=6ex  {
L_1
\ar[d]_{s'_1}
\ar[drr]_(0.75){a'_1}
& 
L
\ar[l]_{t}
\ar[r]^{b}
&
L_2
\ar[dll]^(0.75){s'_2}
\ar[d]^{a'_2}
\\
M 
& 
&
N
} 
\end{equation}
Write
$s := s'_1 \circ t = s'_2 \circ b$,  $a_1 := a'_1 \circ t$ and 
$a_2 := a'_2 \circ b$. 
By Lemma \ref{lem:3} we get 
$q_i = \opn{Q}(a_i) \circ \opn{Q}(s)^{-1}$.
\end{proof}

\begin{prop} \label{prop:3380} 
Let $\cat{A}$ be a category, let $\cat{S} \sub \cat{A}$ be a right denominator 
set, and let $\cat{A}' \sub \cat{A}$ be a full subcategory.  
Define $\cat{S}' := \cat{A}' \cap \, \cat{S}$. 
Assume these two conditions hold\tup{:}
\begin{itemize}
\rmitem{i} $\cat{S}'$ is a right denominator set in $\cat{A}'$.

\rmitem{ii} Let $M \in \opn{Ob}(\cat{A})$. If 
there exists a morphism $s : M \to L'$ in $\cat{S}$ with 
$L' \in \opn{Ob}(\cat{A}')$, 
there exists a morphism $t : K' \to M$ in $\cat{S}$ with 
$K' \in \opn{Ob}(\cat{A}')$.
\end{itemize}
Then the canonical functor 
$\cat{A}'_{\cat{S}'} \to \cat{A}_{\cat{S}}$ 
is fully faithful.
\end{prop}

\begin{proof}
Let's denote the inclusion functor by $F : \cat{A}' \to \cat{A}$. 
We want to prove that its localization 
$F_{\cat{S}'} : \cat{A}'_{\cat{S}'} \to \cat{A}_{\cat{S}}$ 
is fully faithful.

\medskip \noindent
Step 1. Let $L'_1, L'_2 \in \opn{Ob}(\cat{A}')$, and let 
$q : L'_1 \to L'_2$ be a
morphism in $\cat{A}_{\cat{S}}$. Choose a presentation 
$q = \opn{Q}(a) \circ \opn{Q}(s)^{-1}$ with
$s : M \to L'_1$ a morphism in $\cat{S}$
and $a : M \to L'_2$  a morphism in $\cat{A}$. 
This is possible because $\cat{S}$ is a right denominator set in $\cat{A}$. 
By condition (ii) we can find a morphism $t : K' \to M$ in $\cat{S}$ 
with $K' \in \opn{Ob}(\cat{A}')$. See next diagram. 
\[ \UseTips  \xymatrix @C=6ex @R=6ex  {
K'
\ar[r]^{t}
& 
M 
\ar[dl]_{s}
\ar[dr]^{a}
\\
L'_1
\ar@{-->}[rr]^{q}
& &
L'_2
} \]
Then  
$q = \opn{Q}(a \circ t) \circ \opn{Q}(s \circ t)^{-1}$.
But $s \circ t \in \cat{S}'$ and $a \circ t  \in \cat{A}'$, 
so $q$ is in the image of the functor $F_{\cat{S}'}$. We see that 
$F_{\cat{S}'}$ is full. 

\medskip \noindent
Step 2. Let $p', q' : L'_1 \to L'_2$ 
be morphisms in $\cat{A}'_{\cat{S}'}$
such that $F_{\cat{S}'}(p') = F_{\cat{S}'}(q')$.  
Let us denote the localization functor of $\cat{A}'$ by
$\opn{Q}' : \cat{A}' \to \cat{A}'_{\cat{S}'}$. 
Because $\cat{S}'$ is a right denominator set in $\cat{A}'$,
and using Proposition \ref{prop:107}, we can find presentations 
$p' = \opn{Q}'(a') \circ \opn{Q}'(s')^{-1}$
and
$q' = \opn{Q}'(b') \circ \opn{Q}'(s')^{-1}$
with morphisms 
$s' : N' \to L'_1$ in $\cat{S}'$
and 
$a', b' : N' \to L'_2$ in $\cat{A}'$.
Next, since $F_{\cat{S}'}(p') = F_{\cat{S}'}(q')$, Lemma \ref{lem:3} tells us 
that there are morphisms 
$u, v : M \to N'$ in $\cat{A}$ s.t.\ 
$a' \circ u = b' \circ v$, and 
$s' \circ u = s' \circ v \in \cat{S}$. 
Condition (ii), applied to the morphism 
$s' \circ u : M \to L'_1$, says that there is a morphism 
$t : K' \to M$ in $\cat{S}$ with source $K' \in \opn{Ob}(\cat{A}')$.
See diagram below. 
%\begin{figure}[!htb]
\[ \UseTips  \xymatrix @C=6ex @R=6ex  {
K'
\ar[r]^{t}
&
M 
\ar[r]^{u, v}
&
N'
\ar[dl]_{s'}
\ar[dr]^{a', b'}
\\
&
L'_1
\ar@{-->}[rr]^{p', q'}
& &
L'_2
} \]
% \caption{Diagram for the proof of of Proposition \tup{\ref{prop:3380}}}
% \label{fig:4690}
% \end{figure}
Then we have 
\[ \begin{aligned}
&
p' = \opn{Q}'(a' \circ u \circ t) \circ \opn{Q}'(s' \circ u \circ t)^{-1} 
\\ & \quad 
= \opn{Q}'(b' \circ u \circ t) \circ \opn{Q}'(s' \circ u \circ t)^{-1} = q' .
\end{aligned} \] 
This proves that $F_{\cat{S}'}$ is faithful.
\end{proof}

In the next subsection we study Ore localization of linear categories, and give 
rings as examples (see Examples \ref{exa:4207} and \ref{exa:1405}). And in 
Section \ref{sec:der-cat} we study localization of triangulated categories. 
Here is a nonlinear example. 

\begin{exa} \label{exa:4205}
Fix a base commutative ring $\K$ and a topological space (or, more generally, a 
site) $X$. Let $\K_X$ be the constant sheaf of rings $\K$ on $X$.
Consider the category
$\catt{DGRng}^{\leq 0}_{\mrm{sc}} / \K_X$
of sheaves of commutative DG $\K_X$-rings on $X$ (see Definition 
\ref{dfn:3090} and Remark \ref{rem:4945}). For the sake of brevity let us 
just write 
$\catt{DGRng}_{X} := \catt{DGRng}^{\leq 0}_{\mrm{sc}} / \K_X$.
 
The localization of $\catt{DGRng}_X$ with respect to the quasi-isomorphisms in 
it is called the {\em derived category of commutative DG $\K_X$-rings}, and the 
notation is $\dcat{D}(\catt{DGRng}_X)$. 
The categorical localization functor is 
$\opn{Q} : \catt{DGRng}_X \to \dcat{D}(\catt{DGRng}_X)$,
and it is the identity on objects. 
(The unfortunate tradition in homotopy theory is to call such a localization 
the ``homotopy category'', and to denote it by 
``$\opn{Ho}(\catt{DGRng}_X)$''.)

As explained in the lecture notes \cite{Ye15}, every DG $\K_X$-ring $\AA$ 
admits a {\em semi-pseudo-free resolution} $\til{\AA} \to \AA$ in 
$\catt{DGRng}_X$. The semi-pseudo-free DG $\K_X$-rings have good lifting 
properties. Indeed, they behave like cofibrant objects in a Quillen model 
structure (even though there does not seems to be such a structure on 
$\catt{DGRng}_X$, except when $X$ is a discrete space).

There is a congruence on the category $\catt{DGRng}_X$, called the {\em 
quasi-homotopy relation}. (It turns out that this idea was already known to some 
homotopy theorists, see Remark \ref{rem:4205} below.) Let us denote by 
$\dcat{K}(\catt{DGRng}_X)$ the {\em homotopy category}, which is the 
quotient of $\catt{DGRng}_X$ modulo this congruence. So there is a functor 
$\opn{P} : \catt{DGRng}_X \to \dcat{K}(\catt{DGRng}_X)$,
which is the identity on objects, and it sends a DG ring homomorphism to its 
quasi-homotopy class. 

The functors $\opn{Q}$ and $\opn{P}$ fit into this commutative diagram:
\begin{equation} \label{eqn:4913}
\UseTips \xymatrix @C=8ex @R=6ex {
\catt{DGRng}_X
\ar[r]^{\opn{P}}
\ar@(u,u)[rr]^{\opn{Q}}
&
\dcat{K}(\catt{DGRng}_X)
\ar[r]^{\bar{\opn{Q}}}
&
\dcat{D}(\catt{DGRng}_X)
}
\end{equation}
The functor $\bar{\opn{Q}}$ is a {\em faithful right Ore localization}. 
The proof relies on the lifting properties of semi-pseudo-free DG $\K_X$-rings. 
See Remark \ref{rem:4945} for related material. 
\end{exa}

\begin{rem} \label{rem:4205}
Suppose $\cat{C}$ is a category, with a given multiplicatively closed set of 
morphisms $\cat{S} \sub \cat{C}$. If $\cat{C}$ admits a {\em Quillen model 
structure} for which $\cat{S}$ is the set of {\em weak equivalences}, then 
there is automatically a {\em calculus of fractions} for the localization
functor $\opn{Q} : \cat{C} \to \cat{C}_{\cat{S}}$.

For instance, suppose that all objects of $\cat{C}$ are {\em fibrant}. 
In this case $\cat{C}$ is called a {\em category of fibrant objects}. 
Then there is a congruence relation $\cat{H}$ on $\cat{C}$,
and a commutative diagram 
\[ \UseTips \xymatrix @C=8ex @R=6ex {
\catt{C}
\ar[r]^{\opn{P}}
\ar@(u,u)[rr]^{\opn{Q}}
&
\cat{C} / \cat{H}
\ar[r]^{\bar{\opn{Q}}}
&
\cat{C}_{\cat{S}}
} \]
in which, like in (\ref{eqn:4913}), the functor $\opn{P}$ is full, and the 
functor $\bar{\opn{Q}}$ is a right Ore localization. 
These facts are explained in \cite{nLab} (search for the emphasized text in 
this remark). 
\end{rem}

\begin{rem} \label{rem:4970}
This is a continuation of Remark \ref{rem:4960} on strictly unital 
$\mrm{A}_{\infty}$ categories%
\index{Ainfty@$\mrm{A}_{\infty}$ category}, 
and it is related to Example \ref{exa:4205} and Remark \ref{rem:4205}
above, and to Remark \ref{rem:4605} below. The references are \cite{Kel5} and 
\cite{COS}. Recall that the base ring $\K$ is a field. 

Let $\cat{A}$ and $\cat{B}$ be $\mrm{A}_{\infty}$ categories,
i.e.\ objects of the category $\cat{A_{\infty}Cat}$. 
An $\mrm{A}_{\infty}$ functor $F : \cat{A} \to \cat{B}$
is called a {\em quasi-equivalence} if 
$\opn{H}(F) : \opn{H}(\cat{A}) \to \opn{H}(\cat{B})$
is an equivalence of graded categories.
The {\em derived category} of $\cat{A_{\infty}Cat}$
is the localization $\dcat{D}(\cat{A_{\infty}Cat})$ of $\cat{A_{\infty}Cat}$ 
with respect to the quasi-equivalences. 
(As mentioned above, the tradition is to call 
$\dcat{D}(\cat{A_{\infty}Cat})$ the ``homotopy category'' of 
$\cat{A_{\infty}Cat}$, and to denote 
it by ``$\opn{Ho}(\cat{A_{\infty}Cat})$''.)
There is a localization functor 
$\opn{Q} : \cat{A_{\infty}Cat} \to \dcat{D}(\cat{A_{\infty}Cat})$, 
which is the identity on objects. 

Suppose $F, F' : \cat{A} \to \cat{B}$ are $\mrm{A}_{\infty}$ functors. So 
$F$ and $F'$ are objects of the $\mrm{A}_{\infty}$ category 
$\cat{A_{\infty}Fun}(\cat{A}, \cat{B})$.
We say that $F$ and $F'$ are {\em $\mrm{A}_{\infty}$ homotopic} if 
they are isomorphic in the category 
$\opn{H}^0 \bigl( \cat{A_{\infty}Fun}(\cat{A}, \cat{B}) \bigr)$, 
in which case we write $F \approx F'$. 
Define $\dcat{K}(\cat{A_{\infty}Cat})$
to be the quotient of $\cat{A_{\infty}Cat}$ modulo the relation $\approx$. 
There is a full functor 
$\opn{P} : \cat{A_{\infty}Cat} \to \dcat{K}(\cat{A_{\infty}Cat})$, 
which is the identity on objects. 
An $\mrm{A}_{\infty}$ functor $F : \cat{A} \to \cat{B}$
is called an {\em $\mrm{A}_{\infty}$ equivalence} if it is an isomorphism in 
$\dcat{K}(\cat{A_{\infty}Cat})$. This means that there is  
an $\mrm{A}_{\infty}$ functor $G : \cat{B} \to \cat{A}$, 
such that $G \circ F \approx \opn{Id}_{\cat{A}}$ 
and $F \circ G \approx \opn{Id}_{\cat{B}}$.
There is a commutative diagram 
\begin{equation} \label{eqn:4975}
\UseTips \xymatrix @C=8ex @R=6ex {
\cat{A_{\infty}Cat}
\ar[r]^{\opn{P}}
\ar@(u,u)[rr]^{\opn{Q}}
&
\dcat{K}(\cat{A_{\infty}Cat})
\ar[r]^{\bar{\opn{Q}}}
&
\dcat{D}(\cat{A_{\infty}Cat})
}
\end{equation}

An important theorem (see \cite[Theorem 9.2.0.4 ]{Lef}) says that 
an $\mrm{A}_{\infty}$ functor $F$ is a quasi-equivalence iff it is an
 $\mrm{A}_{\infty}$ equivalence. This implies that the functor 
$\bar{\opn{Q}}$ in diagram (\ref{eqn:4975}) is an isomorphism of categories. 

It is worthwhile to mention that the category $\cat{A_{\infty}Cat}$
does not have a Quillen model structure (for which the weak equivalences are 
the quasi-equivalences). 

Let $\cat{DGCat}$ be the category of $\K$-linear DG categories, where the 
morphisms are the DG functors (see Definition \ref{dfn:1065}). 
There is a fully faithful embedding 
$\cat{DGRng} \centover \K \to \cat{DGCat}$,
since the DG rings (see Definition \ref{dfn:1074}) are the 
single-object DG categories.
Because an $\mrm{A}_{\infty}$ category whose higher operations vanish is a DG 
category, there is a faithful (but not full) embedding
$\cat{DGCat} \to \cat{A_{\infty}Cat}$.
The {\em derived category} of $\cat{DGCat}$
is the localization $\dcat{D}(\cat{DGCat})$ of $\cat{DGCat}$ with 
respect to the quasi-equivalences in it. 
(Once more, the traditional notation for $\dcat{D}(\cat{DGCat})$
is ``$\opn{Ho}(\cat{DGCat})$''.)
The embedding of $\cat{DGCat}$ into $\cat{A_{\infty}Cat}$ induces a functor 
$\dcat{D}(\cat{DGCat}) \lb \to \dcat{D}(\cat{A_{\infty}Cat})$,
and this is an equivalence. Together with the 
equivalence $\bar{\opn{Q}}$ from (\ref{eqn:4975}), this implies that
every morphism $F : \cat{A} \to \cat{B}$ in 
$\dcat{D}(\cat{DGCat})$
can be represented by an $\mrm{A}_{\infty}$ functor 
$\til{F} : \cat{A} \to \cat{B}$, which is unique up to a unique isomorphism in 
$\dcat{K}(\cat{A_{\infty}Cat})$. 

Moreover, the category $\dcat{D}(\cat{DGCat})$ is 
symmetric monoidal, with tensor operation $\cat{A} \ot \cat{B}$. 
It turns out that the DG category 
$\cat{A_{\infty}Fun}(\cat{A}, \cat{B})$
is the {\em internal Hom object}, making  
$\dcat{D}(\cat{DGCat})$ into a {\em closed monoidal category}.
\end{rem}





\mysubsection{Localization of Linear Categories} \label{subsec:loc-lin}

Until now in this section we dealt with arbitrary categories. 
In this subsection our categories will be linear over the 
commutative base ring $\K$ (which will be implicit most of the time).
See Convention \ref{conv:2490}. 

For convenience we only talk about right denominator sets here. 
All the statements hold equally for left denominator sets; cf.\ Remark 
\ref{rem:1845} and Proposition \ref{prop:1845}.

\begin{thm} \label{thm:104}
\index{Localization! Ore {\indash} of a linear category}
Let $\cat{A}$ be a $\K$-linear category, let $\cat{S}$ be a right
denominator set in $\cat{A}$, and let $(\cat{A}_{\cat{S}}, \opn{Q})$ be the 
right Ore
localization.

\begin{enumerate}
\item  The category $\cat{A}_{\cat{S}}$ has a unique $\K$-linear structure,
such that $\opn{Q} : \cat{A} \to \cat{A}_{\cat{S}}$ is a $\K$-linear 
functor. 

\item Suppose $\cat{B}$ is another $\K$-linear category, and 
$F : \cat{A} \to \cat{B}$ is a $\K$-linear functor s.t.\
$F(s)$ is invertible for every $s \in \cat{S}$. Let 
$F_{\cat{S}} : \cat{A}_{\cat{S}} \to \cat{B}$ be the localization of $F$. Then 
$F_{\cat{S}}$ is a $\K$-linear functor. 

\item If $\cat{A}$ is an additive category, then so is $\cat{A}_{\cat{S}}$.
\end{enumerate}
\end{thm}

\begin{proof} \mbox{}

\smallskip \noindent
(1) Let $q_1, q_2 : M \to N$ be morphisms in $\cat{A}_{\cat{S}}$.
Choose common denominator presentations
$q_i = \opn{Q}(a_i) \circ \opn{Q}(s)^{-1}$. Since $\opn{Q}$ must be an additive 
functor, we have to define 
\begin{equation} \label{eqn:1430}
\opn{Q}(a_1) + \opn{Q}(a_2) := \opn{Q}(a_1 + a_2) .
\end{equation}
By the distributive law (bilinearity of composition) we must define 
$q_1 + q_2 := \opn{Q}(a_1 + a_2) \circ \opn{Q}(s)^{-1}$.
And for $\la \in \K$ we must define 
$\la \cdot q_i := \opn{Q}(\la \cd a_i) \circ \opn{Q}(s)^{-1}$.

The usual tricks are then used to prove independence of representatives. 
For instance, to prove that (\ref{eqn:1430}) is independent of choices, suppose 
that $\opn{Q}(a_1) = \opn{Q}(a'_1)$ and $\opn{Q}(a_2) = \opn{Q}(a'_2)$. 
Then, by (RO4), there are $t_1, t_2 \in \cat{S}$ such that 
$a_1 \circ t_1 = a'_1 \circ t_1$ and $a_2 \circ t_2 = a'_2 \circ t_2$. 
By (RD1) there exist $b \in \cat{A}$ and $v \in \cat{S}$ s.t.\ 
$t_1 \circ b = t_2 \circ v$. 
Let $t_3 := t_2 \circ v \in \cat{S}$.
Then 
$(a_1 + a_2) \circ t_3 = (a'_1 + a'_2) \circ t_3$,
and hence 
$\opn{Q}(a_1 + a_2) = \opn{Q}(a'_1 + a'_2)$. 

In this way $\cat{A}_{\cat{S}}$ is a $\K$-linear category, and $\opn{Q}$ is a 
$\K$-linear functor. 

\medskip \noindent (2)
The only option for $F_{\cat{S}}$ is 
$F_{\cat{S}}(q_i) := F(a_i) \circ F(s)^{-1}$. 
The usual tricks are used to prove independence of representatives.

\medskip \noindent 
(3) Clear from Propositions \ref{prop:2} and \ref{prop:109}. 
\end{proof}

\begin{exer} \label{exer:5025}
Finish the proofs of items (1) and (2) of the theorem. 
\end{exer}


\begin{rem} \label{rem:4202}
Let $A$ be a ring, which we can think of as a one-object linear category
$\cat{A}$. In this context, Theorem
\ref{thm:104} is one of the most important results in noncommutative ring 
theory. See the books \cite{McRo}, \cite{Row} or \cite{Ste}. 
\end{rem}

Here are somewhat concrete examples. 

\begin{exa} \label{exa:1405}
Suppose $A$ is a noncommutative ring (i.e.\ it is not necessarily commutative), 
and $S$ is a central  multiplicatively closed set in $A$ (i.e.\ 
$S \sub \opn{Cent}(A)$). Because the elements of $S$ commute with all elements 
of $A$, the denominator conditions hold automatically. 
We get a left and right Ore localization of rings 
$\opn{Q} : A \to A_S$.

Note that if $S$ contains a nilpotent element, then the ring $A_S$ is trivial. 
\end{exa}

The last observation should serve as a warning: localization can sometimes 
kill everything. Fortunately, the localization procedure that 
gives rise to the derived category does not cause any catastrophe, as we shall 
see in Proposition \ref{prop:1410}.

\begin{exa} \label{exa:4207}
Suppose $A$ is a noncommutative ring that is both noetherian and an integral 
domain (i.e.\ the only zero-divisor in $A$ is $0$). Let $S$ be the set of 
nonzero elements of $A$. Then $S$ is both a left and a right denominator set. 
See \cite[Theorem 2.1.15]{McRo}. 
\end{exa}
 
\begin{rem} \label{rem:1}
Suppose $A$ is a ring and $S$ is a right denominator set in it. 
Then the right Ore localization $A_S$ is {\em flat} as left $A$-module. 
See \cite[Theorem 3.1.20]{Row}. We have no idea if something like this is true 
for linear categories with more than one object. 
\end{rem}

\begin{prop} \label{prop:1840}
Let $(\cat{K}, \opn{T})$ be a T-additive category, let
$\cat{S}$ be a right denominator set in $\cat{K}$ such that 
$\opn{T}(\cat{S}) = \cat{S}$, and let 
$\opn{Q} : \cat{K} \to  \cat{K}_{\cat{S}}$ be the localization functor.
\begin{enumerate}
\item There is a unique additive automorphism
$\opn{T}_{\cat{S}}$ of the category $\cat{K}_{\cat{S}}$, such that 
$\opn{T}_{\cat{S}} \circ \opn{Q} = \opn{Q} \circ \opn{T}$
as functors $\cat{K} \to  \cat{K}_{\cat{S}}$.

\item Let $\tau$ be the identity automorphism of the functor 
$\opn{Q} \circ \opn{T}$. Then 
$(\opn{Q}, \tau) :  (\cat{K}, \opn{T}) \to 
(\cat{K}_{\cat{S}}, \opn{T}_{\cat{S}})$
is a T-additive functor. 
\end{enumerate}
\end{prop}

\begin{proof} \mbox{}

\smallskip \noindent
(1) By the assumption the functor 
$\opn{Q} \circ \opn{T} : \cat{K} \to  \cat{K}_{\cat{S}}$
sends the morphisms in $\cat{S}$ to isomorphisms. By the property (Loc3) of 
localization in Definition \ref{dfn:1406}, the functor 
$\opn{T}_{\cat{S}} : \cat{K}_{\cat{S}} \to  \cat{K}_{\cat{S}}$ 
satisfying
$\opn{T}_{\cat{S}} \circ \opn{Q} = \opn{Q} \circ \opn{T}$
exists and is unique. 
Similarly, there is a unique functor 
$\opn{T}^{-1}_{\cat{S}} : \cat{K}_{\cat{S}} \to  \cat{K}_{\cat{S}}$ 
satisfying
$\opn{T}^{-1}_{\cat{S}} \circ \opn{Q} = \opn{Q} \circ \opn{T}^{-1}$.
An easy calculation shows that 
$\opn{T}^{-1}_{\cat{S}} \circ \opn{T}_{\cat{S}} = \opn{Id} = 
\opn{T}_{\cat{S}} \circ \opn{T}^{-1}_{\cat{S}}$.
Hence $\opn{T}_{\cat{S}}$ is an automorphism of $\cat{K}_{\cat{S}}$.
By Theorem \ref{thm:104} the functor $\opn{T}_{\cat{S}}$ is additive. 

\medskip \noindent 
(2) This is clear. 
\end{proof}

The composition of T-additive functors was defined in Definition \ref{dfn:1840}.

\begin{prop} \label{prop:1841}
In the situation of Proposition \tup{\ref{prop:1840}}, suppose 
$(\cat{K}', \opn{T}')$ is another  T-additive category, and 
$(F, \nu) :  (\cat{K}, \opn{T}) \to (\cat{K}', \opn{T}')$
is a T-additive functor, such that $F(s)$ is invertible for every 
$s \in \cat{S}$. Let 
$F_{\cat{S}} : \cat{K}_{\cat{S}} \to \cat{K}'$
be the localized functor. 
Then there is a unique isomorphism 
$\nu_{\cat{S}} : F_{\cat{S}} \circ \opn{T}_{\cat{S}} \iso
\opn{T}' \circ \, F_{\cat{S}}$
of functors $\cat{K}_{\cat{S}} \to  \cat{K}'$, such that
$(F, \nu) = (F_{\cat{S}}, \nu_{\cat{S}}) \circ (\opn{Q}, \tau)$ 
as T-additive functors 
$(\cat{K}, \opn{T}) \to (\cat{K}', \opn{T}')$.
\end{prop}

\begin{exer}  \label{exer:1840}
Prove Proposition \ref{prop:1841}. 
\end{exer}




\cleardoublepage
\mysection{The Derived Category 
\texorpdfstring{$\dcat{D}(A, \cat{M})$}{D(A,M)}} 
\label{sec:der-cat}

\AYcopyright

In this section we introduce the main mathematical concept of the book: the 
{\em derived category} $\dcat{D}(A, \cat{M})$ of DG $A$-modules in $\cat{M}$. 
Here $A$ is a central DG $\K$-ring, and $\cat{M}$ is a $\K$-linear abelian 
category. The base ring $\K$ can be any nonzero commutative ring, and it will 
remain implicit most of the time. 


\mysubsection{Localization of Triangulated Categories}

Let $\cat{K}$ be a triangulated category, with translation functor $\opn{T}$. 
As we did in Section \ref{sec:loc-cats}, we shall often write 
$a \in \cat{K}$ for a morphism 
$a \in \cat{K}(M, N) = \opn{Hom}_{\cat{K}}(M, N)$, leaving the objects 
implicit. 

\begin{prop} \label{prop:106}
Suppose 
$H : \cat{K} \to \cat{M}$ is a cohomological functor, where $\cat{M}$ is some
abelian category. Let 
\[ \cat{S} := \bigl\{ s \in \cat{K} \mid H( \opn{T}^i(s)) \tup{ is invertible 
for all } i \in \Z \bigr\} . \] 
Then $\cat{S}$ is a left and right denominator set in $\cat{K}$.
\end{prop}

\begin{proof}
It is clear that $\cat{S}$ is closed under composition and contains the
identity morphisms. So it is a multiplicatively closed set. 

Let's prove that condition (RD1) of Definition \ref{dfn:21} holds. Suppose we 
are given morphisms
$L \xar{a} N \xleftarrow{s} M$ with $s \in \cat{S}$. 
We need to find morphisms 
$L \xleftarrow{t} K \xar{b} M$ with $t \in \cat{S}$ and such that 
$a \circ t = s \circ b$.  
 
Consider the solid commutative diagram 
\[ \UseTips  \xymatrix @C=8ex @R=6ex { 
K
\ar[r]^{t}
\ar@{-->}[d]_{b}
&
L
\ar[r]^{c \circ a}
\ar[d]_{a}
& 
P
\ar[r]^{}
\ar[d]_{=}
& 
\opn{T}(K) 
\ar@{-->}[d]_{\opn{T}(b)}
\\
M
\ar[r]^{s}
&
N
\ar[r]^{c}
& 
P
\ar[r]^{}
&
\opn{T}(M) 
}  \]
where the bottom row is a distinguished triangle built on $M \xar{s} N$,
and and the top row is a distinguished triangle built on 
$L \xar{c \circ a} P$, then turned $120^\circ$ to the right. By axiom (TR3) 
there 
is a morphism $b$ making the diagram commutative. Thus 
$a \circ t = s \circ b$.  Since $H(\opn{T}^i(s))$ are
invertible for all $i \in \Z$, it follows that $H(\opn{T}^i(P)) = 0$. But 
then $H(\opn{T}^i(t))$ are invertible for all $i \in \Z$, so 
$t \in \cat{S}$. 

Next we prove condition (RD2) of Definition \ref{dfn:21}. Because we are in 
an additive category, this condition is simplified: given $a \in \cat{K}$ and 
$s \in \cat{S}$ satisfying 
$s \circ a = 0$, we have to find $t \in \cat{S}$ satisfying 
$a \circ t = 0$. 

Say the objects involved are $L \xar{a} M \xar{s} N$. 
Take a distinguished triangle built on $s$ and then turned:
$P \xar{b} M \xar{s} N \xar{} \opn{T}(P)$.
We get an exact sequence 
\[ \opn{Hom}_{\cat{K}}(L, P) \xar{b \, \circ \, (-)} 
\opn{Hom}_{\cat{K}}(L, M) \xar{s \, \circ \, (-)} 
\opn{Hom}_{\cat{K}}(L, N) . \]
Since $s \circ a = 0$, there is $c : L \to P$ s.t.\ $a = b \circ c$. 
Now look at a distinguished triangle  built on $c$, and then turned:
$K \xar{t} L \xar{c} P \xar{} \opn{T}(K)$.
We know that $c \circ t = 0$; hence 
$a \circ t = b \circ c \circ t = 0$. 
But ($s \in \cat{S}$) $\Rightarrow$ ($H(\opn{T}^i(P)) = 0$ for all $i$)
$\Rightarrow$ ($t \in \cat{S}$).

The left versions (LD1) and (LD2) are proved the same way. 
\end{proof}

\begin{dfn} \label{dfn:1495}
A {\em denominator set of cohomological origin}
\index{Denominator set! of cohomological origin}
in $\cat{K}$ is a denominator set $\cat{S} \sub \cat{K}$ 
that arises from a cohomological functor $H$, as in Proposition \ref{prop:106}.
The morphisms in $\cat{S}$ are called {\em quasi-isomorphisms relative to $H$}.
\end{dfn}

\begin{thm} \label{thm:105}
Let $(\cat{K}, \opn{T})$ be a triangulated category, let $\cat{S}$ be a 
denominator set of cohomological origin in $\cat{K}$, 
\index{Denominator set! of cohomological origin}
and let 
$(\opn{Q}, \tau) : (\cat{K}, \opn{T}) \to
(\cat{K}_{\cat{S}}, \opn{T}_{\cat{S}})$
be the T-additive functor from Proposition \tup{\ref{prop:1840}}. Then the 
T-additive category 
$(\cat{K}_{\cat{S}}, \opn{T}_{\cat{S}})$
has a unique triangulated structure 
\index{Triangulated category}
such that these two properties 
hold\tup{:}

\begin{enumerate}
\rmitem{i} The pair $(\opn{Q}, \tau)$ is a triangulated functor. 
\index{Triangulated functor}

\rmitem{ii} Suppose 
$(\cat{K}', \opn{T}')$ 
is another triangulated category, and 
$(F, \nu) : (\cat{K}, \opn{T}) \to (\cat{K}', \opn{T}')$
is a triangulated functor, such that
$F(s)$ is invertible for every $s \in \cat{S}$. Let 
$(F_{\cat{S}}, \nu_{\cat{S}}) : (\cat{K}_{\cat{S}}, \opn{T}_{\cat{S}}) \to 
(\cat{K}', \opn{T}')$
be the T-additive functor from Proposition \tup{\ref{prop:1841}}.
Then $(F_{\cat{S}}, \nu_{\cat{S}})$ is a triangulated functor. 
\end{enumerate}
\end{thm}

This result is stated as \cite[Proposition I.3.2]{RD} and 
as \cite[Proposition 1.6.9]{KaSc1}. Both sources do not give proofs (just 
hints). 

\begin{proof} 
Since $\cat{S}$ is of cohomological origin we have 
$\opn{T}(\cat{S}) = \cat{S}$. Recall that the translation isomorphism $\tau$ is 
the identity automorphism of the functor $\opn{Q} \circ \opn{T}$; see 
Proposition \ref{prop:1840}. So we will ignore it. 

\medskip \noindent
Step 1.
The distinguished triangles in $\cat{K}_{\cat{S}}$ are defined to be those
triangles that are isomorphic to the images under $\opn{Q}$ of distinguished 
triangles in $\cat{K}$. Let us verify the axioms of triangulated 
category. 

\medskip \noindent
(TR1). By definition every triangle that's isomorphic to a distinguished 
triangle is distinguished; and the triangle 
$M \xar{\opn{id}_M} M \to 0 \to \opn{T}(M)$ 
in $\cat{K}_{\cat{S}}$ is clearly distinguished.

Suppose we are given a morphism $\al : L \to M$ in $\cat{K}_{\cat{S}}$. We have 
to build a distinguished triangle on it. 
Choose a fraction presentation 
$\al = \opn{Q}(a) \circ \opn{Q}(s)^{-1}$.
Using condition (LD1) we can find $b \in \cat{K}$ and $t \in 
\cat{S}$ such that $t \circ a = b \circ s$.
These fit into the solid commutative diagram
\[ \UseTips  \xymatrix @C=8ex @R=4ex { 
L
\ar@(ur,ul)@{-->}[rr]^{\al}
\ar[dr]_{b}
&
K
\ar[l]_{s}
\ar[r]^{a}
&
M
\ar[dl]^{t}
\\
&
\til{K}
}  \]
in $\cat{K}$. (The dashed arrow $\al$ is in $\cat{K}_{\cat{S}}$.) 

Consider the solid commutative diagram below, where
the rows are  distinguished triangles  built on $a$  and $b$ respectively.
\begin{equation} \label{eqn:1435}
 \UseTips  \xymatrix @C=8ex @R=6ex { 
K
\ar[r]^{a}
\ar[d]_{s}
&
M
\ar[r]^{e}
\ar[d]_{t}
& 
N
\ar[r]^{c}
\ar@{-->}[d]_{u}
&
\opn{T}(K)
\ar[d]_{\opn{T}(s)}
\\
L
\ar[r]^{b}
&
\til{K}
\ar[r]^{}
& 
P
\ar[r]^{d}
&
\opn{T}(L)
} 
\end{equation}
By (TR3) there is a morphism $u$ that makes the whole diagram commutative.
Since $s, t \in \cat{S}$ and $H$ is a cohomological functor, it follows that 
$u \in \cat{S}$. 
Applying the functor $\opn{Q}$ to (\ref{eqn:1435}), and using the 
isomorphism 
$\opn{Q}(t) : M \to \til{K}$ 
to replace $\til{K}$ with $M$,  we get the commutative diagram
\[ \UseTips  \xymatrix @C=8ex @R=6ex { 
K
\ar[r]^{\opn{Q}(a)}
\ar[d]_{\opn{Q}(s)}
&
M
\ar[r]^{\opn{Q}(e)}
\ar[d]_{\opn{Q}(\opn{id}_M)}
& 
N
\ar[r]^{\opn{Q}(c)}
\ar[d]_{\opn{Q}(u)}
&
\opn{T}(K)
\ar[d]_{\opn{T}(\opn{Q}(s))}
\\
L
\ar[r]^{\al}
&
M
\ar[r]^{\opn{Q}(u \circ e)}
& 
P
\ar[r]^{\opn{Q}(d)}
&
\opn{T}(L)
}  \]
in $\cat{K}_{\cat{S}}$. The top row is a distinguished triangle, and the
vertical arrows are isomorphisms. So the bottom row is a distinguished 
triangle. This is the triangle we were looking for. 

\medskip \noindent 
(TR2). Turning: this is trivial. 

\medskip \noindent 
(TR3). We are given the solid commutative diagram in $\cat{K}_{\cat{S}}$,
where the rows are distinguished triangles:
\begin{equation} \label{eqn:104}
\UseTips \xymatrix @C=8ex @R=6ex { 
L 
\ar[r]^{\al}
\ar[d]_{\phi}
&
M
\ar[r]^{\be}
\ar[d]_{\psi}
& 
N
\ar[r]^{\ga}
\ar@{-->}[d]_{\chi}
&
\opn{T}(L)
\ar[d]_{\opn{T}(\phi)}
\\
L' 
\ar[r]^{\al'}
&
M'
\ar[r]^{\be'}
& 
N'
\ar[r]^{\ga'}
&
\opn{T}(L') 
} 
\end{equation}
and we have to find $\chi$ to complete the diagram. 

By replacing the rows with isomorphic triangles, we can assume they come from
$\cat{K}$.
Thus we can replace (\ref{eqn:104}) with this diagram:
\begin{equation} \label{eqn:110}
\UseTips \xymatrix @C=8ex @R=6ex { 
L 
\ar[r]^{\opn{Q}(\al)}
\ar[d]_{\phi}
&
M
\ar[r]^{\opn{Q}(\be)}
\ar[d]_{\psi}
& 
N
\ar[r]^{\opn{Q}(\ga)}
\ar@{-->}[d]_{\chi}
&
\opn{T}(L)
\ar[d]_{\opn{T}(\phi)}
\\
L' 
\ar[r]^{\opn{Q}(\al')}
&
M'
\ar[r]^{\opn{Q}(\be')}
& 
N'
\ar[r]^{\opn{Q}(\ga')}
&
\opn{T}(L') 
} 
\end{equation}
in which $\al, \be, \ga, \al', \be', \ga'$ are morphisms in $\cat{K}$. 
It is a commutative diagram. Let us choose fraction presentations 
$\phi = \opn{Q}(a) \circ \opn{Q}(s)^{-1}$ and $\psi = \opn{Q}(b) \circ 
\opn{Q}(t)^{-1}$. 
Then the solid diagram (\ref{eqn:110}) comes from applying $\opn{Q}$ to the 
diagram 
\begin{equation} \label{eqn:103}
\UseTips \xymatrix @C=8ex @R=5.5ex { 
L 
\ar[r]^{\al}
&
M
\ar[r]^{\be}
& 
N
\ar[r]^{\ga}
&
\opn{T}(L)
\\
\til{L} 
\ar[d]_{a}
\ar[u]^{s}
&
\til{M}
\ar[d]_{b}
\ar[u]^{t}
& 
&
\opn{T}(\til{L})
\ar[d]_{\opn{T}(a)}
\ar[u]^{\opn{T}(s)}
\\
L' 
\ar[r]^{\al'}
&
M'
\ar[r]^{\be'}
& 
N'
\ar[r]^{\ga'}
&
\opn{T}(L') 
} \end{equation}
in $\cat{K}$. Here the rows are distinguished triangles in $\cat{K}$; but the 
diagram might fail to be commutative. 

By axiom (RO3) we can find $c \in \cat{K}$ and $u \in \cat{S}$ s.t.\ 
$\opn{Q}(t)^{-1} \circ \opn{Q}(\al) \circ \opn{Q}(s) = \opn{Q}(c) \circ 
\opn{Q}(u)^{-1}$.
This is the solid diagram below:
\[ \UseTips \xymatrix @C=8ex @R5.5ex { 
& &
L 
\ar[r]^{\al}
&
M
\\
\til{L}''
\ar@{-->}[r]^{u'}
&
\til{L}'
\ar[r]^{u}
\ar@(dr,dl)[rr]_(0.2){c}
&
\til{L} 
\ar[d]_(0.3){a}
\ar[u]^{s}
&
\til{M}
\ar[d]_{b}
\ar[u]^{t}
\\
& &
L' 
\ar[r]^{\al'}
&
M'
} \]
Thus 
$\opn{Q}(\al \circ s \circ u) = \opn{Q}(t \circ c)$.
By (RO4) there is a morphism $u' \in \cat{S}$ s.t.\ 
$(\al \circ s \circ u) \circ u' = (t \circ c) \circ u'$.
We get
\[ \phi = \opn{Q}(a) \circ \opn{Q}(s)^{-1} = \opn{Q}(a \circ u \circ u') \circ 
\opn{Q}(s \circ 
u \circ u')^{-1} \]
in $\cat{K}_{\cat{S}}$. 
Thus, after substituting $\til{L} :=  \til{L}''$,
 $s := s \circ u \circ u'$,
$a := a \circ u \circ u'$
and $c := c \circ u'$, we get a new diagram
\begin{equation} \label{eqn:105}
\UseTips \xymatrix @C=8ex @R=5.5ex { 
L 
\ar[r]^{\al}
&
M
\ar[r]^{\be}
& 
N
\ar[r]^{\ga}
&
\opn{T}(L)
\\
\til{L} 
\ar[r]^{c}
\ar[d]_{a}
\ar[u]^{s}
&
\til{M}
\ar[d]_{b}
\ar[u]^{t}
& 
&
\opn{T}(\til{L})
\ar[d]_{\opn{T}(a)}
\ar[u]^{\opn{T}(s)}
\\
L' 
\ar[r]^{\al'}
&
M'
\ar[r]^{\be'}
& 
N'
\ar[r]^{\ga'}
&
\opn{T}(L') 
} 
\end{equation}
in $\cat{K}$ instead of (\ref{eqn:103}). In this new diagram the top left 
square is commutative; but maybe the bottom left square is not commutative. 

When we apply $\opn{Q}$ to the diagram (\ref{eqn:105}), the  whole diagram, 
including the bottom left square, becomes commutative, since (\ref{eqn:110}) is
commutative. Again using condition (RO4), there is 
$v \in \cat{S}$ s.t.\ 
$(\al' \circ a) \circ v = (b \circ c) \circ v$.
In a diagram:
\[ \UseTips \xymatrix @C=8ex @R=5.5ex { 
& 
L 
\ar[r]^{\al}
&
M
\\
\til{L}'
\ar[r]^{v}
&
\til{L} 
\ar[d]_{a}
\ar[u]^{s}
\ar[r]^{c}
&
\til{M}
\ar[d]_{b}
\ar[u]^{t}
\\
& 
L' 
\ar[r]^{\al'}
&
M'
} \]
Performing the replacements
$\til{L} :=  \til{L}'$,
$s := s \circ v$,
$c := c \circ v$ and 
$a := a \circ v$ we now have a commutative square also at the bottom
left of (\ref{eqn:105}). Since $\ga \circ \be = 0$ and 
$\ga' \circ \be' = 0$,  
in fact the whole diagram (\ref{eqn:105}) in $\cat{K}$ is now commutative.

Now by (TR1) we can embed the morphism $c$ in a distinguished
triangle. We get the solid diagram
\begin{equation} \label{eqn:106}
\UseTips \xymatrix @C=8ex @R=5.5ex { 
L 
\ar[r]^{\al}
&
M
\ar[r]^{\be}
& 
N
\ar[r]^{\ga}
&
\opn{T}(L)
\\
\til{L} 
\ar[r]^{c}
\ar[d]_{a}
\ar[u]^{s}
&
\til{M}
\ar[r]^{\til{\be}}
\ar[d]_{b}
\ar[u]^{t}
& 
\til{N}
\ar[r]^{\til{\ga}}
\ar@{-->}[d]_{d}
\ar@{-->}[u]^{w}
&
\opn{T}(\til{L})
\ar[d]_{\opn{T}(a)}
\ar[u]^{\opn{T}(s)}
\\
L' 
\ar[r]^{\al'}
&
M'
\ar[r]^{\be'}
& 
N'
\ar[r]^{\ga'}
&
\opn{T}(L') 
} 
\end{equation}
in $\cat{K}$. The rows are distinguished triangles. 
Since $\til{\ga} \circ \til{\be} = 0$, the solid diagram is commutative.  
By (TR3) there are morphisms $w$ and $d$ that make the whole diagram
commutative. 
Now by the long exact cohomology sequence (see Theorem \ref{thm:3640})
the morphism $w$ belongs to $\cat{S}$. The morphism 
$\chi := \opn{Q}(d) \circ \opn{Q}(w)^{-1} :  N \to N'$
solves the problem.

\medskip \noindent 
(TR4). We will not give any details, as this axiom is not an important feature 
for us (see Remark \ref{rem:1281}).  A proof can be found as part of the proof 
of \cite[Theorem 2.1.8, page 97]{Ne1}, where this axiom is called (TR4$'$). 

\medskip \noindent
Step 2. Suppose $(F, \nu)$ is a triangulated functor as in condition (ii).
By Proposition \tup{\ref{prop:1841}} this extends uniquely to a 
T-additive functor $(F_{\cat{S}}, \nu_{\cat{S}})$. 
The construction of the triangulated structure on 
$(\cat{K}_{\cat{S}}, \opn{T}_{\cat{S}})$ in the previous steps,
and the defining property of the translation isomorphism $\nu_{\cat{S}}$ in 
Proposition \tup{\ref{prop:1841}}, show that
$(F_{\cat{S}}, \nu_{\cat{S}})$ is a triangulated functor. 

\medskip \noindent
Step 3. At this point $(\cat{K}_{\cat{S}}, \opn{T}_{\cat{S}})$ 
is a triangulated category, and conditions (i)-(ii) of the theorem are 
satisfied. We need to prove the uniqueness of the triangulated structure on
$(\cat{K}_{\cat{S}}, \opn{T}_{\cat{S}})$. Condition (i) says that we can't have 
less distinguished triangles than those we declared. We can't have more
distinguished triangles, because of condition (ii). 
\end{proof}

\begin{prop} \label{prop:1410}
Consider the situation of Proposition \tup{\ref{prop:106}} and Theorem 
\tup{\ref{thm:105}}. 
\begin{enumerate}
\item The cohomological functor 
$H : \cat{K} \to \cat{M}$ factors into 
$H = H_{\cat{S}} \circ \opn{Q}$, where 
$H_{\cat{S}} : \cat{K}_{\cat{S}} \to \cat{M}$ is a cohomological functor.

\item Let $M$ be an object of $\cat{K}$. The object $\opn{Q}(M)$ is zero 
in $\cat{K}_{\cat{S}}$ iff the objects $H(\opn{T}^i((M))$ are zero in $\cat{M}$ 
for all $i$. 
\end{enumerate}
\end{prop}

\begin{proof} \mbox{}

\smallskip \noindent
(1) The existence and uniqueness of the functor $H_{\cat{S}}$ are by the 
universal property (Loc3) in Definition \ref{dfn:1406}. 
We leave it as an exercise to show that $H_{\cat{S}}$ is a cohomological 
functor. 

\medskip \noindent 
(2) Since $H_{\cat{S}}$ is an additive functor, 
if $\opn{Q}(M) = 0$, then so is 
$H(M) = H_{\cat{S}}(\opn{Q}(M))$. And of course $\opn{Q}(M) = 0$ iff
$\opn{Q}(\opn{T}^i(M)) = 0$ for all $i$. 

For the converse, let $\phi: 0 \to M$ be the zero morphism in $\cat{K}$. 
If $H(\opn{T}^i(M)) = 0$ for all $i$, then 
$H(\opn{T}^i(\phi)) : 0 \to H(\opn{T}^i(M))$
are isomorphisms for all $i$. Then $\phi \in \cat{S}$, and so 
$\opn{Q}(\phi) : 0 \to \opn{Q}(M)$
is an isomorphism in  $\cat{K}_{\cat{S}}$.
\end{proof}

\begin{prop} \label{prop:1480}
Let $\cat{K}$ be a triangulated category, let 
$\cat{S}$ be a denominator set of cohomological origin in $\cat{K}$, and let 
$\cat{K}'$ be a full triangulated subcategory of $\cat{K}$. 
Then $\cat{S}' := \cat{K}' \cap \, \cat{S}$ is a denominator set of 
cohomological origin in $\cat{K}'$, the Ore localization 
$\cat{K}'_{\cat{S}'}$ exists, and $\cat{K}'_{\cat{S}'}$ is a triangulated 
category. 
\end{prop}

\begin{proof}
Let $H :  \cat{K} \to \cat{M}$ be a cohomological functor that determines 
$\cat{S}$. The functor $H|_{\cat{K}'} : \cat{K}' \to \cat{M}$ is also 
cohomological, and the set of morphisms $\cat{S}'$ satisfies 
\[ \cat{S}' = \bigl\{ s \in  \cat{K}'
\mid  H|_{\cat{K}'}(\opn{T}^i(s))  \tup{ is an isomorphism for all } i \bigr\}  
. \]
Hence Proposition \ref{prop:106} and Theorem \ref{thm:105} apply.
\end{proof}

In the situation of the proposition, the localization functor is denoted by
$\opn{Q}' : \cat{K}' \to \cat{K}'_{\cat{S}'}$.

\begin{prop} \label{prop:1496}
In the situation of Proposition \tup{\ref{prop:1480}}, 
let $F : \cat{K}' \to \cat{E}$ be a triangulated functor into some 
triangulated category $\cat{E}$. Assume that for every 
$s \in \cat{S}'$, the morphism $F(s)$ is an isomorphism in $\cat{E}$. 
Then there is a unique triangulated functor 
$F_{\cat{S}'} : \cat{K}'_{\cat{S}'} \to \cat{E}$
that extends $F$\tup{;} Namely 
$F_{\cat{S}'} \circ \opn{Q}' = F$ as functors $\cat{K}' \to \cat{E}$.
\end{prop}

\begin{proof}
This is part of Theorem \ref{thm:105}. 
\end{proof}

In particular we can look at the functor  
$F : \cat{K}' \xar{\opn{inc}} \cat{K} \xar{\opn{Q}} \cat{K}_{\cat{S}}$, 
and its extension 
$F_{\cat{S}'} : \cat{K}'_{\cat{S}'} \to \cat{K}_{\cat{S}}$.
We are interested in sufficient conditions for the functor 
$F_{\cat{S}'}$ to be fully faithful.

\begin{prop} \label{prop:1320}  
Let $\cat{K}$ be a triangulated category, let 
$\cat{S}$ be a denominator set of cohomological origin in $\cat{K}$,
and let $\cat{K}' \subseteq \cat{K}$ be a full triangulated subcategory.  
Define 
$\cat{S}' := \cat{K}' \cap \, \cat{S}$. 
Assume either of these conditions holds\tup{:}
\begin{itemize}
\rmitem{r} Let $M \in \opn{Ob}(\cat{K})$. If 
there exists a morphism $s : M \to L$ in $\cat{S}$ with 
$L \in \opn{Ob}(\cat{K}')$, 
there exists a morphism $t : K \to M$ in $\cat{S}$ with 
$K \in \opn{Ob}(\cat{K}')$.

\rmitem{l} Let $M \in \opn{Ob}(\cat{K})$. If 
there exists a morphism $s : L \to M$ in $\cat{S}$ with 
$L \in \opn{Ob}(\cat{K}')$, 
there exists a morphism $t : M \to K$ in $\cat{S}$ with 
$K \in \opn{Ob}(\cat{K}')$.
\end{itemize}
Then the functor 
$F_{\cat{S}'} : \cat{K}'_{\cat{S}'} \to \cat{K}_{\cat{S}}$ 
is fully faithful.
\end{prop}

\begin{proof}
By  Proposition \ref{prop:1480}, $\cat{S}'$ is a multiplicatively closed 
set of cohomological origin. So by Proposition \ref{prop:106}
it is a right and left denominator set in $\cat{K}'$. 
According to Proposition \ref{prop:3380}, under condition (r) the functor 
$F_{\cat{S}'}$ is  fully faithful.

For condition (l) we pass to the opposite categories, using Proposition 
\ref{prop:2310}. 
\end{proof}

\begin{rem} \label{rem:4220}
Let $\cat{K}$ be a triangulated category, and let 
$\cat{N} \sub \cat{K}$ be a full triangulated subcategory. Define 
$\cat{S}$ to be the set of morphisms $s$ in $\cat{K}$ such that the 
cone of $s$ (see Definition \ref{dfn:3600}) is isomorphic to an object of 
$\cat{N}$.  It turns out that $\cat{S}$ is a left and right denominator set in 
$\cat{K}$. We call $\cat{S}$ the {\em denominator set associated to $\cat{N}$}. 
The localization $\cat{K}_{\cat{S}}$ is a triangulated 
category, and it is called the {\em Verdier quotient}%
\index{Verdier quotient}
of $\cat{K}$ by $\cat{N}$, with notation 
$\catt{K} / \catt{N} :=  \cat{K}_{\cat{S}}$. 
See \cite[Chapter 2]{Ver} or \cite[Section 1.6]{KaSc1}. 

We want to relate this to Remark \ref{rem:4777}. 
Let $\til{\cat{S}}$ be a left and right denominator set in $\cat{K}$.
A morphism $s$ in $\cat{K}$ is called a {\em divisor} of $\til{\cat{S}}$ 
if there exist morphisms $u, v \in \cat{K}$ such that 
$u \circ s$ and $s \circ v$ belong to $\til{\cat{S}}$.
We call $\til{\cat{S}}$ a {\em saturated denominator set} if 
it contains all its divisors.

Let $\cat{N} \sub \cat{K}$ be a full triangulated subcategory, and let 
$\cat{S}$ be the associated denominator set, as in the first 
paragraph. Let  $\til{\cat{N}}$ be the saturated 
full triangulated subcategory of $\cat{K}$ generated by $\cat{N}$, in the 
sense of Definition \ref{dfn:4916}. Let $\til{\cat{S}}$ be the set of morphisms
gotten by adjoining to $\cat{S}$ all its divisors. 
Then $\til{\cat{S}}$ is a saturated left and right denominator set of morphisms 
in $\cat{K}$, the canonical functor 
$\cat{K}_{\cat{S}} \to \cat{K}_{\til{\cat{S}}}$
is an isomorphism, the kernel of the localization functor 
$\opn{Q} : \cat{K} \to \cat{K}_{\cat{S}}$
(see Corollary \ref{cor:4920}) is $\til{\cat{N}}$,
and $\til{\cat{S}}$ is the denominator set associated to $\til{\cat{N}}$.
This is proved in \cite[Chapter 2]{Ver} and \cite[Proposition 7.1.20]{KaSc2}.

An important example of a Verdier quotient is the passage from 
the homotopy category $\dcat{K}(A, \cat{M})$ to the derived category 
$\dcat{D}(A, \cat{M})$, that we shall study in Subsection 
\ref{subsec:def-der-cat} below. Indeed, the set of quasi-isomorphisms 
$\dcat{S}(A, \cat{M})$ in $\dcat{K}(A, \cat{M})$
can be described as the denominator set associated to the 
full subcategory $\dcat{N}(A, \cat{M}) \sub \dcat{K}(A, \cat{M})$ on the 
acyclic complexes, which is a saturated full triangulated subcategory. 
\end{rem}



\mysubsection{Definition of the Derived Category} \label{subsec:def-der-cat}

We now specialize to the triangulated category 
$\cat{K} = \dcat{K}(A, \cat{M})$, for a central DG $\K$-ring $A$ and 
a $\K$-linear abelian category $\cat{M}$. 

Recall the cohomology functor 
$\opn{H} : \dcat{C}_{\mrm{str}}(A, \cat{M}) \to 
\dcat{G}_{\mrm{str}}(\cat{M})$
from Definition \ref{dfn:2993}. 

\begin{lem} \label{lem:2320}
Suppose $\phi, \psi : M \to N$ are morphisms in 
$\dcat{C}_{\mrm{str}}(A, \cat{M})$ such that 
$\psi - \phi$ is a coboundary in $\opn{Hom}_{A, \cat{M}}(M, N)$. 
Then $\opn{H}(\phi) = \opn{H}(\psi)$, as morphisms 
$\opn{H}(M) \to \opn{H}(N)$ in $\dcat{G}_{\mrm{str}}(\cat{M})$. 
\end{lem}

\begin{exer} \label{exer:2320}
Prove the lemma. 
\end{exer}

It follows that there is an induced functor 
\begin{equation} \label{eqn:4985}
\opn{H} : \dcat{K}(A, \cat{M}) \to \dcat{G}_{\mrm{str}}(\cat{M}) . 
\end{equation}
We shall be interested in its degree $0$ component 
$\mrm{H}^0$. Of course all other components can be recovered from $\mrm{H}^0$
by the translation functor: 
$\mrm{H}^i = \mrm{H}^0 \circ \opn{T}^i$.

\begin{prop} \label{prop:1475}
Let $\cat{M}$ be an abelian category and let $A$ be a DG ring. Then the functor 
$\mrm{H}^0 : \dcat{K}(A, \cat{M}) \to \cat{M}$
is cohomological. 
\end{prop}

\begin{proof}
Clearly $\mrm{H}^0$ is additive. Consider a distinguished triangle 
\begin{equation} \label{eqn:1437}
L \xar{\al} M \xar{\be} N \xar{\ga} \opn{T}(L)
\end{equation}
in $\dcat{K}(A, \cat{M})$. We must prove that 
\[ \opn{H}^{0}(L) \xar{\opn{H}^0(\al)} \opn{H}^{0}(M)
\xar{\opn{H}^0(\be)} \opn{H}^{0}(N) \]
is an exact sequence in $\cat{M}$. 

By Definition \ref{dfn:140} we can assume that the distinguished triangle 
(\ref{eqn:1437}) is the image of a standard triangle in  
$\dcat{C}(A, \cat{M})$. Namely that $N = \opn{Cone}(\al)$, the standard 
cone associated to $\al$, and the morphisms are 
$\be = e_{\al}$ and $\ga = p_{\al}$; see Definitions  \ref{dfn:1172} and 
\ref{dfn:1140}. To be explicit, in matrix notation we have 
$N = \sbmat{M \\[0.2em] \opn{T}(L)}$, 
$\d_N = \sbmat{\d_M & \al \circ \, \opn{t}_L^{-1} 
\\[0.2em] 0 & \d_{\opn{T}(L)}}$ 
and 
$\be = \sbmat{\opn{id}_M \\[0.2em] 0}$.

Let us use the abbreviations $\bar{\be} := \opn{H}^0(\be)$ and
$\bar{\al} := \opn{H}^0(\al)$. 
By Proposition \ref{prop:1280} we know that $\be \circ \al = 0$, and hence 
$\bar{\be} \circ \bar{\al} = 0$. 
It remains to prove that the morphism 
$\bar{\al}  : \opn{H}^0(L) \to \opn{Ker}(\bar{\be})$
in $\cat{M}$ is an epimorphism. 
This will be done using the first sheaf trick (Proposition \ref{prop:3635}(3)):
we will prove that given a section 
$\bar{m} \in \Ga(U, \opn{Ker}(\bar{\be}))$
on an open set $U \in \cat{M}$, there is a refinement
$V' \surj U$ and a section 
$\bar{l} \in \Ga(V', \opn{H}^0(L))$
such that $\bar{\al}(\bar{l}) = \bar{m}$
in $\Ga(V', \opn{Ker}(\bar{\be}))$. 
See Subsection \ref{subsec:sheaf-tricks} for the ``geometric'' terminology 
that's used in the sheaf tricks. 

Take some section 
$\bar{m} \in \Ga(U, \opn{Ker}(\bar{\be}))$
on an open set $U \in \cat{M}$.
This means that 
$\bar{m} \in \Ga(U, \opn{H}^0(M))$ and 
$\bar{\be}(\bar{m}) = 0$. 
There is a covering $V \surj U$ and a section 
$m \in \Ga(V, \opn{Z}^{0}(M))$ such that 
$\pi_{M}(m) = \bar{m}$. Now 
\[ \be(m) = \bmat{m \\[0.1em] 0} \in \Ga(V, \opn{Z}^{0}(N)) \sub 
\Ga(V, N^{0}) = 
\bmat{ \Ga(V, M^{0}) \\[0.1em] \Ga(V, L^{1})}  \]
in column notation. 
The vanishing of $\bar{\be}(\bar{m}) = \pi_N(\be(m))$ in 
$\Ga(V, \opn{H}^{0}(N))$ means that there is a covering 
$V' \surj V$ and a section 
$n \in \Ga(V', N^{-1})$
such that $\d_N(n) = \be(m)$. 
In column notation we have 
$n = \sbmat{m' \\[0.1em] l}$,
with $m' \in \Ga(V', M^{-1})$ and $l \in \Ga(V', L^{0})$. 
The equality $\be(m) = \d_N(n)$ becomes the matrix equality 
\[ \bmat{m \\[0.1em] 0} =
\bmat{\d_M & \al \\[0.1em] 0 & -\d_{L}} \cd 
\bmat{m' \\[0.1em] l} = 
\bmat{ \d_M(m') + \al(l) \\[0.1em] -\d_{L}(l) } . \]
We see that $l \in \Ga(V', \opn{Z}^{0}(L))$. 
Then the cohomology class 
$\bar{l} := \pi_L(l) \in \Ga(V', \opn{H}^{0}(L))$
satisfies 
\[ \bar{\al}(\bar{l}) = \pi_M(\al(l)) = \pi_M(\d_M(m') + \al(l)) = 
\pi_M(m) = \bar{m} \]
in $\Ga(V', \opn{H}^{0}(M))$, as required. 
\end{proof} 

\begin{dfn} \label{dfn:1475}
A morphism $\phi$ in  $\dcat{K}(A, \cat{M})$ is called a 
{\em quasi-isomorphism}
\index{Quasi-isomorphism}
if the morphisms $\opn{H}^i(\phi)$ in $\cat{M}$ are 
isomorphisms for all $i$. 
The set of quasi-isomorphisms in $\dcat{K}(A, \cat{M})$ is denoted by 
$\dcat{S}(A, \cat{M})$. 
\end{dfn}

By Proposition \ref{prop:1475} the functor $\opn{H}^0$ is cohomological.
According to Proposition \ref{prop:106}, $\dcat{S}(A, \cat{M})$ is a 
denominator set of cohomological origin; so Theorem \ref{thm:105} applies to 
it, 
and the next definition makes sense. 

\begin{dfn}
Let $\cat{M}$ be a $\K$-linear abelian category and $A$ a central DG $\K$-ring.
The {\em derived category of DG $A$-modules in $\cat{M}$}%
\index{Derived category! of DG $A$-modules in $\cat{M}$}%
\index{1-D(A,M)@$\dcat{D}(A, \cat{M})$}%
\index{1-D(A)@$\dcat{D}(A)$}%
\index{1-D(M)@$\dcat{D}(\cat{M})$}
is the $\K$-linear triangulated category 
\[ \dcat{D}(A, \cat{M}) := \dcat{K}(A, \cat{M})_{\dcat{S}(A, \cat{M})} . \]
The corresponding triangulated localization functor is 
\[ \opn{Q} : \dcat{K}(A, \cat{M}) \to \dcat{D}(A, \cat{M}) . \]
\end{dfn}

We also have the additive functor
$\opn{P} : \dcat{C}_{\mrm{str}}(A, \cat{M}) \to \dcat{K}(A, \cat{M})$
which sends a strict morphism of DG modules to its homotopy class.

\begin{dfn} \label{dfn:1480}
Let $\cat{M}$ be an abelian category and let $A$ be a DG ring.
Define the functor 
\[ \til{\opn{Q}} := \opn{Q} \circ \opn{P} : 
\dcat{C}_{\mrm{str}}(A, \cat{M}) \to \dcat{D}(A, \cat{M}) . \]
\end{dfn}

We thus get a commutative diagram of additive functors
\begin{equation} \label{eqn:5025}
\UseTips \xymatrix @C=8ex @R=6ex {
\dcat{C}_{\mrm{str}}(A, \cat{M}) 
\ar[r]^{\opn{P}}
\ar@(u,u)[rr]^{\til{\opn{Q}}}
&
\dcat{K}(A, \cat{M})
\ar[r]^{\opn{Q}}
&
\dcat{D}(A, \cat{M})
}
\end{equation}
that are all the identity on objects.

It is sometimes convenient to describe morphisms in $\dcat{D}(A, \cat{M})$
in terms of the functor $\til{\opn{Q}}$. A morphism 
$s \in \dcat{C}_{\mrm{str}}(A, \cat{M})$
is called a quasi-isomorphism if $\opn{P}(s)$ is a quasi-isomorphism
in $\dcat{K}(A, \cat{M})$; i.e.\ if $\opn{H}(s)$ is an isomorphism in
$\dcat{G}_{\mrm{str}}(\cat{M})$. 

\begin{prop} \label{prop:1481}
\mbox{}
\begin{enumerate}
\item Every morphism $\phi$ in $\dcat{D}(A, \cat{M})$ can be written as a right 
fraction 
$\phi = \til{\opn{Q}}(a) \circ \til{\opn{Q}}(s)^{-1}$
where $a, s \in \dcat{C}_{\mrm{str}}(A, \cat{M})$ and 
$s$ is a quasi-isomorphism. 

\item Let $a \in \dcat{C}_{\mrm{str}}(A, \cat{M})$. 
Then $\til{\opn{Q}}(a) = 0$ in $\dcat{D}(A, \cat{M})$ iff 
there exists a quasi-isomorphism $s$ in $\dcat{C}_{\mrm{str}}(A, \cat{M})$
such that $a \circ s$ is a coboundary in $\dcat{C}(A, \cat{M})$.
\end{enumerate}
\end{prop}

\begin{proof} \mbox{}

\smallskip \noindent
(1) This is because of property (RO3) of Definition \ref{dfn:1407} and the fact 
that $\opn{P}$ is full. 

\medskip \noindent 
(2) Let us write $\bar{a} := \opn{P}(a) \in \dcat{K}(A, \cat{M})$.
Since $\opn{Q}(\bar{a}) = 0$, by Property (RO4) of Definition \ref{dfn:1407}
there is a quasi-isomorphism 
$\bar{s} \in \dcat{K}(A, \cat{M})$
such that $\bar{a} \circ \bar{s} = 0$ in $\dcat{K}(A, \cat{M})$.
Choose a quasi-isomorphism $s \in \dcat{C}_{\mrm{str}}(A, \cat{M})$
such that $\bar{s} = \opn{P}(s)$. Then $\opn{P}(a \circ s) = 0$, and 
this means that $a \circ s$ is a coboundary. 
\end{proof}

Of course there is a left version of this proposition.
The next definition is about arbitrary (not necessarily linear) categories. 

\begin{dfn} \label{4980}
A functor $F : \cat{C} \to \cat{D}$ between categories is called  
{\em conservative} if for every morphism $\phi$ in $\cat{C}$, 
$\phi$ is an isomorphism if and only if $F(\phi)$ is an isomorphism. 
\end{dfn}

By Proposition \ref{prop:1410} the cohomology functor (\ref{eqn:4985}) extends 
uniquely to a functor 
\begin{equation} \label{eqn:4986}
\opn{H} : \dcat{D}(A, \cat{M}) \to \dcat{G}_{\mrm{str}}(\cat{M}) . 
\end{equation}

\begin{cor} \label{cor:2145}
The functor $\opn{H}$ in \tup{(\ref{eqn:4986})} is conservative.
\end{cor}

\begin{proof}
One implication is trivial. For the other implication, assume that 
$\phi$ is a morphism in $\dcat{D}(A, \cat{M})$ such that 
$\opn{H}(\phi)$ is an isomorphism. We can write $\phi$ as a right fraction: 
$\phi = \opn{Q}(a) \circ \opn{Q}(s)^{-1}$, 
where $a \in \dcat{K}(A, \cat{M})$ and 
$s \in \dcat{S}(A, \cat{M})$.
Then 
$\opn{H}(\phi) = \opn{H}(\opn{Q}(a)) \circ \opn{H}(\opn{Q}(s))^{-1}$,
and we see that $\opn{H}(\opn{Q}(a))$ is an isomorphism.
But $\opn{H}(a) = \opn{H}(\opn{Q}(a))$, so in fact 
$a \in \dcat{S}(A, \cat{M})$ too. Therefore 
$\opn{Q}(a)$ is an isomorphism in 
$\dcat{D}(A, \cat{M})$. It follows that $\phi$ is an isomorphism in 
$\dcat{D}(A, \cat{M})$.
\end{proof}

\begin{exer} \label{exer:1846}
Here $\cat{M} = \dcat{M}(\K)$, so 
$\dcat{K}(A, \cat{M}) = \dcat{K}(A)$. Show that the functor 
$\opn{H}^0 : \dcat{K}(A) \to \dcat{M}(\K)$ is corepresentable 
by the object $A \in \dcat{K}(A)$ (see Subsection 
\ref{subsec:repfunc}). 
\end{exer}

\begin{rem} \label{rem:4971}
This is a continuation of Remarks \ref{rem:4960} and \ref{rem:4970} on strictly 
unital 
$\mrm{A}_{\infty}$ categories% 
\index{Ainfty@$\mrm{A}_{\infty}$ category}.
The base ring $\K$ is a field. 
The references are \cite{Kel5} and \cite{COS}.
Recall that if $\cat{A}$ is an  
$\mrm{A}_{\infty}$ category and $\cat{B}$ is a DG category, then 
$\cat{A_{\infty}Fun}(\cat{A}, \cat{B})$
is a DG category. 

Consider the special case is when $\cat{B} = \dcat{C}(\K)$, the DG category of 
complexes of $\K$-modules. An $\mrm{A}_{\infty}$ functor 
$M : \cat{A} \to \dcat{C}(\K)$ is called an {\em $\mrm{A}_{\infty}$ 
$\cat{A}$-module}. 
The $\mrm{A}_{\infty}$ $\cat{A}$-modules are the objects of the DG category 
$\dcat{C}_{\infty}(\cat{A}) := \cat{A_{\infty}Fun}(\cat{A}, \dcat{C}(\K))$. 
The strict subcategory $\opn{Str}(\dcat{C}_{\infty}(\cat{A}))$, in the sense of
Definition \ref{dfn:1212}(1), is what most texts call the {\em category 
of $\mrm{A}_{\infty}$ $\cat{A}$-modules}. 

A morphism $\phi : M \to N$ in $\opn{Str}(\dcat{C}_{\infty}(\cat{A}))$
is called a {\em quasi-isomorphism} if for every object 
$x \in \opn{Ob}(\cat{A})$ the morphism 
$\phi(x) : M(x) \to N(x)$ in 
$\opn{Str}(\dcat{C}(\K)) = \dcat{C}_{\mrm{str}}(\K)$ is a 
quasi-isomorphism. The localization of 
$\opn{Str}(\dcat{C}_{\infty}(\cat{A}))$ w.r.t.\ 
the quasi-isomorphisms is the {\em derived category of $\mrm{A}_{\infty}$
$\cat{A}$-modules}, denoted by $\dcat{D}_{\infty}(\cat{A})$.
(As usual, the tradition is to call $\dcat{D}_{\infty}(\cat{A})$ the ``homotopy 
category'' of $\dcat{C}_{\infty}(\cat{A})$, and to denote 
it by $\opn{Ho}(\dcat{C}_{\infty}(\cat{A}))$.)
There is a localization functor 
$\opn{Q} : \opn{Str}(\dcat{C}_{\infty}(\cat{A})) \to 
\dcat{D}_{\infty}(\cat{A})$.

We can also consider the category 
$\dcat{K}_{\infty}(\cat{A}) := \opn{H}^0(\dcat{C}_{\infty}(\cat{A}))$, see 
Definition \ref{dfn:1212}(2). There is a full functor 
$\opn{P} : \opn{Str}(\dcat{C}_{\infty}(\cat{A})) \to 
\dcat{K}_{\infty}(\cat{A})$,
and a commutative diagram 
\begin{equation} \label{eqn:4978}
\UseTips \xymatrix @C=10ex @R=6ex {
\opn{Str}(\dcat{C}_{\infty}(\cat{A}))
\ar[r]^{\opn{P}}
\ar@(u,u)[rr]^{\opn{Q}}
&
\dcat{K}_{\infty}(\cat{A})
\ar[r]^{\bar{\opn{Q}}}
&
\dcat{D}_{\infty}(\cat{A})
}
\end{equation}
All three functors are the identity on objects. 
The categories $\dcat{K}_{\infty}(\cat{A})$ and 
$\dcat{D}_{\infty}(\cat{A})$ are triangulated, and the functor $\bar{\opn{Q}}$ 
is an isomorphism of triangulated categories. 

Finally, let $\cat{A}$ be a DG category. Then the DG category 
$\dcat{C}_{}(\cat{A})$ of DG $\cat{A}$-modules embeds faithfully (but not 
fully) into the DG category $\dcat{C}_{\infty}(\cat{A})$.
The induced triangulated functor 
$\dcat{D}_{}(\cat{A}) \to \dcat{D}_{\infty}(\cat{A})$
is an equivalence. 
\end{rem}




\mysubsection{Boundedness Conditions in
\texorpdfstring{$\dcat{K}(A, \cat{M})$}{K(A,M)}}
\label{subsec:bd-in-homot}

We continue with a central DG $\K$-ring $A$ and a $\K$-linear abelian category 
$\cat{M}$. 
Boundedness conditions in $\dcat{C}(A, \cat{M})$ were introduced in Definition 
\ref{dfn:3220}. Similarly: 

\begin{dfn} \label{dfn:3221}
\index{1-K(A,M)@$\dcat{K}^{\star}(A, \cat{M})$}
\index{Boundedness! condition}
We define $\dcat{K}^-(A, \cat{M})$, $\dcat{K}^+(A, \cat{M})$ and 
$\dcat{K}^{\mrm{b}}(A, \cat{M})$ to be the full subcategories of 
$\dcat{K}(A, \cat{M})$ consisting of bounded above, bounded below and bounded
DG modules respectively. 
\end{dfn}

Of course 
$\dcat{K}^{\mrm{b}}(A, \cat{M}) = \dcat{K}^-(A, \cat{M})
\cap \dcat{K}^+(A, \cat{M})$.
The subcategories $\dcat{K}^{\star}(A, \cat{M})$, 
for $\star \in \{ -, +, \mrm{b} \}$, 
are full triangulated subcategories of $\dcat{K}(A, \cat{M})$; this is 
because the operations of translation 
and cone preserve the various boundedness conditions. 
Note that 
$\dcat{K}^{\star}(A, \cat{M}) = 
\opn{Ho} \bigl( \dcat{C}^{\star}(A, \cat{M}) \bigr)$.
As the next example shows, sometimes the categories 
$\dcat{C}^{\star}(A, \cat{M})$ and 
$\dcat{K}^{\star}(A, \cat{M})$ can be very degenerate. 

\begin{exa} \label{exa:1460}
Let $A$ be the DG ring $\K[t, t^{-1}]$, the ring of Laurent polynomials in the  
variable $t$ of degree $1$, with the zero differential. If 
$M = \{ M^i \}_{i \in \Z}$ is a nonzero object of 
$\dcat{C}(A, \cat{M})$, then $M^i \neq 0$ for all $i$. Therefore 
the categories $\dcat{C}^{\star}(A, \cat{M})$ and 
$\dcat{K}^{\star}(A, \cat{M})$ are zero for 
$\star \in \{ -, +, \mrm{b} \}$. 
\end{exa}

Let 
$\dcat{S}^{\star}(A, \cat{M}) := \dcat{K}^{\star}(A, \cat{M}) \cap
\dcat{S}(A, \cat{M})$,
the category of quasi-isomorphisms in $\dcat{K}^{\star}(A, \cat{M})$. 
As already mentioned, Theorem \ref{thm:105} applies here, so we can 
localize. 

\begin{dfn} \label{dfn:1460}
For $\star \in \{ -, +, \mrm{b} \}$ we define
\[  \dcat{D}^{\star}(A, \cat{M}) := 
\dcat{K}^{\star}(A, \cat{M})_{\dcat{S}^{\star}(A, \cat{M})} , \]
\index{1-D(A,M)@$\dcat{D}^{\star}(A, \cat{M})$}
\index{Boundedness! condition}
the Ore localization of $\dcat{K}^{\star}(A, \cat{M})$ with respect to 
$\dcat{S}^{\star}(A, \cat{M})$.
\end{dfn}

Here is another kind of boundedness condition. 

\begin{dfn} \label{dfn:1461}
For $\star \in \{ -, +, \mrm{b} \}$ we define
$\dcat{D}(A, \cat{M})^{\star}$
\index{1-D(A,M)@$\dcat{D}(A, \cat{M})^{\star}$}
\index{Boundedness! condition}
to be the full subcategory of $\dcat{D}(A, \cat{M})^{}$
on the complexes $M$ whose cohomology $\opn{H}(M)$ is of boundedness type 
$\star$. 
\end{dfn}

Of course $\dcat{D}(A, \cat{M})^{\star}$ is a full triangulated subcategory 
of $\dcat{D}(A, \cat{M})$. 
See Remark \ref{rem:1461} regarding Definitions \ref{dfn:1460} and 
\ref{dfn:1461}. 

Handling boundedness conditions requires the use of {\em truncations}. These 
will be introduced after the next proposition. 

\begin{prop} \label{prop:2165}
Let 
$0 \to L \xar{\phi} M \xar{\psi} N \to 0$ 
be a short exact sequence in $\dcat{C}_{\mrm{str}}(A, \cat{M})$. 
Then there is a morphism $\th : N \to \opn{T}(L)$ in $\dcat{D}(A, \cat{M})$ 
such that 
$L \xar{\opn{Q}(\phi)} M \xar{\opn{Q}(\psi)} N \xar{\th} \opn{T}(L)$
is a distinguished triangle in $\dcat{D}(A, \cat{M})$.
\end{prop}

\begin{proof}
We are following the proof of \cite[Proposition 1.7.5]{KaSc1}.
Let $\til{N}$ be the standard cone on $\phi$. In matrix notation, as 
in Definition \ref{dfn:1172}, we have
$\til{N} = \sbmat{M \\[0.2em] \opn{T}(L)}$
and 
$\d_{\til{N}} =
\sbmat{\d_M & \, \phi \, \circ \, \opn{t}^{-1} 
\\[0.2em] 0 & \, \d_{\opn{T}(L)}}$.
The object $\til{N}$ sits inside the standard triangle 
$L \xar{\phi} M \xar{\til{\psi}} \til{N} \xar{\til{\chi}} \opn{T}(L)$
in  $\dcat{C}_{\mrm{str}}(A, \cat{M})$, where 
$\til{\psi} := \sbmat{\opn{id} \\[0.2em] 0}$ 
and 
$\til{\chi} := \sbmat{0 & \, \opn{id}}$
in matrix notation. Define the morphism $\ga = \til{N} \to N$ to be the matrix
$\ga := \sbmat{\psi & \, 0}$. We get a commutative diagram
\[ \UseTips \xymatrix @C=8ex @R=5.5ex {
L
\ar[r]^{\phi}
&
M
\ar[r]^{\til{\psi}}
\ar[dr]_{\psi} 
&
\til{N}
\ar[r]^{\til{\chi}} 
\ar[d]^{\ga}
&
\opn{T}(L)
\\
&
&
N
} \]
in  $\dcat{C}_{\mrm{str}}(A, \cat{M})$.
We shall prove below that $\ga$ is a quasi-isomorphism. 
Then the morphism 
$\th := \opn{Q}(\til{\chi}) \circ \opn{Q}(\ga)^{-1} : N \to \opn{T}(L)$ 
will work. 

Let $\til{K}$ be the standard cone on $\opn{id}_L$, and let 
$\til{\be} : \til{K} \to \til{N}$ be the matrix morphism 
$\sbmat{\phi & \, 0 \\[0.1em] 0 & \, \opn{id}} : 
\sbmat{L \\[0.2em] \opn{T}(L)} \to
\sbmat{M \\[0.2em] \opn{T}(L)}$.
This fits into a short exact sequence 
$0 \to \til{K} \xar{\til{\be}} \til{N} \xar{\ga} N \to 0$
in $\dcat{C}_{\mrm{str}}(A, \cat{M})$.
But the DG module $\til{K}$ is acyclic, and therefore $\ga$ is a 
quasi-isomorphism, as claimed. 
\end{proof}

\begin{dfn} \label{dfn:2320}
Let $M \in \dcat{C}(\cat{M})$ and $i \in \Z$. The 
{\em smart truncation of $M$ below $i$}%
\index{Truncation of a complex! smart}%
\index{1-Smtleq)@$\opn{smt}^{\leq i}(M)$}
is the complex 
\[ \opn{smt}^{\leq i}(M) := 
\bigl( \cdots \to M^{i-2} \xar{\d} M^{i-1} \xar{\d} \opn{Z}^i(M) \to 0 \to
\cdots \bigr) .  \]
The {\em smart truncation of $M$ above $i$}%
\index{1-Smtgeq)@$\opn{smt}^{\geq i}(M)$}
is the complex 
\[ \opn{smt}^{\geq i}(M) := 
\bigl( \cdots \to 0 \to \opn{Y}^{i}(M) \xar{\d} M^{i+1} \xar{\d} M^{i+2}
\to \cdots \bigr) . \]
\end{dfn}

In the definition we use the objects of cocycles $\opn{Z}^i(M)$
and decocycles $\opn{Y}^i(M)$ from Definition \ref{dfn:2993}. Note that 
$\opn{smt}^{\leq i}(M)$ is a subcomplex of $M$, whereas 
$\opn{smt}^{\geq i}(M)$ a quotient complex of $M$. Denoting the inclusion and 
the projection by $e$ and $p$ respectively, we get a diagram 
\begin{equation} \label{eqn:3035}
\opn{smt}^{\leq i}(M) \xar{e} M \xar{p} \opn{smt}^{\geq i + 1}(M) 
\end{equation}
in $\dcat{C}_{\mrm{str}}(\cat{M})$. 
This diagram is functorial in $M \in \dcat{C}_{\mrm{str}}(\cat{M})$.
But diagram (\ref{eqn:3035}) is not an exact sequence in general  -- there is a 
defect in degree $i$. 

Recall that a DG ring $A$ is called nonpositive if $A^i = 0$ for all $i > 0$. 

\begin{prop} \label{prop:3170}
Assume $A$ is a nonpositive DG ring. 
\begin{enumerate}
\item The differential of every 
$M \in \dcat{C}_{\mrm{str}}(A, \cat{M})$ is $A^0$-linear.

\item The smart truncations from Definition \tup{\ref{dfn:2320}}
are functors from $\dcat{C}_{\mrm{str}}(A, \cat{M})$ to itself. 
\end{enumerate}
\end{prop}

\begin{exer} \label{exer:3170}
Prove this proposition. 
\end{exer}

\begin{prop} \label{prop:2320}
Assume $A$ is nonpositive. 
For every $M \in \dcat{C}(A, \cat{M})$ there is a distinguished triangle
\[ \opn{smt}^{\leq i}(M) \xar{e} M \xar{p} \opn{smt}^{\geq i + 1}(M) 
\xar{\th} \opn{T}(\opn{smt}^{\leq i}(M)) \]
in $\dcat{D}(A, \cat{M})$. Also, 
$\opn{H}^j(e) :  \opn{H}^j(\opn{smt}^{\leq i}(M))  \to \opn{H}^j(M)$
is an isomorphism in $\cat{M}$ for all $j \leq i$, and 
$\opn{H}^j(p) : \opn{H}^j(M) \to \opn{H}^j(\opn{smt}^{\geq i + 1}(M))$
is an isomorphism in $\cat{M}$ for all $j \geq i + 1$.
\end{prop}

\begin{proof}
The claims about $\opn{H}^j(e)$ and $\opn{H}^j(p)$ are trivial to verify. 
Now there is a short exact sequence 
\begin{equation} \label{eqn:3036}
0 \to \opn{smt}^{\leq i}(M) \xar{e} M \xar{p'} N \to 0 
\end{equation}
in $\dcat{C}_{\mrm{str}}(A, \cat{M})$, where 
\[ N := \bigl( \cdots \to 0 \to M^i / \opn{Z}^{i}(M) \xar{\d} M^{i+1} 
\xar{\d} M^{i+2} \to \cdots \bigr) . \]
According to Proposition \ref{prop:2165} we get a distinguished triangle 
\[ \opn{smt}^{\leq i}(M) \xar{e} M \xar{p'} N 
\xar{\th'} \opn{T}(\opn{smt}^{\leq i}(M)) \]
in $\dcat{D}(A, \cat{M})$. Next, there is an obvious quasi-isomorphism 
$\phi : N \to \opn{smt}^{\geq i + 1}(M)$
in $\dcat{C}_{\mrm{str}}(A, \cat{M})$
such that $p = \phi \circ p'$. We define the morphism 
\[ \th := \th' \circ \opn{Q}(\phi)^{-1} : \opn{smt}^{\geq i + 1}(M)
\to \opn{T}(\opn{smt}^{\leq i}(M)) \]
in $\dcat{D}(A, \cat{M})$. 
\end{proof}

\begin{prop} \label{prop:1463}
Assume $A$ is nonpositive.
For $\star \in \{ -, +, \mrm{b} \}$ the canonical functor
$\dcat{D}^{\star}(A, \cat{M}) \to \dcat{D}(A, \cat{M})^{\star}$ 
is an equivalence of triangulated categories. 
\end{prop}

\begin{proof} It is done in four steps. 

\smallskip \noindent
Step 1. Here we prove that the functor 
$F^- : \dcat{D}^-(A, \cat{M}) \to \dcat{D}(A, \cat{M})$ 
is fully faithful.
Let $s : M \to L$ be a quasi-isomorphism in 
$\dcat{K}(A, \cat{M})$ with 
$L \in \dcat{K}^-(A, \cat{M})$.
Say $L$ is concentrated in degrees $\leq i$.
Then $\opn{H}^j(M) = \opn{H}^j(L) = 0$ for all $j > i$. 
The smart truncation $\opn{smt}^{\leq i}(M)$ belongs to 
$\dcat{K}^-(A, \cat{M})$, and the inclusion
$t : \opn{smt}^{\leq i}(M) \to M$ is a quasi-isomorphism. 
According to Proposition \ref{prop:1320}, with 
$\cat{K} = \dcat{K}(A, \cat{M})$ and 
$\cat{K}' = \dcat{K}^-(A, \cat{M})$, and with condition (r), we see that 
$F^-$ is fully faithful.

\medskip \noindent 
Step 2. Here we prove that the functor 
$F^+ : \dcat{D}^+(\cat{M}) \to \dcat{D}(\cat{M})$ 
is fully faithful.
Let $s : L \to M$ be a quasi-isomorphism in 
$\dcat{K}(A, \cat{M})$ with $L \in \dcat{K}^+(A, \cat{M})$.
Say $L$ is concentrated in degrees $\geq i$. 
Then $\opn{H}^j(M) = \opn{H}^j(L) = 0$ for all $j < i$. 
The smart truncation 
$\opn{smt}^{\geq i}(M)$ belongs to $\dcat{K}^+(A, \cat{M})$, 
and the projection 
$t : M \to \opn{smt}^{\geq i}(M)$ is a quasi-isomorphism. 
According to Proposition \ref{prop:1320}, with condition (l), we see that 
$F^+$ is fully faithful.

\medskip \noindent 
Step 3. The arguments in step 1 we show that 
$\dcat{D}^{\mrm{b}}(A, \cat{M}) \to \dcat{D}^{+}(A, \cat{M})$
is fully faithful. And by step 2, 
$\dcat{D}^{+}(A, \cat{M}) \to \dcat{D}(A, \cat{M})$ is fully faithful. 
Therefore 
$\dcat{D}^{\mrm{b}}(A, \cat{M}) \to \dcat{D}(A, \cat{M})$
is fully faithful.

\medskip \noindent 
Step 4. Smart truncation shows that the functor
$\dcat{D}^{\star}(A, \cat{M}) \to \dcat{D}(A, \cat{M})^{\star}$ 
is essentially surjective on objects. 
\end{proof}

\begin{rem} \label{rem:1461}
\index{Boundedness! condition}
\index{1-D(M)@$\dcat{D}^{\star}(\cat{M})$}
\index{1-D(M)@$\dcat{D}(\cat{M})^{\star}$}
Definition \ref{dfn:1460} with $A = \K$, namely $\dcat{D}^{\star}(\cat{M})$,
is standard usage of this notation -- starting from \cite[Section I.4]{RD}, all 
the way to the recent book \cite[Definition 13.1.2]{KaSc2}. An exception is 
\cite[Definition tag=05RU]{SP}. 

On the other hand, definition \ref{dfn:1461} with $A = \K$, namely 
$\dcat{D}(\cat{M})^{\star}$, 
is almost never used (an exception is \cite{PSY3}). 

We are going to require both distinct concepts,  
$\dcat{D}^{\star}(A, \cat{M})$ and $\dcat{D}(A, \cat{M})^{\star}$,
until the end of Section \ref{sec:exist-resol}. Beginning with 
Section \ref{sec:adj-equ-cohdim}, in Definition \ref{dfn:2122}, the notation 
$\dcat{D}(A, \cat{M})^{\star}$ will be eliminated, and 
$\dcat{D}^{\star}(A, \cat{M})$ will be redefined to replace it. 
\end{rem}


% forcing page break here
\newpage
\mysubsection{Thick Subcategories of \texorpdfstring{$\cat{M}$}{M}}
\label{subsec:thick-subcats}

We begin by recalling a definition about abelian categories. 

\begin{dfn} \label{dfn:2323}
Let $\cat{M}$ be an abelian category. A 
{\em thick abelian subcategory}%
\index{Abelian category! thick abelian subcategory of}
of $\cat{M}$ is a full abelian subcategory $\cat{N}$ that is closed under
extensions. Namely if 
\[ 0 \to M' \to M \to M'' \to 0 \]
is a short exact sequence in $\cat{M}$ with $M', M'' \in \cat{N}$, then 
$M \in \cat{N}$ too. 
\end{dfn}

\begin{prop} \label{prop:3171} 
Let $\cat{M}$ be an abelian category, and let $\cat{M}' \sub \cat{M}$ be a 
thick abelian subcategory. Suppose 
$M_1 \to M_2 \to N \to M_3 \to M_4$ 
is an exact sequence in $\cat{M}$, and the objects $M_i$ belong $\cat{M}'$. 
Then 
$N \in \cat{M}'$ too. 
\end{prop}

\begin{exer} \label{exer:2881}
Prove this proposition.
\end{exer}

\begin{dfn} \label{dfn:2321}
Let $\cat{M}$ be an abelian category and $\cat{N} \sub \cat{M}$ a thick abelian 
subcategory. We denote by $\dcat{D}_{\cat{N}}(\cat{M})$ 
\index{1-DN(M)@$\dcat{D}_{\cat{N}}(\cat{M})$}
the full subcategory of $\dcat{D}(\cat{M})$
consisting of complexes $M$ such that 
$\opn{H}^i(M) \in \cat{N}$ for every $i$. 

Given a boundedness condition $\star$, 
\index{Boundedness! condition}
\index{1-DNstar(M)@$\dcat{D}^{\star}_{\cat{N}}(\cat{M})$}
we write 
$\dcat{D}_{\cat{N}}^{\star}(\cat{M}) := 
\dcat{D}_{\cat{N}}(\cat{M}) \cap \dcat{D}^{\star}(\cat{M})$
and
$\dcat{D}_{\cat{N}}(\cat{M})^{\star} := 
\dcat{D}_{\cat{N}}(\cat{M}) \cap \dcat{D}(\cat{M})^{\star}$.
\end{dfn}

\begin{prop} \label{prop:3052}
If $\cat{N}$ is a thick abelian subcategory of $\cat{M}$ then
$\dcat{D}_{\cat{N}}(\cat{M})$ is a full triangulated subcategory of 
$\dcat{D}(\cat{M})$.
\end{prop}

\begin{proof}
Clearly $\dcat{D}_{\cat{N}}(\cat{M})$ is closed under translations. 
Now suppose 
$M' \to M \to M'' \xar{\, \triangle\, }$ \, 
is a distinguished triangle in $\dcat{D}(\cat{M})$ such that 
$M', M \in \dcat{D}_{\cat{N}}(\cat{M})$; we have to show that $M''$
is also in $\dcat{D}_{\cat{N}}(\cat{M})$.
Consider the  exact sequence 
$\opn{H}^i(M') \to \opn{H}^i(M) \to
\opn{H}^i(M'') \to \opn{H}^{i+1}(M') \to \opn{H}^{i+1}(M)$.
The four outer objects belong to $\cat{N}$. 
Since $\cat{N}$ is a thick abelian subcategory of $\cat{M}$ it follows (using 
Proposition \ref{prop:3171}) that 
$\opn{H}^i(M'') \in \cat{N}$.
\end{proof}


\begin{exa}
Let $A$ be a noetherian commutative ring. 
The category $\cat{Mod}_{\mrm{f}} A$ of finitely generated
modules is a thick abelian subcategory of $\cat{Mod} A$.
\end{exa}

\begin{exa}
Consider $\cat{Mod} \Z = \cat{Ab}$. 
As above we have the thick abelian subcategory 
$\cat{Ab}_{\mrm{f}} = \cat{Mod}_{\mrm{f}} \Z$ of finitely generated
abelian groups.
There is also the thick abelian subcategory 
$\cat{Ab}_{\mrm{tors}}$ of torsion abelian groups (every element has a finite
order). The intersection of $\cat{Ab}_{\mrm{tors}}$ and 
$\cat{Ab}_{\mrm{f}}$ is the category $\cat{Ab}_{\mrm{fin}}$ of finite abelian
groups. This is also thick. 
\end{exa}

\begin{exa}
Let $X$ be a noetherian scheme (e.g.\ an algebraic variety over an 
algebraically closed field). Consider the abelian category 
$\cat{Mod} \mcal{O}_X$ of $\mcal{O}_X$-modules. In it there is the thick abelian
subcategory $\cat{QCoh} \mcal{O}_X$ of quasi-coherent sheaves, and in 
that there is the thick abelian
subcategory $\cat{Coh} \mcal{O}_X$ of coherent sheaves.

Passing to complexes, let us write 
$\dcat{C}(X) := \dcat{C}(\cat{Mod} \mcal{O}_X)$,
the DG category of unbounded complexes. Its strict subcategory is 
$\dcat{C}_{\cat{str}}(X)$, and for each boundedness condition $\star$ we have 
the full subcategory 
$\dcat{C}^{\star}_{\cat{str}}(X) \sub \dcat{C}_{\cat{str}}(X)$. 
For the derived category we write 
$\dcat{D}(X) := \dcat{D}(\cat{Mod} \mcal{O}_X)$.
Inside it there are the full triangulated subcategories
\[ \dcat{D}_{\mrm{c}}(X) := 
\dcat{D}_{\cat{Coh} \mcal{O}_X}(\cat{Mod} \mcal{O}_X) \sub
\dcat{D}_{\mrm{qc}}(X) := 
\dcat{D}_{\cat{QCoh} \mcal{O}_X}(\cat{Mod} \mcal{O}_X) . \]

One can show (see \cite[Corollary II.7.19]{RD}) that the canonical functor \lb
$\dcat{D}^{+}(\cat{QCoh} \mcal{O}_X) \to \dcat{D}^{+}_{\mrm{qc}}(X)$
is an equivalence. The proof is based on the deep fact 
(see \cite[Theorem II.7.18]{RD}) 
that $\cat{QCoh} \mcal{O}_X$ has enough injectives relative to 
$\cat{Mod} \mcal{O}_X$, in the sense of Definition \ref{dfn:2880}. 
Then Theorem \ref{thm:2880} says that every complex 
$\MM \in \dcat{D}^{+}_{\mrm{qc}}(X)$ admits a quasi-isomorphism 
$\MM \to \II$ in $\dcat{C}_{\cat{str}}(X)$, where $\II$ is a bounded below 
complex of injective quasi-coherent $\OO_X$-modules. 

Then, using smart truncations and the fact that every 
quasi-coherent sheaf on $X$ is the direct limit of its coherent subsheaves, one 
can show that the canonical functor 
$\dcat{D}^{\mrm{b}}(\cat{Coh} \mcal{O}_X) \to \dcat{D}^{\mrm{b}}_{\mrm{c}}(X)$
is an equivalence. See \cite[Proposition 3.5]{Huy}; when $X = 
\opn{Spec}(A)$ is affine see also Proposition \ref{prop:200} below. The 
category 
$\dcat{D}^{\mrm{b}}_{\mrm{c}}(X)$
is the focus of a significant part of the research in modern algebraic geometry 
(e.g.\ birational geometry, see \cite{Huy}). 

A totally different approach yields this result: if $X$ is separated and 
quasi-compact, then the canonical functor
$\dcat{D}^{}(\cat{QCoh} \mcal{O}_X) \to \dcat{D}^{}_{\mrm{qc}}(X)$
is an equivalence. See \cite[Corollary 5.5]{BoNe}.
\end{exa}

For a left noetherian ring $A$ we write 
$\dcat{D}_{\mrm{f}}(\cat{Mod} A) :=
\dcat{D}_{\cat{Mod}_{\mrm{f}} A}(\cat{Mod} A)$.

\begin{prop} \label{prop:200}
Let $A$ be a left noetherian ring and $\star \in \{ -, \mrm{b} \}$. Then the
canonical functor 
$\dcat{D}^{\star}(\cat{Mod}_{\mrm{f}} A) \to 
\dcat{D}_{\mrm{f}}(\cat{Mod} A)^{\star}$
is an equivalence of triangulated categories.
\end{prop}

\begin{proof}
Consider the functor 
$F : \dcat{D}^{-}(\cat{Mod}_{\mrm{f}} A) \to \dcat{D}(\cat{Mod} A)$.
Suppose $s : M \to L$ is a quasi-isomorphism in $\dcat{K}(\cat{Mod} A)$,
such that $L \in \dcat{K}^{-}(\cat{Mod}_{\mrm{f}} A)$. 
Then $M \in \dcat{D}_{\mrm{f}}(\cat{Mod} A)^{-}$.
A bit later (in Theorem \ref{thm:3340}) we will prove that $M$ 
admits a free resolution $P \to M$, where $P$ is a
bounded above complex of finitely generated free modules. 
Thus we get a quasi-isomorphism $t : P \to M$ with 
$P \in \dcat{K}^{-}(\cat{Mod}_{\mrm{f}} A)$.  By Proposition \ref{prop:1320} 
with condition (r) we conclude that $F$ is fully faithful. This also shows 
that the essential image of $F$ is 
$\dcat{D}_{\mrm{f}}(\cat{Mod} A)^{-}$.
 
Next consider the functor 
$G : \dcat{D}^{\mrm{b}}(\cat{Mod}_{\mrm{f}} A) \to 
\dcat{D}^{-}(\cat{Mod}_{\mrm{f}} A)$.
Suppose $s : L \to M$ is a quasi-isomorphism in 
$\dcat{K}^-(\cat{Mod}_{\mrm{f}} A)$
with $L \in \dcat{K}^{\mrm{b}}(\cat{Mod}_{\mrm{f}} A)$. 
Say $\opn{H}(L)$ is concentrated in the integer interval $[d_0, d_1]$. 
Then 
$t : M \to \opn{smt}^{\geq d_0}(M)$
is a quasi-isomorphism, and 
$\opn{smt}^{\geq d_0}(M) \in \dcat{K}^{\mrm{b}}(\cat{Mod}_{\mrm{f}} A)$. 
By Proposition \ref{prop:1320} with condition (l) we conclude that $G$ is 
fully faithful. Therefore the composition 
$F \circ G : \dcat{D}^{\mrm{b}}(\cat{Mod}_{\mrm{f}} A) \to 
\dcat{D}(\cat{Mod} A)$
is fully faithful. Suitable truncations ($\opn{smt}^{\geq d_0}$ and 
$\opn{smt}^{\leq d_1}$) show that the essential image of $F \circ G$ is 
$\dcat{D}_{\mrm{f}}(\cat{Mod} A)^{\mrm{b}}$.
\end{proof}




\mysubsection{The Embedding of \texorpdfstring{$\cat{M}$}{M} in 
\texorpdfstring{$\dcat{D}(\cat{M})$}{D(M)}}

Here again we only consider an abelian category $\cat{M}$. 

For $M, N \in \cat{M}$ there is no difference between the $\K$-modules
$\opn{Hom}_{\cat{M}}(M, N)$, $\opn{Hom}_{\dcat{C}(\cat{M})}(M, N)$ and 
$\opn{Hom}_{\dcat{K}(\cat{M})}(M, N)$. Thus the canonical functors 
$\cat{M} \to \dcat{C}(\cat{M})$ and 
$\cat{M} \to \dcat{K}(\cat{M})$ are fully faithful. 
The same is true for $\dcat{D}(\cat{M})$, but this requires a proof.

Let $\dcat{D}(\cat{M})^0$ be the full subcategory of 
$\dcat{D}(\cat{M})$ consisting of complexes whose cohomology is 
concentrated in degree $0$. This is an additive subcategory of 
$\dcat{D}(\cat{M})$.

\begin{prop} \label{prop:121}
The canonical functor $\cat{M} \to \dcat{D}(\cat{M})^0$ 
is an equivalence.
\end{prop}

\begin{proof}
Let's denote the canonical functor $\cat{M} \to \dcat{D}(\cat{M})^0$ 
by $F$. Under the fully faithful embedding 
$\cat{M} \subseteq \dcat{C}_{\mrm{str}}(\cat{M})$,
$F$ is just the restriction of the functor $\til{\opn{Q}}$ from Definition 
\ref{dfn:1480}.

The functor $\opn{H}^0 : \dcat{D}(\cat{M}) \to \cat{M}$
satisfies $\opn{H}^0 \circ \, F = \opn{Id}_{\cat{M}}$. This implies 
that $F$ is faithful.

Next we prove that $F$ is full. Take any objects $M, N \in \cat{M}$ and a 
morphism $q : M \to N$ in $\dcat{D}(\cat{M})$. 
By Proposition \ref{prop:1481} we know that  
$q = \til{\opn{Q}}(a) \circ \til{\opn{Q}}(s)^{-1}$ for some
morphisms $a : L \to N$ and $s : L \to M$ in 
$\dcat{C}_{\mrm{str}}(\cat{M})$, with $s$
a quasi-isomorphism. Let $L' := \opn{smt}^{\leq 0}(L)$, the smart truncation of 
$L$. Since $L \in \dcat{D}(\cat{M})^0$, the inclusion 
$u : L' \to L$ is a quasi-isomorphism  in $\dcat{C}_{\mrm{str}}(\cat{M})$. 
Writing $a' := a \circ u$ and $s' := s \circ u$, we see that $s'$ is a 
quasi-isomorphism, and
$q = \til{\opn{Q}}(a') \circ \til{\opn{Q}}(s')^{-1}$. 

Next let $L'' := \opn{smt}^{\geq 0}(L')$, the other smart truncation. 
The projection $v : L' \to L''$ is a surjective quasi-isomorphism in 
$\dcat{C}_{\mrm{str}}(\cat{M})$. 
Because $L''$ is a complex concentrated in degree $0$, we can view it as 
an object of $\cat{M}$. The morphisms $a'$ and $s'$ factor 
as $a' = a'' \circ v$ and $s' = s'' \circ v$, where 
$a'' : L'' \to N$ and $s'' : L'' \to M$ are morphisms 
in $\cat{M}$. But $s''$ is a quasi-isomorphism in 
$\dcat{C}_{\mrm{str}}(\cat{M})$, and so it is actually an isomorphism in 
$\cat{M}$. Therefore we have a morphism 
$a'' \circ (s'')^{-1} : M \to N$
in $\cat{M}$, and 
\[ \til{\opn{Q}}(a'' \circ (s'')^{-1}) = 
\til{\opn{Q}}(a'') \circ \til{\opn{Q}}(s'')^{-1} = 
\til{\opn{Q}}(a') \circ \til{\opn{Q}}(s')^{-1} =  q  . \]

Finally we have to prove that every $L \in \dcat{D}(\cat{M})^0$ is isomorphic, 
in $\dcat{D}(\cat{M})$, to a complex $L''$ that's concentrated in degree $0$. 
But we already showed it in the previous paragraphs.
\end{proof}

\begin{prop} \label{prop:1440}
Let $\cat{M}$ be an abelian category and let 
$0 \to L \xar{\phi}  M \xar{\psi} N \to 0$ 
be a diagram in $\cat{M}$. The following conditions are equivalent\tup{:}
\begin{enumerate}
\rmitem{i} The diagram is an exact sequence. 
\rmitem{ii}  There is a distinguished triangle 
$L \xar{ \, \til{\opn{Q}}(\phi) \, }  M \xar{ \, \til{\opn{Q}}(\psi) \, } N 
\xar{ \, \th \, } \opn{T}(L)$ 
in $\dcat{D}(\cat{M})$.
\end{enumerate}
\end{prop}

\begin{exer} \label{exer:1917}
Prove Proposition \ref{prop:1440}. 
\end{exer}

The last two propositions say that the abelian category $\cat{M}$, with its 
kernels and cokernels, can be recovered from the triangulated category 
$\dcat{D}(\cat{M})$. 

\begin{rem} \label{rem:2320}
Assume that the diagram in Proposition \ref{prop:1440} is an exact sequence.
By Proposition \ref{prop:3605}, this sequence is split if and only if the 
morphism $\th$ is zero. Furthermore, if $\cat{M}$ has enough injectives, then 
there is a canonical bijection between the set of 
isomorphism classes of extensions of $N$ by $L$, and the set 
$\opn{Hom}_{\dcat{D}(\cat{M})}(N, \opn{T}(L))$. See \cite[Section I.6]{RD}, or 
use Exercise \ref{exer:2125}, which says that 
$\opn{Hom}_{\dcat{D}(\cat{M})}(N, \opn{T}(L)) \cong 
\opn{Ext}^1_{\cat{M}}(N, L)$.
\end{rem}



\mysubsection{The Opposite Derived Category is Triangulated}
\label{subsec:opp-dercat-triang}

Here we deal with a DG $\K$-ring $A$ and a $\K$-linear 
abelian category $\cat{M}$. 
In Subsection \ref{subsec:opp-hom-triang} we put a canonical triangulated 
structure on the opposite homotopy category 
$\dcat{K}(A, \cat{M})^{\mrm{op}}$. This structure is such that the flip functor 
\[ \opn{Flip} : \dcat{K}(A, \cat{M})^{\mrm{op}} \to
\dcat{K}(A^{\mrm{op}}, \cat{M}^{\mrm{op}}) =  
\dcat{K}(A, \cat{M})^{\mrm{flip}} \]
is an isomorphism of triangulated categories. 
Here, unlike in Subsection \ref{subsec:opp-hom-triang},
we are omitting the overline decoration from $\opn{Flip}$. 
In the current subsection we push this triangulated structure to 
the opposite derived category $\dcat{D}(A, \cat{M})^{\mrm{op}}$.

But first let us present two types of full subcategories
$\cat{K} \sub \dcat{K}(A, \cat{M})$ that are triangulated, and also 
$\cat{K}^{\mrm{op}} \sub \dcat{K}(A, \cat{M})^{\mrm{op}}$ 
is triangulated. 

For any boundedness condition $\star$ (see Definition \ref{dfn:3221}), we know 
that $\dcat{K}^{\star}(A, \cat{M})$
is a full triangulated subcategory of $\dcat{K}^{}(A, \cat{M})$. 
The notation 
$\dcat{K}^{\star}(A, \cat{M})^{\mrm{op}}$
refers to the opposite category of 
$\dcat{K}^{\star}(A, \cat{M})$;
thus, $\dcat{K}^{\star}(A, \cat{M})^{\mrm{op}}$ is the full subcategory of 
$\dcat{K}^{}(A, \cat{M})^{\mrm{op}}$
on the DG modules satisfying the boundedness condition $\star$. 

\begin{prop} \label{prop:3050}
For every boundedness condition $\star$, the subcategory 
$\dcat{K}^{\star}(A, \cat{M})^{\mrm{op}}$
is triangulated in $\dcat{K}^{}(A, \cat{M})^{\mrm{op}}$. 
\end{prop}

\begin{proof}
Let's write $\cat{K} := \dcat{K}^{\star}(A, \cat{M})$ and
$\cat{K}^{\mrm{op}} := \dcat{K}^{\star}(A, \cat{M})^{\mrm{op}}$.
Define  
$\cat{K}^{\mrm{flip}} \lb := \opn{Flip}(\cat{K}^{\mrm{op}}) \sub 
\dcat{K}(A^{\mrm{op}}, \cat{M}^{\mrm{op}})$.
According to Theorem \ref{thm:2495}(2) we know that
$\cat{K}^{\mrm{flip}} = \dcat{K}^{-\star}(A^{\mrm{op}}, \cat{M}^{\mrm{op}})$,
where $-\star$ is the reversed boundedness condition. 
Thus $\cat{K}^{\mrm{flip}}$ is a full triangulated subcategory of
$\dcat{K}(A^{\mrm{op}}, \cat{M}^{\mrm{op}})$. 
By definition of the triangulated structure on 
$\dcat{K}^{}(A, \cat{M})^{\mrm{op}}$
we conclude that $\cat{K}^{\mrm{op}}$ is triangulated.
\end{proof}

Let $\cat{N} \sub \cat{M}$ be a thick abelian subcategory. Then 
$\dcat{K}_{\cat{N}}(\cat{M})$, see Definition \ref{dfn:2321},  
is a full triangulated subcategory of $\dcat{K}(\cat{M})$.
This is just like Proposition \ref{prop:3052}.
The opposite category 
$\dcat{K}_{\cat{N}}(\cat{M})^{\mrm{op}}$,
on the same set of objects, is a full additive subcategory of 
$\dcat{K}(\cat{M})^{\mrm{op}}$. 

\begin{prop} \label{prop:3051}
Let $\cat{N} \sub \cat{M}$ be a thick abelian subcategory.
Then 
$\dcat{K}_{\cat{N}}(\cat{M})^{\mrm{op}}$
is a full triangulated subcategory of $\dcat{K}(\cat{M})^{\mrm{op}}$. 
\end{prop}

\begin{proof}
Let's write 
$\cat{K} := \dcat{K}_{\cat{N}}(\cat{M})$. Then 
$\cat{K}^{\mrm{flip}} := \opn{Flip}(\cat{K}^{\mrm{op}}) \sub 
\dcat{K}(\cat{M}^{\mrm{op}})$
is, by Theorem \ref{thm:2495}(4), the category  
$\dcat{K}_{\cat{N}^{\mrm{op}}}(\cat{M}^{\mrm{op}})$. 
Thus $\cat{K}^{\mrm{flip}}$ is a triangulated subcategory. 
By definition of the triangulated structure on 
$\dcat{K}^{}(\cat{M})^{\mrm{op}}$
we conclude that $\cat{K}^{\mrm{op}}$ is triangulated.
\end{proof}

Later, in Section \ref{sec:resol},  we will see that several {\em resolving 
subcategories} are also full triangulated subcategories of 
$\dcat{K}(A, \cat{M})^{\mrm{op}}$.
These are relevant to Theorems \ref{thm:1470} and \ref{thm:1500} below. 

Taking intersections, we get more full triangulated subcategories of  
$\dcat{K}^{}(A, \cat{M})^{\mrm{op}}$.

Recall that 
$\dcat{S}(A, \cat{M})$ is the set of quasi-isomorphisms in 
$\dcat{K}(A, \cat{M})$. Let $\dcat{S}(A, \cat{M})^{\mrm{op}}$ be the 
set of quasi-isomorphisms in $\dcat{K}(A, \cat{M})^{\mrm{op}}$.
Note that being a quasi-isomorphism has nothing to do with the 
triangulated structure. Since a morphism $\psi$ in 
$\dcat{K}(A, \cat{M})$ is a quasi-isomorphism iff 
$\opn{Op}(\psi)$ is a quasi-isomorphism in 
$\dcat{K}(A, \cat{M})^{\mrm{op}}$, we have a bijection
$\opn{Op} : \dcat{S}(A, \cat{M}) \iso \dcat{S}(A, \cat{M})^{\mrm{op}}$.
We know that $\dcat{S}(A, \cat{M})$ is a left and right 
denominator set in 
$\dcat{K}(A, \cat{M})$, by Proposition \ref{prop:106}. 
Therefore, by Proposition \ref{prop:1845},
$\dcat{S}(A, \cat{M})^{\mrm{op}}$ is left and right denominator set in 
$\dcat{K}(A, \cat{M})^{\mrm{op}}$.
We need to know more: 

\begin{prop} \label{prop:2356}
Let $\cat{K}^{\mrm{op}}$ be a full triangulated subcategory of 
$\dcat{K}(A, \cat{M})^{\mrm{op}}$, and define
$\cat{S}^{\mrm{op}} := \cat{K}^{\mrm{op}} \cap \, 
\dcat{S}(A, \cat{M})^{\mrm{op}}$,
the set of quasi-isomorphisms in $\cat{K}^{\mrm{op}}$. 
Then\tup{:}
\begin{enumerate}
\item $\cat{S}^{\mrm{op}}$ is a denominator set of cohomological origin in 
$\cat{K}^{\mrm{op}}$.

\item The localized category 
$\cat{D}^{\mrm{op}} := (\cat{K}^{\mrm{op}})_{\cat{S}^{\mrm{op}}}$
has a unique triangulated structure s.t.\ the localization functor 
$\opn{Q}^{\mrm{op}} : \cat{K}^{\mrm{op}} \to \cat{D}^{\mrm{op}}$
is a triangulated functor. 
\end{enumerate}
\end{prop}

\begin{proof}
On the flip side, i.e.\ in 
$\dcat{K}(A^{\mrm{op}}, \cat{M}^{\mrm{op}})$,
we know that the set of quasi-iso\-morphisms
$\dcat{S}(A^{\mrm{op}}, \cat{M}^{\mrm{op}})$
is a denominator set of cohomological origin. But
by Theorem \ref{thm:2495}(4) we have 
$\dcat{S}(A, \cat{M})^{\mrm{op}} = 
\opn{Flip}^{-1} \bigl( \dcat{S}(A^{\mrm{op}}, \cat{M}^{\mrm{op}}) \bigr)$,
so this too is a denominator set of cohomological origin. 
Now we can use Theorem \ref{thm:105}. 
\end{proof}


Observe that 
$\cat{D}^{\mrm{op}} = (\cat{K}_{\cat{S}})^{\mrm{op}}$,
by Proposition \ref{prop:3030}. 
The situation is summarized by the following commutative diagram of functors:
\begin{equation} \label{eqn:3060}
\UseTips \xymatrix @C=10ex @R=6ex {
\cat{C}_{\mrm{str}}^{\mrm{op}}
\ar@{->>}[r]^{\opn{P}^{\mrm{op}}}
\ar[d]_{\opn{Flip}}^{\cong}
&
\cat{K}^{\mrm{op}}
\ar[r]^{\opn{Q}^{\mrm{op}}}
\ar[d]_{\opn{Flip}}^{\cong}
&
\cat{D}^{\mrm{op}}
\ar[d]_{\opn{Flip}}^{\cong}
\\
\cat{C}_{\mrm{str}}^{\mrm{flip}}
\ar@{->>}[r]^{\opn{P}^{\mrm{flip}}}
&
\cat{K}^{\mrm{flip}}
\ar[r]^{\opn{Q}^{\mrm{flip}}}
&
\cat{D}^{\mrm{flip}} 
} 
\end{equation}
Here $\cat{C}$ is the full subcategory of 
$\dcat{C}(A, \cat{M})$ on the objects of $\cat{K}$,
and 
$\cat{C}^{\mrm{flip}} := \opn{Flip}(\cat{C}) \sub 
\dcat{C}(A, \cat{M})^{\mrm{flip}} = 
\dcat{C}(A^{\mrm{op}}, \cat{M}^{\mrm{op}})$.
All these functors are bijective on objects. The vertical ones are bijective on 
morphisms too. The functors marked $\surj$ are surjective on morphisms (i.e.\ 
they are full). The first vertical arrow is an isomorphism of abelian 
categories, and the other two vertical arrows are isomorphisms of triangulated 
categories.

Warning: Given $\cat{K}^{\mrm{op}}$ like in Proposition \ref{prop:2356}, 
there is no reason for the canonical triangulated functor 
$\cat{D}^{\mrm{op}} \to \dcat{D}(A, \cat{M})^{\mrm{op}}$
to be fully faithful; see Proposition \ref{prop:1320} for sufficient 
conditions. 


\begin{dfn} \label{dfn:3225}
Let $\cat{K} \sub \dcat{K}(A, \cat{M})$ be a full additive subcategory s.t.\ 
$\cat{K}^{\mrm{op}}$ is a  triangulated subcategory of 
$\dcat{K}(A, \cat{M})^{\mrm{op}}$, and let 
$\cat{S} := \cat{K} \cap \, \dcat{S}(A, \cat{M})$.
The category 
$\cat{D}^{\mrm{op}} := (\cat{K}^{\mrm{op}})_{\cat{S}^{\mrm{op}}}$
is given the triangulated structure from Proposition \ref{prop:2356}. 
\index{Derived category! opposite}
\end{dfn}

For $\cat{K} = \dcat{K}(A, \cat{M})$ we get a triangulated structure on 
$\cat{D}^{\mrm{op}} = \dcat{D}(A, \cat{M})^{\mrm{op}}$.

\begin{dfn} \label{dfn:3226} 
Let $\cat{K} \sub \dcat{K}(A, \cat{M})$ be a full additive subcategory s.t.\ 
$\cat{K}^{\mrm{op}}$ is a triangulated subcategory of 
$\dcat{K}(A, \cat{M})^{\mrm{op}}$.
Define
$\cat{D} := \cat{K}_{\cat{S}}$
and 
$\cat{D}^{\mrm{op}} := (\cat{K}^{\mrm{op}})_{\cat{S}^{\mrm{op}}}$, 
where $\cat{S} := \cat{K} \cap \, \dcat{S}(A, \cat{M})$.
Let $\cat{E}$ be some triangulated category. A 
{\em contravariant triangulated functor}%
\index{Triangulated functor! contravariant}
from $\cat{D}$ to $\cat{E}$ is, by definition, a 
triangulate functor 
$F : \cat{D}^{\mrm{op}}  \to \cat{E}$,
where $\cat{D}^{\mrm{op}}$ has the canonical triangulated structure from 
Definition \ref{dfn:3225}.
\end{dfn}







\cleardoublepage
\mysection{Derived Functors} \label{sec:der-funcs}

\AYcopyright

Suppose 
$F : \dcat{K}(A, \cat{M}) \to \cat{E}$
is a triangulated functor. Here $A$ is a DG ring, $\cat{M}$ is an abelian 
category, and $\cat{E}$ is a triangulated category. 
In this section we define the right and left derived functors
$\mrm{R} F, \, \mrm{L} F : \dcat{D}(A, \cat{M}) \to \cat{E}$ 
of $F$. These are also triangulated functors, satisfying certain universal 
properties. We shall prove the uniqueness of the derived functors, and their 
existence under suitable assumptions.
 
The universal properties of the derived functors are best stated in 
$2$-categorical language. This will be explained is the first two subsections. 
In Subsection \ref{subsec:abstr-der-funs} we define derived functors in the 
abstract setting (as opposed to the triangulated setup), and prove the main 
results for them. These results will then be specialized to various settings:
triangulated functors (Subsection \ref{subsec:tri-der-funs}), contravariant 
triangulated functors (Subsection \ref{subsec:cntr-tri-der-funs}), and 
triangulated bifunctors (Subsection \ref{subsec:tri-bifun}). 



\mysubsection{\texorpdfstring{$2$}{2}-Categorical Notation} 
\label{subsec:2-cat-not}

In this section we are going to do a lot of work with morphisms of functors 
(i.e.\ natural transformations). The language and notation of ordinary category 
theory that we used so far is not adequate for this purpose. Therefore we will 
now introduce notation from the theory of {\em $2$-categories}.
(We will not give a definition of a $2$-category here; but it is basically the 
structure of $\cat{Cat}$ that is mentioned below.) 
For more details on $2$-categories the reader 
can look at \cite{Mac2} or \cite[Section 1]{Ye7}.

Consider the set $\cat{Cat}$ of all categories. The set-theoretical aspects 
are neglected, as explained in Subsection \ref{subsec:set-theor}.
(Briefly, the precise solution is this: $\cat{Cat}$ is the set of all 
$\cat{U}$-categories; so $\cat{Cat}$ is a subset of a bigger Grothendieck 
universe, say $\cat{V}$, and it is a $\cat{V}$-category.)

The set $\cat{Cat}$ is the set of objects of a $2$-category.
\index{Two@$2$-category} 
This means that in $\cat{Cat}$ there are two kinds of morphisms: 
{\em $1$-morphisms} between objects, and {\em $2$-morphisms} between 
$1$-morphisms. There are several kinds of compositions, and these have several 
properties. All this will be explained below. 

Suppose $\cat{C}_0$, $\cat{C}_1, \ldots$ are categories, namely 
objects of $\cat{Cat}$. The $1$-morph\-isms between them are the functors.
The notation is as usual: $F : \cat{C}_0 \to \cat{C}_1$ denotes a functor.

Suppose $F, G : \cat{C}_0 \to \cat{C}_1$ are functors. The $2$-morphisms from 
$F$ to $G$ are the morphisms of functors 
(i.e.\ the natural transformations), and the notation is 
$\eta : F \twoto G$. The double arrow is the distinguishing notation for  
$2$-morphisms. When specializing to an object
$M \in \cat{C}_0$ we revert to the single arrow notation, namely 
$\eta_M : F(M) \to G(M)$ is the corresponding morphism in $\cat{C}_1$. 
The diagram depicting this is
\[ \UseTips  \xymatrix @C=12ex @R=6ex {
\cat{C}_0
\ar@(ur,ul)[r]^(0.5){F} _(0.5){}="eta-s"
\ar@(dr,dl)[r]_(0.5){G}  ^(0.5){}="eta-t"
\ar@{=>}  "eta-s";"eta-t" _(0.5){\eta}
&
\cat{C}_1
} \]
We shall refer to such a diagram as a {\em $2$-diagram}. 
\index{Two@$2$-diagram} 

Each object (category) $\cat{C}$ has its identity $1$-morphism  (functor)
$\opn{Id}_{\cat{C}} : \cat{C} \to \cat{C}$.  
Each $1$-morphism $F$ has its identity $2$-morphism (natural transformation)
$\opn{id}_F : F \twoto F$.  

Now we consider compositions. For functors there is nothing new: given functors 
$F_1 : \cat{C}_{0} \to \cat{C}_{1}$ and 
$F_2 : \cat{C}_{1} \to \cat{C}_{2}$, their composition, that we now call 
{\em horizontal composition}, is the functor
$F_2 \circ F_1 : \cat{C}_{0} \to \cat{C}_{2}$. The diagram is 
\[ \UseTips  \xymatrix @C=12ex @R=6ex {
\cat{C}_0
\ar[r]^(0.5){F_1} 
\ar@(dr,dl)[rr]_(0.5){F_2 \, \circ \, F_1}  
&
\cat{C}_1
\ar[r]^(0.5){F_2} 
&
\cat{C}_2
} \]
This can be viewed as a commutative $1$-diagram, or as a shorthand for the 
$2$-diagram
\[ \UseTips  \xymatrix @C=12ex @R=6ex {
\cat{C}_0
\ar[r]^(0.5){F_1}
\ar @/_3em/ [rr]_(0.5){F_2 \, \circ \, F_1} _(0.5){}="t"
&
\cat{C}_1
\ar[r]^(0.5){F_2} 
\ar@{=>}  ; "t" _(0.5){\opn{id}}
&
\cat{C}_2
} \]
in which $\opn{id}$ is the identity $2$-morphism of $F_2 \circ F_1$. 

The complication begins with compositions of $2$-morphisms. 
Suppose we are given $1$-morphisms
$F_i, G_i : \cat{C}_{i - 1} \to \cat{C}_{i}$
and $2$-morphisms $\eta_i : F_i \twoto G_i$. In a diagram: 
\[ \UseTips  \xymatrix @C=12ex @R=6ex {
\cat{C}_0
\ar@(ur,ul)[r]^(0.5){F_1} _(0.5){}="eta1-s"
\ar@(dr,dl)[r]_(0.5){G_1}  ^(0.5){}="eta1-t"
\ar@{=>}  "eta1-s";"eta1-t" _(0.5){\eta_1}
&
\cat{C}_1
\ar@(ur,ul)[r]^(0.5){F_2} _(0.5){}="eta2-s"
\ar@(dr,dl)[r]_(0.5){G_2}  ^(0.5){}="eta2-t"
\ar@{=>}  "eta2-s";"eta2-t" _(0.5){\eta_2}
&
\cat{C}_2
} \]
The {\em horizontal composition} is the morphism of functors 
$\eta_2 \circ \eta_1 : F_2 \circ F_1  \twoto G_2 \circ G_1$.
The diagram is 
\[ \UseTips  \xymatrix @C=18ex @R=6ex {
\cat{C}_0
\ar@(ur,ul)[r]^(0.5){F_2 \, \circ \, F_1} _(0.5){}="eta-s"
\ar@(dr,dl)[r]_(0.5){G_2 \, \circ \, G_1}  ^(0.5){}="eta-t"
\ar@{=>}  "eta-s";"eta-t" _(0.5){\eta_2 \, \circ \, \eta_1}
&
\cat{C}_2
} \] 

\begin{exer} \label{exer:1314}
For an object $M \in \cat{C}_0$, give an explicit formula for the morphism 
\[ (\eta_2 \circ \eta_1)_M : (F_2 \circ F_1)(M) \to (G_2 \circ G_1)(M) \]
in the category $\cat{C}_2$.
\end{exer}

Suppose we are given $1$-morphisms 
$E, F, G : \cat{C}_{0} \to \cat{C}_{1}$, and $2$-morphisms \lb 
$\ze : E \twoto F$ and $\eta : F \twoto G$. 
The diagram depicting this is
\[ \UseTips  \xymatrix @C=16ex @R=6ex {
\cat{C}_0
\ar@(u,u)[r]^(0.5){E} _(0.5){}="zeta-s"
\ar[r]_(0.7){F}  ^(0.5){}="zeta-t" _(0.5){}="eta-s"
\ar@(d,d)[r]_(0.5){G}  ^(0.5){}="eta-t"
\ar@{=>}  "zeta-s";"zeta-t" _(0.5){\zeta}
\ar@{=>}  "eta-s";"eta-t" _(0.5){\eta}
&
\cat{C}_1
} \] 
The {\em vertical composition} of $\zeta$ and $\eta$ is the $2$-morphism  
$\eta * \zeta : E \to G$.
Notice the new symbol for this operation. 
The corresponding diagram is 
\[ \UseTips  \xymatrix @C=18ex @R=6ex {
\cat{C}_0
\ar@(ur,ul)[r]^(0.5){E} _(0.5){}="eta-s"
\ar@(dr,dl)[r]_(0.5){G}  ^(0.5){}="eta-t"
\ar@{=>}  "eta-s";"eta-t" _(0.5){\eta \, * \, \ze}
&
\cat{C}_1
} \] 

\begin{exer} \label{exer:1311}
For an object $M \in \cat{C}_0$, give an explicit formula for the morphism
$(\eta * \zeta)_M : E(M) \to G(M)$
in the category $\cat{C}_{1}$.
\end{exer}

Something intricate occurs in the situation shown in the next diagram. 
\[ \UseTips  \xymatrix @C=16ex @R=6ex {
\cat{C}_0
\ar@(u,u)[r]^(0.5){E_1} _(0.5){}="zeta1-s"
\ar[r]_(0.7){F_1}  ^(0.5){}="zeta1-t" _(0.5){}="eta1-s"
\ar@(d,d)[r]_(0.5){G_1}  ^(0.5){}="eta1-t"
\ar@{=>}  "zeta1-s";"zeta1-t" _(0.5){\zeta_1}
\ar@{=>}  "eta1-s";"eta1-t" _(0.5){\eta_1}
&
\cat{C}_1
\ar@(u,u)[r]^(0.5){E_2} _(0.5){}="zeta2-s"
\ar[r]_(0.7){F_2}  ^(0.5){}="zeta2-t" _(0.5){}="eta2-s"
\ar@(d,d)[r]_(0.5){G_2}  ^(0.5){}="eta2-t"
\ar@{=>}  "zeta2-s";"zeta2-t" _(0.5){\zeta_2}
\ar@{=>}  "eta2-s";"eta2-t" _(0.5){\eta_2}
&
\cat{C}_2
} \] 
It turns out that 
\[ (\eta_2 * \zeta_2) \circ (\eta_1 * \zeta_1) = 
(\eta_2 \circ \eta_1) * (\zeta_2 \circ \zeta_1) \]
as morphisms $E_2 \circ E_1 \twoto G_2 \circ G_1$. 
This is called the {\em exchange property}. 

\begin{exer} \label{exer:1312}
Prove the exchange property.
\end{exer}

Just like abstract categories, we can talk about 
triangulated categories. There is the $2$-category $\cat{TrCat}$ 
of all ($\K$-linear) triangulated categories. The objects here are 
the triangulated categories $(\cat{K}, \opn{T})$; the $1$-morphisms
are the triangulated functors $(F, \tau)$; and the $2$-morphisms
are the morphisms of triangulated functors $\eta$. This is what we are going 
to use later in this section.




\mysubsection{Functor Categories} \label{subsec:fun-cats}

In this subsection we isolate a part of $2$-category theory. This simplifies 
the discussion greatly. All set theoretical issues (sizes of sets) are 
neglected. As before, this can be treated by introducing a bigger universe. 

\begin{dfn} \label{dfn:2025}
Given categories $\cat{C}$ and $\cat{D}$, let 
$\cat{Fun}(\cat{C}, \cat{D})$
\index{1-Fun@$\cat{Fun}(\cat{C}, \cat{D})$}
be the category whose objects are the functors $F : \cat{C} \to \cat{D}$.
Given objects $F, G \in \cat{Fun}(\cat{C}, \cat{D})$,
the morphisms $\eta : F \twoto G$ in 
$\cat{Fun}(\cat{C}, \cat{D})$ are the morphisms of functors, i.e.\ the  natural 
transformations. 
\end{dfn}

For the case $\cat{D} = \cat{Set}$ this concept was already mentioned in 
Subsection \ref{subsec:repfunc}. 

In terms of the $2$-category $\cat{Cat}$ from the previous subsection,
$\cat{C}$ and $\cat{D}$ are objects of $\cat{Cat}$; the objects of 
$\cat{Fun}(\cat{C}, \cat{D})$ are $1$-morphisms in 
$\cat{Cat}$; and the morphisms  of 
$\cat{Fun}(\cat{C}, \cat{D})$ are $2$-morphisms in 
$\cat{Cat}$. Thus we made a ``reduction of order'', from $2$ to $1$, by 
passing from $\cat{Cat}$ to $\cat{Fun}(\cat{C}, \cat{D})$.

Suppose $G : \cat{C}' \to \cat{C}$ and 
$H : \cat{D} \to \cat{D}'$ are functors. There is an induced functor
\begin{equation} \label{eqn:2025}
\cat{Fun}(G, H) : \cat{Fun}(\cat{C}, \cat{D}) \to 
\cat{Fun}(\cat{C}', \cat{D}')
\end{equation}
defined by 
$\cat{F}(G, H)(F) := H \circ F \circ G$. 

\begin{prop} \label{prop:2025}
If $G$ and $H$ are equivalences, then the functor $\cat{Fun}(G, H)$ in 
\tup{(\ref{eqn:2025})} is an equivalence.
\end{prop}

\begin{exer} \label{exer:2035}
Prove Proposition \ref{prop:2025}.
\end{exer}

Recall that for a category $\cat{C}$ and a multiplicatively closed set of 
morphisms $\cat{S} \sub \cat{C}$ we denote by $\cat{C}_{\cat{S}}$ the 
localization. It comes with the localization functor 
$\opn{Q} : \cat{C} \to \cat{C}_{\cat{S}}$,
which is the identity on objects. See Definition \ref{dfn:1406}. 

For a category $\cat{E}$ let $\cat{E}^{\times} \sub \cat{E}$
be the category of isomorphisms; it has all the objects, but its  
morphisms are just the isomorphisms in $\cat{E}$. 

\begin{dfn} \label{dfn:2037}
Given categories $\cat{C}$ and $\cat{E}$, a multiplicatively closed set of 
morphisms $\cat{S} \sub \cat{C}$, and a functor $F : \cat{C} \to \cat{E}$,
we say that {\em $F$ is localizable to $\cat{S}$} if 
$F(\cat{S}) \sub  \cat{E}^{\times}$. 
We denote by
$\cat{Fun}_{\cat{S}}(\cat{C}, \cat{E})$%
\index{1-Fun@$\cat{Fun}_{\cat{S}}(\cat{C}, \cat{D})$}
the full subcategory of $\cat{Fun}(\cat{C}, \cat{E})$
on the functors that are localizable to $\cat{S}$.
\end{dfn}

Here is a useful formulation of the universal property of localization
of categories (see Definition \ref{dfn:1406}). 
Recall that a functor is an isomorphism of categories iff it is an equivalence 
that is bijective on sets of objects. 

\begin{prop} \label{prop:2026} 
Let $\cat{C}$ and $\cat{E}$ be categories, and let $\cat{S} \sub \cat{C}$ be a 
multiplicatively closed set of morphisms. Then the functor 
\[ \cat{Fun}(\opn{Q}, \opn{Id}_{\cat{E}}) : 
\cat{Fun}(\cat{C}_{\cat{S}}, \cat{E}) \to 
\cat{Fun}_{\cat{S}}(\cat{C}, \cat{E}) \]
is an isomorphism of categories.
\end{prop}

\begin{exer} \label{exer:2036}
Prove Proposition \ref{prop:2026}.
\end{exer}

By definition a bifunctor 
$F : \cat{C} \times \cat{D} \to \cat{E}$
is a functor from the product category $\cat{C} \times \cat{D}$.
See Subsection \ref{subsec:bifunc}.
It will be useful to retain both meanings; so we shall write 
\begin{equation} \label{eqn:2035}
\cat{BiFun}(\cat{C} \times \cat{D}, \cat{E}) := 
\cat{Fun}(\cat{C} \times \cat{D}, \cat{E}) ,
\end{equation}
where in the first expression we recall that $\cat{C} \times \cat{D}$ is a 
product. 

The next proposition describes bifunctors in a non-symmetric fashion. 

\begin{prop} \label{prop:2027} 
Let $\cat{C}$, $\cat{D}$ and $\cat{E}$ be categories. There is an isomorphism 
of categories 
\[ \Xi : \cat{Fun}(\cat{C} \times \cat{D}, \cat{E}) \to 
\cat{Fun}(\cat{C}, \cat{Fun}(\cat{D}, \cat{E})) \]
with the following formula\tup{:}
for a functor $F : \cat{C} \times \cat{D} \to \cat{E}$,
the functor
\[ \Xi(F) : \cat{C} \to \cat{Fun}(\cat{D}, \cat{E}) \]
is $\Xi(F)(C) := F(C, -)$.
\end{prop}

\begin{exer} \label{exer:2037}
Prove Proposition \ref{prop:2027}.
\end{exer}

\begin{prop}  \label{prop:2028} 
Let $\cat{C}$ and $\cat{D}$ be categories, and let 
$\cat{S} \sub \cat{C}$ and $\cat{T} \sub \cat{D}$
be multiplicatively closed sets of morphisms. Then
$\cat{S} \times \cat{T}$ is a multiplicatively closed set of morphisms
in $\cat{C} \times \cat{D}$, and the canonical functor 
\[ \Theta:  (\cat{C} \times \cat{D})_{\cat{S} \times \cat{T}} \to 
\cat{C}_{\cat{S}} \times \cat{D}_{\cat{T}} \]
is an isomorphism of categories. 
\end{prop}

\begin{proof}
It is clear that
$\cat{S} \times \cat{T} \sub \cat{C} \times \cat{D}$  
is closed under compositions. Let's spell out how the functor 
$\Theta$ arises. The functor
\[ \opn{Q}_{\cat{S} \times \cat{T}} : \cat{C} \times \cat{D} \to 
(\cat{C} \times \cat{D})_{\cat{S} \times \cat{T}} \]
is universal for functors 
$F : \cat{C} \times \cat{D} \to \cat{E}$
that invert $\cat{S} \times \cat{T}$; namely functors $F$ such that 
$F(s, t) \in \cat{E}^{\times}$ for all $s \in \cat{S}$ and $t \in \cat{T}$.
Since the functor 
\[ \opn{Q}_{\cat{S}} \times \opn{Q}_{\cat{T}} : \cat{C} \times \cat{D} \to 
\cat{C}_{\cat{S}} \times \cat{D}_{\cat{T}} \]
inverts $\cat{S} \times \cat{T}$, we get the functor $\Theta$. We will prove 
that the functor $\opn{Q}_{\cat{S}} \times \opn{Q}_{\cat{T}}$ 
has the same universal property as $\opn{Q}_{\cat{S} \times \cat{T}}$; this 
will imply that $\Theta$ is an isomorphism. 

Consider an arbitrary category $\cat{E}$. 
Invoking Propositions \ref{prop:2027} and \ref{prop:2026} we get a commutative 
diagram 

\begin{equation} \label{eqn:3383}
\UseTips  \xymatrix @C=5ex @R=8ex {
\cat{Fun}(\cat{C}_{\cat{S}} \times \cat{D}_{\cat{T}}, \cat{E})
\ar[d]^{\cat{Fun}
(\opn{Q}_{\cat{S}} \times \opn{Q}_{\cat{T}}, \opn{Id}_{\cat{E}})}
\ar[r]^(0.44){\cong}
&
\cat{Fun}(\cat{C}_{\cat{S}}, \cat{Fun}(\cat{D}_{\cat{T}}, \cat{E})) 
\ar[r]^{\cong}
&
\cat{Fun}_{\cat{S}}(\cat{C}, \cat{Fun}_{\cat{T}}(\cat{D}, \cat{E}))
\ar[d]^{\tup{f.f.\ emb.}}
\\
\cat{Fun}_{}(\cat{C} \times \cat{D}, \cat{E})
\ar[rr]^{\Xi}_{\cong}
&
&
\cat{Fun}_{}(\cat{C}, \cat{Fun}_{}(\cat{D}, \cat{E}))
}
\end{equation}

\noindent in which the right vertical arrow is a fully faithful embedding. 
Take any functor $F$ in the bottom left corner of (\ref{eqn:3383}), 
and let $F' := \Xi(F)$, which lives in the bottom right corner. 
Now $F'$ belongs to the top right corner of (\ref{eqn:3383})
iff $F'(s)(t) \in \cat{E}^{\times}$ for all $s \in \cat{S}$ and $t \in \cat{T}$.
But $F'(s)(t) = F(s, t)$, so this happens iff $F$ inverts 
$\cat{S} \times \cat{T}$.
\end{proof}

Denominator sets were introduced in Definition \ref{dfn:21}.

\begin{prop} \label{prop:2100}
In the situation of Proposition \tup{\ref{prop:2028}}, the following two 
conditions are equivalent\tup{:}
\begin{enumerate}
\rmitem{i} The multiplicatively closed sets  
$\cat{S} \sub \cat{C}$ and $\cat{T} \sub \cat{D}$ are left \tup{(}resp.\ 
right\tup{)} denominator sets. 

\rmitem{ii} The multiplicatively closed set
$\cat{S} \times \cat{T} \sub \cat{C} \times \cat{D}$ is a left 
\tup{(}resp.\ right\tup{)} denominator set. 
\end{enumerate}
\end{prop}

\begin{exer} \label{exer:2100}
Prove Proposition \ref{prop:2100}.
\end{exer}




\mysubsection{Abstract Derived Functors} \label{subsec:abstr-der-funs}

Here we deal with right and left derived functors in an abstract setup (as 
opposed to the triangulated setup). 

\begin{dfn}[Right Derived Functor] \label{dfn:2080}
Consider a category $\cat{K}$ and a multiplicatively closed set of morphisms 
$\cat{S} \sub \cat{K}$, with localization functor 
$\opn{Q} : \cat{K} \to \cat{K}_{\cat{S}}$. 
Let $F : \cat{K} \to \cat{E}$ be a functor. A
{\em right derived functor of $F$ with respect to $\cat{S}$}
\index{Derived functor! abstract right}
is a pair $(\mrm{R} F, \eta^{\mrm{R}})$, where 
\[ \mrm{R} F : \cat{K}_{\cat{S}} \to \cat{E} \]
is a functor, and 
\[ \eta^{\mrm{R}} : F \twoto \mrm{R} F \circ \opn{Q} \] 
is a morphism of functors, such that the following universal property holds:
\begin{itemize}
\item[(R)] Given any pair $(G, \th)$, consisting of a functor
$G : \cat{K}_{\cat{S}} \to \cat{E}$ 
and a morphism of functors 
$\th : F \twoto G \circ \opn{Q}$,
there is a unique morphism of functors 
$\mu : \mrm{R}^{} F \twoto G$
such that 
$\th = (\mu \circ \opn{id}_{\opn{Q}}) * \eta^{\mrm{R}}$.
\end{itemize}
\end{dfn}

Pictorially: there is a $2$-diagram 
\[ \UseTips  \xymatrix @C=12ex @R=8ex  {
\cat{K}
\ar[r]^{F}_(0.45){}="q" 
\ar[d]_{\opn{Q}}
& 
\cat{E} 
\\
\cat{K}_{\cat{S}}
\ar[ur]_{\mrm{R} F} ^(0.45){}="f"
\ar@{=>}  "q";"f" _{\eta^{\mrm{R}}}
} \]
For any other pair $(G, \th)$ there is a unique morphism $\mu$ that sits in 
this $2$-diagram 
\[ \UseTips  \xymatrix @C=14ex @R=10ex {
\cat{K}
\ar[r]^{F} _(0.45){}="q"  ^(0.55){}="q2" 
\ar[d]_{\opn{Q}}
& 
\cat{E} 
\\
\cat{K}_{\cat{S}}
\ar[ur] ^(0.45){}="f" ^(0.5){}="f1" 
\ar@{=>}  "q";"f" _{\eta^{\mrm{R}}}
\ar@(r,d)[ru]_{G} ^(0.55){}="g"  ^(0.48){}="g1"
\ar@{=>}  "q2";"g" ^(0.7){\th}
\ar@{=>}  "f";"g1" _(0.45){\mu}
} \]
such that the diagram of $2$-morphisms (with $*$ composition) 
\[ \UseTips \xymatrix @C=8ex @R=6ex { 
F
\ar@{=>}[d]_(0.45){\eta^{\mrm{R}}}
\ar@{=>}[dr]^{\th}
&
\\
\mrm{R} F \circ \opn{Q}
\ar@{=>}[r]_{\mu \, \circ \, \opn{id}_{\opn{Q}}}
&
G \circ \opn{Q}
} \]
is commutative. 

\begin{prop} \label{prop:1975}
If a right derived functor $(\mrm{R} F, \eta^{\mrm{R}})$ exists, then it is
unique, up to a unique isomorphism. Namely, if  $(G, \th)$ is another right 
derived functor of $F$, then there is a unique isomorphism of functors 
$\mu : \mrm{R} F \twoiso G$ such that 
$\th = (\mu \circ \opn{id}_{\opn{Q}}) * \eta^{\mrm{R}}$.
\end{prop}

\begin{proof}
Despite the apparent complication of the situation, the usual argument for 
uniqueness of universals applies. 
To be explicit, given another right derived functor $(G, \th)$, let 
$\mu : \mrm{R} F \twoto G$
be the unique morphism that's guaranteed by property (R) of the pair 
$(\mrm{R} F, \eta^{\mrm{R}})$. Then let 
$\nu : G \twoto \mrm{R} F$ be the unique morphism that's guaranteed by 
property (R) of the pair $(G, \th)$. Then the morphisms 
$\opn{id}_{\mrm{R} F}, \nu * \mu : \mrm{R} F \twoto \mrm{R} F$ 
satisfy 
\[ (\opn{id}_{\mrm{R} F} \circ \opn{id}_{\opn{Q}}) * \mu = \mu = 
((\nu * \mu) \circ \opn{id}_{\opn{Q}}) * \mu . \]
The uniqueness in property (R) of the pair $(\mrm{R} F, \eta^{\mrm{R}})$
implies that 
$\opn{id}_{\mrm{R} F} = \nu * \mu$. Likewise 
$\opn{id}_{G} = \mu * \nu$. 
Hence $\mu$ is an isomorphism of functors. 
\end{proof}

Here is a rather general existence result. 
Left denominator sets were introduced in Definition \ref{dfn:2311}. 

\begin{thm}[Existence of Right Derived Functor] \label{thm:2080}
In the situation of Definition \tup{\ref{dfn:2080}}, 
assume there is a full subcategory $\cat{J} \subseteq \cat{K}$
such the following three conditions hold\tup{:}
\begin{enumerate}
\rmitem{a} The multiplicatively closed set $\cat{S}$ 
is a left denominator set in $\cat{K}$. 

\rmitem{b} For every object $M \in \cat{K}$ there is a morphism 
$\rho : M \to I$ in $\cat{S}$, with target $I \in \cat{J}$. 

\rmitem{c} If $\psi$ is a morphism in $\cat{S} \cap \cat{J}$, then 
$F(\psi)$ is an isomorphism in $\cat{E}$.
\end{enumerate}
Then the right derived functor
\[ (\mrm{R} F, \eta^{\mrm{R}}) : \cat{K}_{\cat{S}} \to \cat{E} \]
exists. Moreover, for every object $I \in \cat{J}$ the morphism  
\[ \eta^{\mrm{R}}_{I} : F(I) \to \mrm{R} F(I) \]
in $\cat{E}$ is an isomorphism.
\end{thm}

Condition (b) says that {\em $\cat{K}$ has enough right $\cat{J}$-resolutions}. 
Condition (c) says that {\em $\cat{J}$ is an $F$-acyclic category}. 

Theorem \ref{thm:2080} is \cite[Proposition 7.3.2]{KaSc2}. However their 
notation is different: what we call ``left denominator set'', they call ``right 
multiplicative system''. 

We need a definition and a few lemmas before giving the proof of the theorem.
In them we assume the situation of the theorem. 

\begin{dfn} \label{dfn:2082}
In the situation of Theorem \ref{thm:2080}, by a {\em system of right 
$\cat{J}$-resolutions} we mean a pair $(I, \rho)$, where 
$I : \opn{Ob}(\cat{K}) \to \opn{Ob}(\cat{J})$
is a function, and 
$\rho = \{ \rho_M \}_{M \in \opn{Ob}(\cat{K})}$ 
is a collection of morphisms
$\rho_M : M \to I(M)$ in $\cat{S}$.
Moreover, if $M \in \opn{Ob}(\cat{J})$, then 
$I(M) = M$ and $\rho_M = \opn{id}_M$. 
\end{dfn}

Since here $\cat{K}$ has enough right $\cat{J}$-resolutions, it follows that
systems of right $\cat{J}$-resolutions $(I, \rho)$ exist.

Let us introduce some new notation that will make the proofs more readable:
\begin{equation} \label{eqn:2098}
\cat{K}' := \cat{J}, \quad \cat{S}' := \cat{J} \cap \cat{S}, 
\quad \cat{D} := \cat{K}_{\cat{S}} \quad \tup{and} 
\quad \cat{D}' := \cat{K}_{\cat{S}'}' . 
\end{equation}
The inclusion functor is $U : \cat{K}' \to \cat{K}$, and its localization is 
$V : \cat{D}' \to \cat{D}$.
These sit in a commutative diagram
\begin{equation} \label{eqn:2083}
\UseTips  \xymatrix @C=6ex @R=6ex {
\cat{K}'
\ar[r]^{U} _(0.5){}="s1"
\ar[d]_{\opn{Q}'}
&
\cat{K}
\ar[d]^{\opn{Q}}
\\
\cat{D}'
\ar[r]^{V} ^(0.5){}="t1"
&
\cat{D} 
}
\end{equation}

\begin{lem} \label{lem:2085}
The multiplicatively closed set $\cat{S}'$ is a left denominator set in 
$\cat{K}'$. 
\end{lem}

\begin{proof}
We need to verify conditions (LD1) and (LD2) in Definition \ref{dfn:2311}.

\medskip \noindent
(LD1): Given morphisms $a' : L' \to N'$ in $\cat{K}'$ and 
$s' : L' \to M'$ in $\cat{S}'$, we must find 
morphisms $b' : M' \to K'$ in $\cat{K}'$ and 
$t' : N' \to K'$ in $\cat{S}'$, such that $t' \circ a' = b' \circ s'$.
Because $\cat{S} \sub \cat{K}$ satisfies this condition, we can find 
morphisms $b : M' \to K$ in $\cat{K}$ and $t : N' \to K$ in $\cat{S}$
such that $t \circ a' = b \circ s'$. There is a morphism $\rho : K \to K'$ in 
$\cat{S}$ with target $K' \in \cat{K}'$. Then the morphisms $t' := \rho \circ 
t$ 
and $b' := \rho \circ b$ satisfy $t' \circ a' = b' \circ s'$, and 
$t' \in \cat{S}'$.

\medskip \noindent
(LD2): Given morphisms $a', b' : M' \to N'$ in $\cat{K}'$ and 
$s' : L' \to M'$ in $\cat{S}'$, that satisfy 
$a' \circ s' = b' \circ s'$, we must find 
a morphism $t' : N' \to K'$ in $\cat{S}'$ such that 
$t' \circ a' = t' \circ  b'$.
Because $\cat{S} \sub \cat{K}$ satisfies this condition, we can find 
a morphism $t : N' \to K$ in $\cat{S}$ such that 
$t \circ a' = t \circ  b'$.
There is a morphism $\rho : K \to K'$ in $\cat{S}$ with target 
$K' \in \cat{K}'$. Then the morphism $t' := \rho \circ t$ has the required 
property. 
\end{proof}

\begin{lem} \label{lem:2086}
The functor $V : \cat{D}' \to \cat{D}$ is an equivalence. 
\end{lem}

\begin{proof}
Condition (b) of the theorem implies that $V$ is essentially surjective on 
objects. We need to prove that $V$ is fully faithful. 
We shall use the left version of Proposition \ref{prop:3380}, namely $\cat{S}$ 
and $\cat{S}'$ are left denominator sets, and in condition (ii) the morphisms 
are $s : L' \to M$ and $t : M \to K'$. (See Remark \ref{rem:1845} and 
Propositions \ref{prop:3030} and \ref{prop:2310} regarding side changes.)
By Lemma \ref{lem:2085} the left version of condition (i) of 
Proposition \ref{prop:3380} holds. Condition (b) of the theorem implies 
the left version of condition (ii) of Proposition \ref{prop:3380}. 
Then the left version of Proposition \ref{prop:3380} says that 
$V$ is fully faithful. 
\end{proof}

\begin{lem} \label{lem:2081}
Suppose a system  of right $\cat{K}'$-resolutions 
$(I, \rho)$ has been chosen. Then the function 
$I : \opn{Ob}(\cat{K}) \to \opn{Ob}(\cat{K}')$ extends uniquely to a 
functor $I : \cat{D} \to \cat{D}'$, such that 
$I \circ V = \opn{Id}_{\cat{D}'}$, and 
$\opn{Q}(\rho) : \opn{Id}_{\cat{D}} \twoto V \circ I$
is an isomorphism of functors.
Therefore the functor $I$ is a a quasi-inverse of $V$. 
\end{lem}

The relevant $2$-diagram is this: 
\[ \UseTips  \xymatrix @C=10ex @R=8ex {
\cat{K}'
\ar[r]^{\opn{Q}'}
\ar[d]_{U}
&
\cat{D}'
\ar[dr]^{V} _(0.45){}="t" 
\ar[r]^{\opn{Id}}
&
\cat{D}'
\\
\cat{K}
\ar[r]^{\opn{Q}} _{}="q"
&
\cat{D} 
\ar[u]^{I}
\ar[r]_{\opn{Id}} ^(0.45){}="s"
\ar@{=>}  "s";"t" ^(0.45){\opn{Q}(\rho)}
&
\cat{D} 
\ar[u]_{I}
} \]
Recall that in a $2$-diagram, an empty polygon means it is commutative, namely 
it can be filled with 
$\xtwoto{\opn{id}}$. 

\begin{proof}
Consider a morphism $\psi : M \to N$ in $\cat{D}$.
Since $V : \cat{D}' \to \cat{D}$ is an equivalence, and since 
$V(I(M)) = I(M)$ and $V(I(N)) = I(N)$, there is a unique morphism 
$I(\psi) : I(M) \to I(N)$ 
in $\cat{D}'$ satisfying 
\begin{equation} \label{eqn:2085}
V(I(\psi)) := \opn{Q}(\rho_N) \circ \psi \circ \opn{Q}(\rho_M)^{-1} . 
\end{equation}
in $\cat{D}$. 

Let us check that $I : \cat{D} \to \cat{D}'$ is really a functor. 
Suppose $\phi : L \to M$ and $\psi : M \to N$ are morphisms in $\cat{D}$. 
Then 
\[ \begin{aligned}
& V \bigl( I(\psi) \circ I(\phi) \bigr) =  V(I(\psi)) \circ V(I(\phi))
\\ & \qquad 
 = \bigl( \opn{Q}(\rho_N) \circ \psi \circ \opn{Q}(\rho_M)^{-1} \bigr) \circ 
\bigl( \opn{Q}(\rho_M) \circ \phi \circ \opn{Q}(\rho_L)^{-1} \bigr)
\\ & \qquad 
= \opn{Q}(\rho_N) \circ (\psi \circ \phi) \circ \opn{Q}(\rho_L)^{-1} 
= V \bigl( I(\psi \circ \phi) \bigr) . 
\end{aligned} \]
Since $V$ is an equivalence, it follows that 
$I(\psi) \circ I(\phi) = I(\psi \circ \phi)$. 

Because $\rho_{M'} : M' \to I(M')$ is the identity for any object 
$M' \in \cat{K}'$, we see that there is equality 
$I \circ V = \opn{Id}_{\cat{D}'}$. 
By the defining formula (\ref{eqn:2085}) of $I(\psi)$ we have a commutative 
diagram 
\[ \UseTips  \xymatrix @C=10ex @R=6ex {
V(I(M))
\ar[r]^{V(I(\psi))}
&
V(I(N))
\\
M
\ar[r]^{\psi}
\ar[u]^{\opn{Q}(\rho_M)}
&
N
\ar[u]_{\opn{Q}(\rho_N)}
} \]
in $\cat{D}$. Hence $\opn{Q}(\rho) : \opn{Id}_{\cat{D}} \twoto V \circ I$
is an isomorphism of functors.
\end{proof}

Diagram (\ref{eqn:2083}) induces a commutative diagram of categories:
\begin{equation} \label{eqn:2084}
\UseTips \xymatrix @C=16ex @R=5.5ex {
\cat{Fun}(\cat{K}', \cat{E})
&
\cat{Fun}(\cat{K}, \cat{E})
\ar[l]_{\cat{Fun}(U, \opn{Id})}
\\
\cat{Fun}_{\cat{S}'}(\cat{K}', \cat{E})
\ar[u]_{\mrm{f.f. \, emb}}
&
\cat{Fun}_{\cat{S}}(\cat{K}, \cat{E})
\ar[u]^{\mrm{f.f. \, emb}}
\ar[l]^{\mrm{equiv}}_{\cat{Fun}(U, \opn{Id})}
\\
\cat{Fun}(\cat{D}', \cat{E})
\ar[u]^{\cat{Fun}(\opn{Q}', \opn{Id})}_{\mrm{isom}}
&
\cat{Fun}(\cat{D}, \cat{E})
\ar[l]_{\cat{Fun}(V, \opn{Id})}^{\mrm{equiv}}
\ar[u]_{\cat{Fun}(\opn{Q}, \opn{Id})}^{\mrm{isom}}
} 
\end{equation}
The vertical arrows marked ``f.f.\ emb'' are fully faithful embeddings by 
definition. According to Proposition \ref{prop:2026}  the vertical arrows 
marked ``isom'' are isomorphisms of categories. And by Lemma \ref{lem:2086} and 
Proposition \ref{prop:2025} the arrow $\cat{Fun}(V, \opn{Id})$  in the bottom 
row is an equivalence. As a consequence, the arrow $\cat{Fun}(U, \opn{Id})$ in 
the middle row is also an equivalence.

Let's introduce the notation 
$F' := F \circ U : \cat{K}' \to \cat{E}$.
This functor is an object of the category in the middle left of diagram 
(\ref{eqn:2084}), since, by condition (c) of Theorem \ref{thm:2080}, it inverts 
$\cat{S}'$.  

\begin{lem} \label{lem:3390}
Let $G : \cat{D} \to \cat{E}$ be a functor. Given a morphism 
$\th' : F' \twoto G \circ \opn{Q} \circ \, U$
of functors $\cat{K}' \to \cat{E}$, there is a unique morphism 
$\th : F \twoto G \circ \opn{Q}$
of functors $\cat{K} \to \cat{E}$ s.t.\ 
$\th' = \th \circ \opn{id}_{U}$. 
\end{lem}

Note that $G$ is an object in the category in the bottom right of diagram 
(\ref{eqn:2084}). The morphisms $\th'$ and $\th$ are in the middle left and top 
right respectively of this diagram. 

\begin{proof}
For every object $M \in \cat{K}$ there is a morphism 
$\rho : M \to I$ in $\cat{S}$ with target $I \in \cat{K}'$. 
A morphism of functors $\th$ satisfying $\th' = \th \circ \opn{id}_{U}$
will make this diagram commutative:
\begin{equation} \label{eqn:3390}
\UseTips  \xymatrix @C=8ex @R=6ex {
F(M)
\ar[r]^(0.4){\th_M}
\ar[d]_{F(\rho)}
&
(G \circ \opn{Q})(M)
\ar[d]^{(G \circ \opn{Q})(\rho)}_{\cong}
\\
F(I)
\ar[r]^(0.4){\th'_I}
&
(G \circ \opn{Q})(I)
} 
\end{equation}
We are using the facts that $I = U(I)$ and that $\opn{Q}(\rho)$ is an 
isomorphism. This proves the uniqueness of $\th$. 

For existence, let us choose a system  of right $\cat{K}'$-resolutions 
$(I, \rho)$, and define $\th = \{ \th_M \}$ using (\ref{eqn:3390}), namely 
\begin{equation} \label{eqn:3394}
\th_M := (G \circ \opn{Q})(\rho_M)^{-1} \circ \th'_{I(M)} \circ 
F(\rho_M) .
\end{equation}
We must prove that this is indeed a morphism of functors 
$\cat{K} \to \cat{E}$. Namely, for a given morphism 
$\phi : M \to N$ in $\cat{K}$, we have to prove that the diagram 
\begin{equation} \label{eqn:3391}
\UseTips  \xymatrix @C=8ex @R=6ex {
F(M)
\ar[r]^(0.4){\th_M}
\ar[d]_{F(\phi)}
&
(G \circ \opn{Q})(M)
\ar[d]^{(G \circ \opn{Q})(\phi)}
\\
F(N)
\ar[r]^(0.4){\th_N}
&
(G \circ \opn{Q})(N)
}
\end{equation}
in $\cat{E}$ is commutative. 

Lemma \ref{lem:2085} tells us that the morphism $I(\opn{Q}(\phi))$ in 
$\cat{D}'$ can be written as a left fraction
\begin{equation} \label{eqn:4715}
I(\opn{Q}(\phi)) = \opn{Q}'(\psi_1)^{-1} \circ \opn{Q}'(\psi_0)
\end{equation}
of morphisms $\psi_0 \in \cat{K}'$ and $\psi_1 \in \cat{S}'$. 
We get these diagrams: 
\begin{equation} \label{eqn:2072}
\UseTips \xymatrix @C=4ex @R=5.5ex {
M
\ar[rr]^{\phi}
\ar[d]_{\rho_M}
&
&
N
\ar[d]^{\rho_N}
\\
I(M)
\ar[dr]_{\psi_0}
&
&
I(N)
\ar[dl]^{\psi_1}
\\
&
J
} 
\quad \quad 
\UseTips \xymatrix @C=4,5ex @R=5.5ex {
M
\ar[rr]^{\opn{Q}(\phi)}
\ar[d]_{\opn{Q}(\rho_M)}
&
&
N
\ar[d]^{\opn{Q}(\rho_N)}
\\
I(M)
\ar[rr]^{ V(I(\opn{Q}(\phi))) }
\ar[dr]_{\opn{Q}(\psi_0)}
&
&
I(N)
\ar[dl]^{\opn{Q}(\psi_1)}
\\
&
J
} 
\end{equation}
The first diagram is in the category $\cat{K}$, and it might fail to be 
commutative. The second diagram is in the category $\cat{D}$, and it is 
commutative: the bottom triangle commutes by formula (\ref{eqn:4715}) and the 
equality $V \circ \opn{Q}' = \opn{Q}$; and the top square commutes by 
formula (\ref{eqn:2085}). By condition (LO4) of the left Ore localization 
$\opn{Q} : \cat{K} \to \cat{D}$, there is a morphism 
$\psi : J \to L$ in $\cat{S}$ such that 
$\psi \circ \psi_0 \circ \rho_M = \psi \circ \psi_1 \circ \rho_N \circ \phi$
in $\cat{K}$. There is the morphism 
$\rho_L : L \to I(L)$ in $\cat{S}$, whose target $I(L)$ belongs to $\cat{K}'$. 
Thus, after replacing the object $J$ with $I(L)$, the morphism 
$\psi_0$ by $\rho_L \circ \psi \circ \psi_0$, and the morphism 
$\psi_1$ by $\rho_L \circ \psi \circ \psi_1$, and noting that the latter is a 
morphism in $\cat{S}'$, we can now assume that the first diagram in 
(\ref{eqn:2072}) commutative too.

Now we embed (\ref{eqn:3391}) in the bigger diagram (\ref{eqn:3392})
in $\cat{E}$. 
Since $\rho_N \in \cat{S}$ and $\psi_1 \in \cat{S}'$, 
the morphisms $(G \circ \opn{Q})(\rho_N)$, $(G \circ \mrm{Q})(\psi_1)$ and 
$F(\psi_1)$ are isomorphisms. The top and bottom squares in (\ref{eqn:3392}) 
are commutative by the definition of $\th_M$ and $\th_N$, see formula 
(\ref{eqn:3394}). The left and right rounded shapes (those involving $J$) are 
commutative because the first diagram in (\ref{eqn:2072}) is commutative. 
\begin{equation} \label{eqn:3392}
\UseTips  \xymatrix @C=6ex @R=5.5ex {
&
F(I(M))
\ar[r]^(0.45){\th'_{I(M)}}
\ar @/_2.5em/ [ddl]_{F(\psi_0)}
&
(G \circ \opn{Q})(I(M))
\ar @/^3em/ [ddr]^{(G \circ \mrm{Q})(\psi_0)}
\\
&
F(M)
\ar[u]^{F(\rho_M)}
\ar[r]^(0.4){\th_M}
\ar[d]_{F(\phi)}
&
(G \circ \opn{Q})(M)
\ar[u]_{(G \circ \opn{Q})(\rho_M)}^{\cong}
\ar[d]^{(G \circ \opn{Q})(\phi)}
\\
F(J)
&
F(N)
\ar[r]^(0.4){\th_N}
\ar[d]_{F(\rho_N)}
&
(G \circ \opn{Q})(N)
\ar[d]^{(G \circ \opn{Q})(\rho_N)}_{\cong}
&
(G \circ \opn{Q})(J)
\\
&
F(I(N))
\ar@(l,d)[ul]^{F(\psi_1)}_{\cong}
\ar[r]^(0.45){\th'_{I(N)}}
&
(G \circ \opn{Q})(I(N))
\ar@(r,d)[ur]_{\ (G \circ \mrm{Q})(\psi_1)}^{\cong}
}
\end{equation}

\medskip 
Since $\th' : F' \twoto G \circ \opn{Q} \circ \, U$ is a morphism of functors, 
we have a commutative diagram 
\begin{equation} \label{eqn:4700}
\UseTips  \xymatrix @C=8ex @R=5.5ex {
F(I(M))
\ar[r]^(0.45){\th'_{I(M)}}
\ar[d]_{F(\psi_0)}
&
(G \circ \opn{Q})(I(M))
\ar[d]^{(G \circ \opn{Q})(\psi_0)}
\\
F(J)
\ar[r]^(0.45){\th'_{J}}
&
(G \circ \opn{Q})(J)
\\
F(I(N))
\ar[u]^{F(\psi_1)}
\ar[r]^(0.45){\th'_{I(N)}}
&
(G \circ \opn{Q})(I(N))
\ar[u]_{(G \circ \opn{Q})(\psi_1)}
}
\end{equation}

\medskip \noindent
When we erase the arrow $\th'_{J}$ from diagram (\ref{eqn:4700}), we obtain the 
outer boundary of diagram (\ref{eqn:3392}). Therefore diagram 
(\ref{eqn:3392}) is commutative. In particular, the middle square of diagram 
(\ref{eqn:3392}) is commutative, and it is precisely diagram (\ref{eqn:3391}).
\end{proof}

\begin{proof}[Proof of Theorem \tup{\ref{thm:2080}}] 
The proof is divided into four steps.  

\smallskip \noindent 
Step 1. Recall that the functor $F' = F \circ U$ lives in the middle left term 
in diagram (\ref{eqn:2084}). Because the arrow $\cat{Fun}(\opn{Q}', \opn{Id})$
is an isomorphism, there is a unique functor 
$\mrm{R} F'$ living in the bottom left term of diagram (\ref{eqn:2084}) that 
satisfies $\mrm{R} F' \circ \opn{Q}' = F'$. See next commutative diagram. 
\begin{equation} \label{eqn:2091}
\UseTips \xymatrix @C=10ex @R=6ex {
\cat{K}'
\ar[r]^{F'}
\ar[d]_{\opn{Q}'}
&
\cat{E}
\\
\cat{D}'
\ar[ur]_{\mrm{R} F'}
} 
\end{equation}

Let $\eta' := \opn{id}_{F'}$.
We claim that the pair $(\mrm{R} F', \eta')$ is a right derived functor of 
$F'$. Indeed, suppose we are given a pair $(G', \th')$, 
where $G'$ is a functor in the bottom left corner of diagram (\ref{eqn:2084}), 
and $\th' : F' \twoto G'  \circ \opn{Q}'$ 
is a morphism in the top left corner of that diagram. See the $2$-diagram
(\ref{eqn:2093}). 
Because the function 
\begin{equation} \label{eqn:2092} 
\opn{Hom}_{\cat{Fun}(\cat{D}', \cat{E})}
(\mrm{R} F', G') \to 
\opn{Hom}_{\cat{Fun}(\cat{K}', \cat{E})}
(F', G' \circ \opn{Q}') 
\end{equation}
is bijective -- this is the left edge of diagram (\ref{eqn:2084}) -- 
there is a unique morphism 
$\mu' : \mrm{R} F' \twoto G'$ that goes to $\th'$ under (\ref{eqn:2092}). 
\begin{equation} \label{eqn:2093} 
\UseTips  \xymatrix @C=14ex @R=10ex {
\cat{K}'
\ar[r]^{F'} _(0.45){}="q"  ^(0.6){}="q2" 
\ar[d]_{\opn{Q}'}
& 
\cat{E} 
\\
\cat{D}'
\ar[ur]^(0.3){\mrm{R} F'} ^(0.45){}="f" ^(0.5){}="f1" 
\ar@{=>}  "q";"f" _{\eta'}
\ar@(r,d)[ru]_{G'} ^(0.55){}="g"  ^(0.48){}="g1"
\ar@{=>}  "q2";"g" ^(0.7){\th'}
\ar@{=>}  "f";"g1" _(0.45){\mu'}
} 
\end{equation}


\medskip \noindent 
Step 2. Now we choose a system of right $\cat{K}'$-resolutions 
$(I, \rho)$, in the sense of Definition \ref{dfn:2082}.
By Lemma \ref{lem:2081}  we get an equivalence of categories
$I : \cat{D} \to \cat{D}'$, that is a quasi-inverse to $V$, and 
an isomorphism of functors 
$\opn{Q}(\rho) : \opn{Id}_{\cat{D}} \iso V \circ I$. 
See the following $2$-diagram (the solid arrows). 
\begin{equation} \label{eqn:2094} 
\UseTips  \xymatrix @C=10ex @R=8ex {
\cat{K}'
\ar[r]^{\opn{Q}'}
\ar[d]_{U}
\ar@(u,u)[rrr]^{F'}
&
\cat{D}'
\ar[dr]^{V} _(0.45){}="t" 
\ar[r]^{\opn{Id}}
&
\cat{D}'
\ar[r]^{\mrm{R} F'}
&
\cat{E}
\\
\cat{K}
\ar[r]_{\opn{Q}} _{}="q"
&
\cat{D} 
\ar[u]^{I}
\ar[r]_{\opn{Id}} ^(0.45){}="s"
\ar@{=>}  "s";"t" ^(0.45){\opn{Q}(\rho)}
&
\cat{D} 
\ar[u]_{I}
\ar@{-->}[ur]_{\mrm{R} F}
} 
\end{equation}

Define the functor 
\begin{equation} \label{eqn:2095}  
\mrm{R} F := \mrm{R} F' \circ I :  \cat{D} \to \cat{E} .
\end{equation}
It is the dashed arrow in diagram (\ref{eqn:2094}). 
So the functor $\mrm{R} F$ lives in the bottom right corner of 
(\ref{eqn:2084}), and $\mrm{R} F' = \mrm{R} F \circ V$.

\medskip \noindent 
Step 3. Consider Lemma \ref{lem:3390}, with the functor 
$G := \mrm{R} F$, and the morphism of functors
\[ \eta' = \opn{id}_{F'} : F' \twoto \mrm{R} F' \circ \opn{Q}' = 
\mrm{R} F \circ \opn{Q} \circ \, U . \]
The lemma says that there is a unique morphism of functors 
$\eta^{\mrm{R}} : F \twoto \mrm{R} F \circ \opn{Q}$
s.t.\ $\eta^{\mrm{R}} \circ \opn{id}_U = \eta'$. 

We can give an explicit formula for the morphism of functors $\eta^{\mrm{R}}$. 
Take an object $M \in \cat{K}$. Then the morphism 
$\eta^{\mrm{R}}_M : F(M) \to \mrm{R} F(M) = F(I(M))$
in $\cat{E}$ is nothing but 
\begin{equation} \label{eqn:2096} 
\eta^{\mrm{R}}_M := F(\rho_M) . 
\end{equation}
Here is the calculation. By formula (\ref{eqn:3394}) we have
\[ \eta^{\mrm{R}}_M = (\mrm{R} F \circ \opn{Q})(\rho_M)^{-1} \circ \eta'_{I(M)} 
\circ F(\rho_M) . \]
But 
$\eta'_{I(M)} = \opn{id}_{F(I(M))}$, 
\[ I(\opn{Q}(\rho_M)) = \opn{Q}(\rho_{I(M)}) \circ \opn{Q}(\rho_M) \circ 
\opn{Q}(\rho_M)^{-1} = \opn{id}_{I(M)} \]
by (\ref{eqn:2085}), and 
\[ (\mrm{R} F \circ \opn{Q})(\rho_M) = 
(\mrm{R} F' \circ I \circ \opn{Q})(\rho_M) = 
\mrm{R} F'(\opn{id}_{I(M)}) = \opn{id}_{F(I(M))} . \]
So everything else gets canceled, and we are left with (\ref{eqn:2096}). 

\medskip \noindent 
Step 4. It remains to prove that the pair $(\mrm{R} F, \eta^{\mrm{R}})$ is a 
right derived functor of $F$. Suppose $(G, \th)$ is a pair, where $G$ is a 
functor in the category in the bottom right corner of diagram (\ref{eqn:2084}), 
and $\th : F \twoto G  \circ \opn{Q}$ 
is a morphism in the top right corner of the diagram.
We are looking for a morphism 
$\mu : \mrm{R} F \twoto G$
in the bottom right category in diagram (\ref{eqn:2084}) for which 
$\th = (\mu \circ \opn{id}_{\opn{Q}}) * \eta^{\mrm{R}}$. 
Let $G' := G \circ V$, and let $\th' : F' \twoto G' \circ \opn{Q}'$
be the morphism in the top left corner of (\ref{eqn:2084}) corresponding to 
$\th$. Because of the equivalence $\cat{Fun}(V, \opn{Id})$, finding such $\mu$ 
is the same as finding a morphism
$\mu' : \mrm{R} F' \twoto G'$
in the bottom left category in diagram (\ref{eqn:2084}), satisfying
\begin{equation} \label{eqn:2075}
\th' = (\mu' \circ \opn{id}_{\opn{Q}'}) * \eta' .
\end{equation}

Finally, by step 1 the pair $(\mrm{R} F', \eta')$ is a right derived functor 
of $F'$. This says that there is a unique morphism 
$\mu'$ satisfying (\ref{eqn:2075}). 
\end{proof}

\begin{dfn} \label{dfn:1350}
The construction of the right derived functor $(\mrm{R} F, \eta^{\mrm{R}})$ in 
the proof of Theorem \ref{thm:2080}, and specifically formulas (\ref{eqn:2095}) 
and (\ref{eqn:2096}), is called a {\em presentation of 
$(\mrm{R} F, \eta^{\mrm{R}})$ by the system of right $\cat{J}$-resolutions 
$(I, \rho)$}. 
\end{dfn}

Of course any other right derived functor of $F$ (perhaps presented  
by another system of right $\cat{J}$-resolutions) is uniquely isomorphic to 
$(\mrm{R} F, \eta^{\mrm{R}})$. This is according to Proposition 
\ref{prop:1350}. In Sections \ref{sec:resol} and \ref{sec:exist-resol} we shall 
give several existence results for right resolutions by suitable acyclic 
objects. 

Now to left derived functors. 

\begin{dfn}[Left Derived Functor] \label{dfn:2090}
Consider a category $\cat{K}$ and a multiplicatively closed set of morphisms 
$\cat{S} \sub \cat{K}$, with localization functor 
$\opn{Q} : \cat{K} \to \cat{K}_{\cat{S}}$. 
Let $F : \cat{K} \to \cat{E}$ be a functor. A
{\em left derived functor of $F$ with respect to $\cat{S}$} 
\index{Derived functor! abstract left}
is a pair $(\mrm{L} F, \eta^{\mrm{L}})$, where 
\[ \mrm{L} F : \cat{K}_{\cat{S}} \to \cat{E} \]
is a functor, and 
\[ \eta^{\mrm{L}} : \mrm{L} F \circ \opn{Q} \twoto F \] 
is a morphism of functors, such that the following universal property holds:
\begin{itemize}
\item[(L)] Given any pair $(G, \th)$, consisting of a functor
$G : \cat{K}_{\cat{S}} \to \cat{E}$ 
and a morphism of functors 
$\th : G \circ \opn{Q} \twoto F$,
there is a unique morphism of functors 
$\mu : G  \twoto \mrm{L}^{} F$
such that 
$\th = \eta^{\mrm{L}} * (\mu \circ \opn{id}_{\opn{Q}})$.
\end{itemize}
\end{dfn}

Pictorially: there is a $2$-diagram 
\[ \UseTips  \xymatrix @C=12ex @R=8ex  {
\cat{K}
\ar[r]^{F} _(0.45){}="q" 
\ar[d]_{\opn{Q}}
& 
\cat{E} 
\\
\cat{K}_{\cat{S}}
\ar[ur]_{\mrm{L} F} ^(0.45){}="f"
\ar@{=>}  "f";"q" ^{\eta^{\mrm{L}}}
} \]
For any other pair $(G, \th)$ there is a unique morphism $\mu$ that sits in 
this $2$-diagram
\[ \UseTips  \xymatrix @C=14ex @R=10ex {
\cat{K}
\ar[r]^{F} _(0.45){}="q"  ^(0.55){}="q2" 
\ar[d]_{\opn{Q}}
& 
\cat{E} 
\\
\cat{K}_{\cat{S}}
\ar[ur] ^(0.45){}="f" ^(0.5){}="f1" 
\ar@{=>}  "f";"q" ^{\eta^{\mrm{L}}}
\ar@(r,d)[ru]_{G} ^(0.55){}="g"  ^(0.48){}="g1"
\ar@{=>}  "g";"q2" _(0.7){\th}
\ar@{=>}  "g1";"f" ^(0.45){\mu}
} \]
such that the diagram of $2$-morphisms (with $*$ composition) 
\[ \UseTips \xymatrix @C=8ex @R=6ex { 
F
&
\\
\mrm{L} F \circ \opn{Q}
\ar@{=>}[u]^{\eta^{\mrm{L}}}
&
G \circ \opn{Q}
\ar@{=>}[ul]_{\th}
\ar@{=>}[l]^(0.4){\mu \, \circ \, \opn{id}_{\opn{Q}}}
} \]
is commutative. 

We see that a left derived functor with target $\cat{E}$ amounts to a right 
derived functor with target $\cat{E}^{\mrm{op}}$. It means that new proofs are 
not needed. 

\begin{prop} \label{prop:2090}
If a left derived functor $(\mrm{L} F, \eta^{\mrm{L}})$ exists, then it is
unique, up to a unique isomorphism. Namely, if  $(G, \th)$ is another left 
derived functor of $F$, then there is a unique isomorphism of functors 
$\mu : G \twoiso \mrm{L} F$ such that 
$\th = \eta^{\mrm{L}} * (\mu \circ \opn{id}_{\opn{Q}})$.
\end{prop}

The proof is the same as that of Proposition \ref{prop:1975}, only some arrows 
have to be reversed. 

\begin{thm}[Existence of Left Derived Functor] \label{thm:2091}
In the situation of Definition \tup{\ref{dfn:2090}}, 
assume there is a full subcategory $\cat{P} \sub \cat{K}$
such the following three conditions hold\tup{:}
\begin{enumerate}
\rmitem{a} The multiplicatively closed set $\cat{S}$ 
is a right denominator set in $\cat{K}$. 

\rmitem{b} For every object $M \in \cat{K}$ there is a morphism 
$\rho : P \to M$ in $\cat{S}$, with source $P \in \cat{P}$. 

\rmitem{c} If $\psi$ is a morphism in $\cat{P} \cap \cat{S}$, then 
$F(\psi)$ is an isomorphism in $\cat{E}$.
\end{enumerate}
Then the left derived functor
\[ (\mrm{L} F, \eta^{\mrm{L}}) : \cat{K}_{\cat{S}} \to \cat{E} \]
exists. Moreover, for every object $P \in \cat{P}$ the morphism  
\[ \eta^{\mrm{L}}_{P} : \mrm{L} F(P) \to F(P) \]
in $\cat{E}$ is an isomorphism.
\end{thm}

Condition (b) says that {\em $\cat{K}$ has enough left $\cat{P}$-resolutions}. 
Condition (c) says that {\em $\cat{P}$ is an $F$-acyclic category}. 

The proof is the same as that of Theorem \ref{thm:2080}, only some arrows 
have to be reversed. We leave this as an exercise. 

\begin{exer} \label{exer:2330}
Prove Theorem \ref{thm:2091}, including the lemmas leading to it. 
(Hint: replace $\cat{E}$ with $\cat{E}^{\mrm{op}}$.)
\end{exer}

For reference we give the next definition.  

\begin{dfn} \label{dfn:2095}
In the situation of Theorem \ref{thm:2091}, by a {\em system of left 
$\cat{P}$-resolutions} we mean a pair $(P, \rho)$, where 
$P : \opn{Ob}(\cat{K}) \to \opn{Ob}(\cat{P})$
is a function, and 
$\rho = \{ \rho_M \}_{M \in \opn{Ob}(\cat{K})}$ 
is a collection of morphisms
$\rho_M : P(M) \to M$ in $\cat{S}$.
Moreover, if $M \in \opn{Ob}(\cat{P})$, then 
$P(M) = M$ and $\rho_M = \opn{id}_M$. 
\end{dfn}

When $\cat{K}$ has enough left $\cat{P}$-resolutions, there exists
a system of left $\cat{P}$-resolutions $(P, \rho)$ .

\begin{dfn} \label{dfn:2331}
The construction of the left derived functor $(\mrm{L} F, \eta^{\mrm{L}})$ in 
the proof of Theorem \ref{thm:2091} -- i.e.\ the left variant of the proof of 
Theorem \ref{thm:2080}, and specifically the left versions of formulas 
(\ref{eqn:2095}) and (\ref{eqn:2096}) -- 
is called a {\em presentation of $(\mrm{L} F, \eta^{\mrm{L}})$ by the system of 
left $\cat{P}$-resolutions $(P, \rho)$}. 
\end{dfn}

Of course any other left derived functor of $F$ (perhaps presented  
by another system of left $\cat{P}$-resolutions) is uniquely isomorphic to 
$(\mrm{L} F, \eta^{\mrm{L}})$. This is according to Proposition 
\ref{prop:2090}. In Sections \ref{sec:resol} and \ref{sec:exist-resol} we shall 
give several existence results for left resolutions by suitable acyclic 
objects. 

\begin{rem} \label{rem:4696}
The right derived functor $\mrm{R} F$ from Definition \ref{dfn:2080} is a {\em 
left Kan extension} of $F$ along $\opn{Q}$. Likewise, the left derived functor 
$\mrm{L} F$ from Definition \ref{dfn:2090} is a {\em right Kan extension} of 
$F$ along $\opn{Q}$. See \cite[Chapter X]{Mac2}.
\end{rem}




\mysubsection{Triangulated Derived Functors} \label{subsec:tri-der-funs}

In this subsection we specialize the definitions and results of the previous 
subsection to the case of triangulated functors between triangulated 
categories. 
There is a fixed nonzero commutative base ring $\K$, and all categories and 
functors here are $\K$-linear. 

Triangulated functors and morphisms between them were introduced in Definition
\ref{dfn:1276}. Recall that a triangulated functor is a pair $(F, \tau)$, 
consisting of a linear functor $F$ and a translation isomorphism $\tau$, 
that respect the distinguished triangles. 

Here is the triangulated version of Definition \ref{dfn:2025}. 

\begin{dfn} \label{dfn:2335}
Let $\cat{K}$ and $\cat{L}$ be triangulated categories. We define 
$\cat{TrFun}(\cat{K}, \cat{L})$
\index{1-TrFun@$\cat{TrFun}(\cat{K}, \cat{L})$}
to be the category whose objects are the triangulated functors 
$(F, \tau) : \cat{K} \to \cat{L}$. Given objects 
$(F, \tau)$ and $(G, \nu)$ of $\cat{TrFun}(\cat{K}, \cat{L})$,
the morphisms 
$\al : (F, \tau) \twoto (G, \nu)$ 
in  $\cat{TrFun}(\cat{K}, \cat{L})$
are the morphisms of triangulated functors. 
\end{dfn}

\begin{lem} \label{lem:2370}  
Let 
$(F, \tau) : (\cat{K}, \opn{T}_{\cat{K}}) \to (\cat{L}, \opn{T}_{\cat{L}})$
be a triangulated functor between 
triangulated categories. Assume $F$ is an equivalence \tup{(}of 
abstract categories\tup{)}, with quasi-inverse 
$G : \cat{L} \to \cat{K}$, and with adjunction isomorphisms 
$\al : G \circ F \twoiso \opn{Id}_{\cat{K}}$ and  
$\be : F \circ G \twoiso \opn{Id}_{\cat{L}}$. 
Then there is an isomorphism of functors 
$\nu : G \circ \opn{T}_{\cat{L}} \twoiso \opn{T}_{\cat{K}} \circ \, G$ 
such that 
$(G, \nu) : (\cat{L}, \opn{T}_{\cat{L}}) \to (\cat{K}, \opn{T}_{\cat{K}})$
is a triangulated functor, and $\al$ and $\be$ are isomorphisms of 
triangulated functors.
\end{lem}

\begin{proof}
It is well-known that $G$ is additive (or in our case, $\K$-linear); but 
since 
the proof is so easy, we shall reproduce it. Take any pair of objects 
$M, N \in \cat{L}$. We have to prove that the bijection
\[ G_{M, N} : \opn{Hom}_{\cat{L}}(M, N) \to 
\opn{Hom}_{\cat{K}} \bigl( G(M), G(N) \bigr) \]
is linear. But 
\[ G_{M, N} = (F_{G(M), G(N)})^{-1} \circ
\opn{Hom}_{\cat{L}}(\be_M, \be_N^{-1}) \]
as bijections (of sets) between these modules. Since 
$\opn{Hom}_{\cat{L}}(\be_M, \be_N^{-1})$
and $F^{-1}_{G(M), G(N)}$ are $\K$-linear, so is $G_{M, N}$.

We define the isomorphism of functors $\nu$ by the formula 
\[ \nu := (\al \circ \opn{id}_{\opn{T}_{\cat{K}} \circ \, G}) * 
(\opn{id}_G \circ \, \tau \circ \opn{id}_G)^{-1} * 
(\opn{id}_{G \circ \opn{T}_{\cat{L}}} \circ \, \be)^{-1} , \]
in terms of the $2$-categorical notation.
This gives rise to a commutative diagram of isomorphisms
\[ \UseTips \xymatrix @C=12ex @R=6ex { 
G \circ \opn{T}_{\cat{L}} \circ \, F \circ G
\ar@{=>}[d]_{\opn{id} \circ \, \be}
& 
G \circ F \circ \opn{T}_{\cat{K}} \circ \, G
\ar@{=>}[l]_{\opn{id} \circ \, \tau \, \circ \, \opn{id}}
\ar@{=>}[d]^{\al \, \circ \, \opn{id}}
\\
G \circ \opn{T}_{\cat{L}}
\ar@{=>}[r]^{\nu}
&
\opn{T}_{\cat{K}} \circ \, G
} \]
of additive functors $\cat{L} \to \cat{K}$. 
So the pair $(G, \nu)$ is a T-additive functor. 

The verification that $(G, \nu)$ preserves triangles (in the sense of 
Definition \ref{dfn:1276}(1)) is done like the proof of the additivity of 
$G$, but now using axiom (TR1)(a) from Definition \ref{dfn:1158}.
Also $\al$ and $\be$ are morphisms of triangulated functors 
(Definitions \ref{dfn:1276}(2) and \ref{dfn:1160}). 
We leave these verifications as an exercise. 
\end{proof}

\begin{exer}
Finish the proof above (the last assertions).  
\end{exer}

From here on in this subsection we shall usually keep the translation 
isomorphisms (such as $\tau$ and $\nu$ above) implicit. 

Suppose we are given triangulated functors 
$U : \cat{K}' \to \cat{K}$ and 
$V : \cat{L} \to \cat{L}'$. There is an induced functor 
\[ \cat{TrFun}(U, V) : 
\cat{TrFun}(\cat{K}, \cat{L}) \to \cat{TrFun}(\cat{K}', \cat{L}') ; \]
the formula is the same as in (\ref{eqn:2025}). 

\begin{lem} \label{lem:2335}
If $U$ and $V$ are equivalences, then the functor $\cat{TrFun}(U, V)$ is an 
equivalence.
\end{lem}

\begin{proof}
Use Proposition \ref{prop:2025} and  Lemma \ref{lem:2370}. 
\end{proof}

\begin{dfn} \label{dfn:2336}
Let $\cat{K}$ and $\cat{L}$ be triangulated categories,
and let $\cat{S} \sub \cat{K}$ be a denominator set of cohomological origin. 
We define 
$\cat{TrFun}_{\cat{S}}(\cat{K}, \cat{L})$
\index{1-TrFunS@$\cat{TrFun}_{\cat{S}}(\cat{K}, \cat{L})$}
to be the full subcategory of 
$\cat{TrFun}(\cat{K}, \cat{L})$
whose objects are the functors that are localizable to $\cat{S}$, in the 
sense of Definition \ref{dfn:2037}.
\end{dfn}

We know that the localization functor 
$\opn{Q} : \cat{K} \to \cat{K}_{\cat{S}}$ is a left and right Ore 
localization. 

\begin{lem} \label{lem:2336} 
Let $\cat{K}$ and $\cat{E}$ be triangulated categories,
and let $\cat{S} \sub \cat{K}$ be a denominator set of cohomological origin. 
Then the functor 
\[ \cat{TrFun}(\opn{Q}, \opn{Id}_{\cat{E}}) : 
\cat{TrFun}(\cat{K}_{\cat{S}}, \cat{E}) \to 
\cat{TrFun}_{\cat{S}}(\cat{K}, \cat{E}) \]
is an isomorphism of categories.
\end{lem}

\begin{proof}
Use Proposition \ref{prop:2026} and Theorem \ref{thm:105}. 
\end{proof}

Here is the triangulated version of Definition \ref{dfn:2080}.

\begin{dfn}  \label{dfn:131}
Let $F : \cat{K} \to \cat{E}$ be a triangulated functor between  
triangulated categories, and let 
$\cat{S} \sub \cat{K}$ be a denominator set of cohomological origin.    
A {\em triangulated right derived functor of $F$ with respect to $\cat{S}$} 
\index{Derived functor! triangulated right}
\index{Triangulated functor! right derived}
is a triangulated functor 
$\mrm{R}^{} F : \cat{K}_{\cat{S}} \to \cat{E}$,
together with a morphism 
$\eta^{\mrm{R}} : F \twoto \mrm{R}^{} F \circ \opn{Q}$ 
of triangulated functors $\cat{K} \to \cat{E}$. The pair 
$(\mrm{R} F, \eta^{\mrm{R}})$ must have this universal property:
\begin{itemize}
\item[(R)] Given any pair $(G, \th)$, consisting of a triangulated 
functor $G : \cat{K}_{\cat{S}} \to \cat{E}$ and a morphism of triangulated 
functors $\th : F \twoto G \circ \opn{Q}$, 
there is a unique morphism of triangulated functors 
$\mu : \mrm{R}^{} F \twoto G$
such that $\th = (\mu \circ \opn{id}_{\opn{Q}}) * \eta^{\mrm{R}}$. 
\end{itemize}
\end{dfn}

\begin{prop} \label{prop:1350}
If a triangulated right derived functor $(\mrm{R} F, \eta^{\mrm{R}})$ exists, 
then it is unique, up to a unique isomorphism. Namely, if  $(G, \th)$ is 
another triangulated right derived functor of $F$, then there is a unique 
isomorphism of 
triangulated functors $\mu : \mrm{R} F \twoiso G$ such that 
$\th = (\mu \circ \opn{id}_{\opn{Q}}) * \eta^{\mrm{R}}$. 
\end{prop}

The proof is the same as that of Proposition \ref{prop:1975}.
Now for the triangulated version of Theorem \ref{thm:2080}.

\begin{thm} \label{thm:1470}
In the situation of Definition \tup{\ref{dfn:131}}, 
assume there is a full triangulated subcategory 
$\cat{J} \subseteq \cat{K}$
with these two properties\tup{:}
\begin{enumerate}
\rmitem{a} Every object $M \in \cat{K}$ admits a morphism 
$\rho : M \to I$ in $\cat{S}$ with target $I \in \cat{J}$.

\rmitem{b} If $\psi$ is a morphism in $\cat{J} \cap \cat{S}$, 
then $F(\psi)$ is an isomorphism in $\cat{E}$.
\end{enumerate}
Then the triangulated right derived functor
$(\mrm{R} F, \eta^{\mrm{R}}) : \cat{K}_{\cat{S}} \to \cat{E}$ 
of $F$ with respect to $\cat{S}$ exists. 
\index{Derived functor! triangulated right}
\index{Triangulated functor! right derived}
Moreover, for every object $I \in \cat{J}$ the morphism  
$\eta^{\mrm{R}}_I : F(I) \to \mrm{R} F(I)$ 
in $\cat{E}$ is an isomorphism.
\end{thm}

\begin{proof}
It will be convenient to change notation to that used in the proof of Theorem 
\ref{thm:2080}. Let's define 
$\cat{K}' := \cat{J}$, $\cat{S}' := \cat{K}' \cap \cat{S}$ and
$\cat{D}' := \cat{K}'_{\cat{S}'}$. 
The localization functor of $\cat{K}'$ is 
$\opn{Q}' : \cat{K}' \to \cat{D}'$. 
The inclusion functor is 
$U : \cat{K}' \to \cat{K}$, and its localization is the functor
$V : \cat{D}' \to \cat{D}$.
We have this commutative diagram of triangulated functors between 
triangulated categories:
\begin{equation} \label{eqn:2340}
\UseTips \xymatrix @C=6ex @R=6ex {
\cat{K}'
\ar[r]^{U}
\ar[d]_{\opn{Q}'}
&
\cat{K}
\ar[d]^{\opn{Q}}
\\
\cat{D}'
\ar[r]^{V}
&
\cat{D}
} 
\end{equation}
The functor $U$ is fully faithful. 
By Lemma \ref{lem:2086} the functor $V$ is an equivalence.

Diagram (\ref{eqn:2340}) induces a commutative diagram of linear categories:
\begin{equation} \label{eqn:2346}
\UseTips \xymatrix @C=14ex @R=6ex {
\cat{TrFun}(\cat{K}', \cat{E})
&
\cat{TrFun}(\cat{K}, \cat{E})
\ar[l]_{\cat{Fun}(U, \opn{Id})}
\\
\cat{TrFun}_{\cat{S}'}(\cat{K}', \cat{E})
\ar[u]_{\mrm{f.f.\, emb}}
&
\cat{TrFun}_{\cat{S}}(\cat{K}, \cat{E})
\ar[u]^{\mrm{f.f.\, emb}}
\ar[l]^{\mrm{equiv}}_{\cat{Fun}(U, \opn{Id})}
\\
\cat{TrFun}(\cat{D}', \cat{E})
\ar[u]^{\cat{Fun}(\opn{Q}', \opn{Id})}_{\mrm{isom}}
&
\cat{TrFun}(\cat{D}, \cat{E})
\ar[l]_{\cat{Fun}(V, \opn{Id})}^{\mrm{equiv}}
\ar[u]_{\cat{Fun}(\opn{Q}, \opn{Id})}^{\mrm{isom}}
} 
\end{equation}

\noindent By definition the arrows marked ``f.f.\ emb'' are fully faithful 
embeddings. According to Lemma \ref{lem:2336}  the arrows 
$\cat{Fun}(\opn{Q}, \opn{Id})$ and $\cat{Fun}(\opn{Q}', \opn{Id})$
are isomorphisms of categories. By Lemma \ref{lem:2335} the 
arrow $\cat{Fun}(V, \opn{Id})$ is an equivalence. 
It follows that the arrow $\cat{Fun}(U, \opn{Id})$ in the middle row is an 
equivalence too. 

We know that $\cat{S} \sub \cat{K}$ is a left denominator set.
Condition (b) of the theorem says that $F$ sends morphisms in 
$\cat{S}'$ to isomorphisms in $\cat{E}$. 
Condition (a) there says that there are enough right 
$\cat{K}'$-resolutions in $\cat{K}$.
Thus we are in a position to use the abstract Theorem \ref{thm:2080}.
It says that there is an abstract right derived functor 
$(\mrm{R} F, \eta^{\mrm{R}}) : \cat{D} = \cat{K}_{\cat{S}} \to \cat{E}$ 
of $F$ with respect to $\cat{S}$. 
However, going over the proof of Theorem \ref{thm:2080},
we see that all constructions there can be made within the triangulated 
setting, namely in diagram (\ref{eqn:2346}) instead of in  diagram 
(\ref{eqn:2084}). 
Therefore $\mrm{R} F$ is an object of the category in the 
bottom right corner of (\ref{eqn:2346}), and the morphism 
$\eta^{\mrm{R}} : F \twoto \mrm{R} F \circ \opn{Q}$
is in the category in the top right corner of (\ref{eqn:2346}).
The triangulated variant of step 4 in the proof of Theorem \ref{thm:2080} 
shows that $(\mrm{R} F, \eta^{\mrm{R}})$ satisfies condition (R) of Definition 
\ref{dfn:131}.
\end{proof}

By slight abuse of notation, in the situation of Definition \ref{dfn:131} 
we sometimes refer to the triangulated right derived functor $\mrm{R} F$
just as ``the right derived functor of $F$''. As the next corollary shows, this 
does really cause a problem. 

\begin{cor} \label{cor:3225}
In the situation of Theorem \tup{\ref{thm:1470}}, the triangulated right 
derived functor $(\mrm{R} F, \eta^{\mrm{R}})$ is also an abstract right derived 
functor of $F$ with respect to $\cat{S}$, in the sense of Definition 
\tup{\ref{dfn:2080}}.
\end{cor}

\begin{proof}
This is seen in the proof of Theorem \tup{\ref{thm:1470}}.
\end{proof}

In the situation of Theorem \ref{thm:1470}, let $\cat{K}^{\star}$ be a full 
triangulated subcategory of $\cat{K}$. Define 
$\cat{S}^{\star} := \cat{K}^{\star} \cap \cat{S}$ and 
$\cat{J}^{\star} := \cat{K}^{\star} \cap \cat{J}$. 
Denote by $W : \cat{K}^{\star} \to \cat{K}$
the inclusion functor, and by 
$W_{\cat{S}^{\star}} : \cat{K}^{\star}_{\cat{S}^{\star}} \to \cat{K}_{\cat{S}}$
its localization. 
Warning: the functor $W_{\cat{S}^{\star}}$ is not necessarily fully faithful; 
cf.\ Proposition \ref{prop:1320}. 

\begin{prop} \label{prop:1355}
Assume that every object $M \in \cat{K}^{\star}$ admits a morphism 
$\rho : M \to I$ in $\cat{S}^{\star}$ with target $I \in \cat{J}^{\star}$. 
Then the pair 
$(\mrm{R}^{\star} F, \eta^{\star}) := 
(\mrm{R} F \circ W_{\cat{S}^{\star}}, \eta^{\mrm{R}} \circ \opn{id}_W)$
is a right derived functor of 
$F^{\star} := F \circ W : \cat{K}^{\star} \to \cat{E}$.
\end{prop}

\begin{exer} \label{exer:1355}
Prove the last proposition.  (Hint: Start by choosing a system of right 
$\cat{J}^{\star}$-resolutions in $\cat{K}^{\star}$. Then extend it 
to a system of right $\cat{J}^{}$-resolutions in $\cat{K}^{}$. 
Now follow the proofs of Theorems \ref{thm:1470} and \ref{thm:2080}.)
\end{exer}

Here is the specialization of Definition \ref{dfn:131} to categories of DG 
modules. First:  

\begin{setup} \label{set:3236}
We are given a DG ring $A$, an abelian category
$\cat{M}$, and a full triangulated subcategory
$\cat{K} \sub \dcat{K}(A, \cat{M})$. We write 
$\cat{S} := \cat{K} \cap \, \dcat{S}(A, \cat{M})$, 
the set of quasi-isomorphisms in $\cat{K}$, and 
$\cat{D} := \cat{K}_{\cat{S}}$, the derived category of $\cat{K}$. 
\end{setup}

\begin{dfn} \label{dfn:3236}
Under Setup \ref{set:3236}, let
$\cat{E}$ be a triangulated category, and let 
$F : \cat{K} \to \cat{E}$ 
be a triangulated functor.  
A {\em triangulated right derived functor} 
\index{Derived functor! triangulated right}
\index{Triangulated functor! right derived}
of $F$ is a triangulated right derived functor 
$(\mrm{R} F, \eta^{\mrm{R}}) : \cat{D} \to \cat{E}$
of $F$ with respect to $\cat{S}$, in the sense 
of Definition \ref{dfn:131}. 
\end{dfn}

\begin{exa} \label{exa:1305}
Suppose we start from an additive functor 
$F^0 : \cat{M} \to \cat{N}$
between abelian categories. We know how to extend it to a DG functor 
\[ \bcat{C}^{+}(F^0) : \bcat{C}^{+}(\cat{M}) \to \bcat{C}^{+}(\cat{N}) , \]
and this induces a triangulated functor 
\[ \bcat{K}^{+}(F^0) : \bcat{K}^{+}(\cat{M}) \to  \bcat{K}^{+}(\cat{N}) . \]
By composing with the localization functor 
$\opn{Q}^+_{\cat{N}} : \bcat{K}^{+}(\cat{N}) \to \bcat{D}^{+}(\cat{N})$
we get a triangulated functor 
\begin{equation} \label{eqn:5030}
F := \opn{Q}^+_{\cat{N}} \circ \, \bcat{K}^{+}(F^0) : \bcat{K}^{+}(\cat{M}) 
\to \bcat{D}^{+}(\cat{N}) .
\end{equation}
Define $\cat{K} := \bcat{K}^{+}(\cat{M})$,
$\cat{S} := \bcat{S}^{+}(\cat{M})$ and  
$\cat{E} := \bcat{D}^{+}(\cat{N})$, so in this notation we have a triangulated 
functor $F : \cat{K} \to \cat{E}$, and we are in the situation of Definition 
\ref{dfn:3236}. Note that the restriction of $F$ to the full subcategory 
$\cat{M} \sub \bcat{K}^{+}(\cat{M})$ coincides with the original functor $F^0$. 

Assume that the abelian category $\cat{M}$ has enough injectives; namely 
that every object $M \in \cat{M}$ admits a monomorphism to an injective object. 
Let $\cat{J}$ be the full subcategory of $\cat{K}$ on the bounded below 
complexes of injective objects. 
We will prove later (see Corollary \ref{cor:1922} and Theorem \ref{thm:2880}) 
that properties (a) and (b) of Theorem \ref{thm:1470} hold 
for this $\cat{J}$, regardless what $F$ is. 
Therefore the triangulated right derived functor 
\begin{equation} \label{eqn:5031}
(\mrm{R}^+ F, \eta^+) : \bcat{D}^{+}(\cat{M}) \to \bcat{D}^{+}(\cat{N})
\end{equation}
exists, and for every complex $I \in \cat{J}$ the morphism 
\begin{equation} \label{eqn:2347}
\eta^+_I : F(I) \to \mrm{R}^+ F(I) 
\end{equation}
in $\bcat{D}^{+}(\cat{N})$ is an isomorphism. 

Now assume the original functor $F^0 : \cat{M} \to \cat{N}$ is left exact.
For each $q \geq 0$ we have the classical right derived functor
$\mrm{R}^q F^0 : \cat{M} \to \cat{N}$, and $\mrm{R}^0 F^0 \cong F^0$. 
Formula (\ref{eqn:2347}) shows that for every $M \in \cat{M}$ and $q \geq 0$ 
there is an isomorphism
$\mrm{R}^q F^0(M) \cong \opn{H}^q(\mrm{R}^+ F(M))$  
in $\cat{N}$. As the object $M \in \cat{M}$ moves,
this becomes an isomorphism 
$\mrm{R}^q F^0 \cong \opn{H}^q \circ \, \mrm{R}^+ F$
of additive functors $\cat{M} \to \cat{N}$.

In this example we were careful to use distinct notations for the functor 
$F^0$ between abelian categories, and the induced triangulated functor 
$F = \opn{Q}^+_{\cat{N}} \circ \, \bcat{K}^{+}(F^0)$.
Later we will use the same notation for $F^0$ and $F$. 
\end{exa}

Our treatment of triangulated left derived functors will be 
brief: we will state the definitions and the main results, but won't give 
proofs, beyond a hint here and there on the passage from right to left derived 
functors. 

\begin{dfn}  \label{dfn:1360}
Let $F : \cat{K} \to \cat{E}$ be triangulated functor between  
triangulated categories, and let 
$\cat{S} \sub \cat{K}$ be a denominator set of cohomological origin.    
A {\em triangulated left derived functor of $F$ with respect to $\cat{S}$} 
\index{Derived functor! triangulated left}
\index{Triangulated functor! left derived}
is a triangulated functor 
$\mrm{L}^{} F : \cat{K}_{\cat{S}} \to \cat{E}$,
together with a morphism 
$\eta^{\mrm{L}} : \mrm{L}^{} F \circ \opn{Q} \twoto F$ 
of triangulated functors $\cat{K} \to \cat{E}$. The pair 
$(\mrm{L} F, \eta^{\mrm{L}})$ must have this universal property:
\begin{itemize}
\item[(L)] Given any pair $(G, \th)$, consisting of a triangulated 
functor $G : \cat{K}_{\cat{S}} \to \cat{E}$ and a morphism of triangulated 
functors $\th : G \circ \opn{Q} \twoto F$, 
there is a unique morphism of triangulated functors 
$\mu : G \twoto \mrm{L}^{} F$
such that $\th = \eta^{\mrm{L}} * (\mu \circ \opn{id}_{\opn{Q}})$.
\end{itemize}
\end{dfn}

\begin{prop} \label{prop:1360}
If a triangulated left derived functor $(\mrm{L} F, \eta^{\mrm{L}})$ exists, 
then it is unique, up to a unique isomorphism. Namely, if  $(G, \th)$ is 
another triangulated left derived functor of $F$, then there is a unique 
isomorphism of triangulated functors $\mu : G \twoiso \mrm{L} F$ such that 
$\th = \eta^{\mrm{L}} * (\mu \circ \opn{id}_{\opn{Q}})$. 
\end{prop}

The proof is the same as that of Proposition \ref{prop:1350}, with direction of 
arrows in $\cat{E}$ reversed. 

\begin{thm} \label{thm:1500}
In the situation of Definition \tup{\ref{dfn:1360}}, 
assume there is a full triangulated subcategory 
$\cat{P} \subseteq \cat{K}$
with these two properties\tup{:}
\begin{enumerate}
\rmitem{a} Every object $M \in \cat{K}$ admits a morphism 
$\rho : P \to M$ in $\cat{S}$ with source $P \in \cat{P}$.

\rmitem{b} If $\psi$ is a morphism in $\cat{P} \cap \cat{S}$, 
then $F(\psi)$ is an isomorphism in $\cat{E}$.
\end{enumerate}
Then the triangulated left derived functor
$(\mrm{L} F, \eta^{\mrm{L}}) : \cat{K}_{\cat{S}} \to \cat{E}$ 
\index{Derived functor! triangulated left}
\index{Triangulated functor! left derived}
of $F$ with respect to $\cat{S}$ exists. 
Moreover, for every object $P \in \cat{P}$ the morphism  
$\eta^{\mrm{L}}_P : \mrm{L} F(P) \to F(P)$
in $\cat{E}$ is an isomorphism.
\end{thm}

The proof is the same as that of Theorem \ref{thm:1470}, with direction of 
arrows in $\cat{E}$ reversed.

\begin{cor} \label{cor:3226}
In the situation of Theorem \tup{\ref{thm:1500}}, the triangulated left derived 
functor $(\mrm{L} F, \eta^{\mrm{L}})$ is also an abstract left derived functor 
of $F$ with respect to $\cat{S}$, in the sense of Definition 
\tup{\ref{dfn:2090}}.
\end{cor}

The proof is the same as that of Corollary \ref{cor:3225}, with direction of 
arrows in $\cat{E}$ reversed.

In the situation of Theorem \ref{thm:1500}, let $\cat{K}^{\star}$ be a full 
triangulated subcategory of $\cat{K}$. Define 
$\cat{S}^{\star} := \cat{K}^{\star} \cap \cat{S}$ and 
$\cat{P}^{\star} := \cat{K}^{\star} \cap \cat{P}$. 
Denote by $W : \cat{K}^{\star} \to \cat{K}$
the inclusion functor, and by 
$W_{\cat{S}^{\star}} : \cat{K}^{\star}_{\cat{S}^{\star}} \to \cat{K}_{\cat{S}}$
its localization. 

\begin{prop} \label{prop:1502}
Assume that every $M \in \cat{K}^{\star}$ admits a morphism 
$\rho : P \to M$ in $\cat{S}^{\star}$ with source $P \in \cat{P}^{\star}$. 
Then the pair 
$(\mrm{L}^{\star}, \eta^{\star}) := 
(\mrm{L} F \circ W_{\cat{S}^{\star}}, \eta^{\mrm{L}} \circ \opn{id}_W)$ 
is a triangulated left derived functor of 
$F^{\star} := F \circ W : \cat{K}^{\star} \to \cat{E}$.
\end{prop}

The proof is just like that of Proposition \ref{prop:1355} (which was an 
exercise).

Here is the specialization of Definition \ref{dfn:1360} to categories of DG 
modules.

\begin{dfn} \label{dfn:3237}
Under Setup \ref{set:3236}, let
$\cat{E}$ be a triangulated category, and let 
$F : \cat{K} \to \cat{E}$ be a triangulated functor.  
A {\em triangulated left derived functor}
\index{Derived functor! triangulated left}
\index{Triangulated functor! left derived}
of $F$ is a triangulated left derived functor 
$(\mrm{L} F, \eta^{\mrm{L}}) : \cat{D} \to \cat{E}$
of $F$ with respect to $\cat{S}$, in the sense 
of Definition \ref{dfn:1360}. 
\end{dfn}

\begin{exa} \label{exa:1501}
This is analogous to Example \ref{exa:1305}.
We start with an additive functor 
$F^0 : \cat{M} \to \cat{N}$ between abelian categories. We know how to extend 
it 
to a DG functor 
\[ \bcat{C}^{-}(F^0) : \bcat{C}^{-}(\cat{M}) \to \bcat{C}^{-}(\cat{N}) , \]
and this induces a triangulated functor 
\[ \bcat{K}^{-}(F^0) : \bcat{K}^{-}(\cat{M}) \to  \bcat{K}^{-}(\cat{N}) . \]
By composing with $\opn{Q}^-_{\cat{N}}$ we get a triangulated functor
\begin{equation} \label{eqn:5032}
F := \opn{Q}^-_{\cat{N}} \circ \, \bcat{K}^{-}(F^0) : \bcat{K}^{-}(\cat{M}) 
\to \bcat{D}^{-}(\cat{N}) . 
\end{equation}
Defining $\cat{K} := \bcat{K}^{-}(\cat{M})$,
$\cat{S} := \bcat{S}^{+}(\cat{M})$ and 
$\cat{E} := \bcat{D}^{-}(\cat{N})$,
we have a triangulated functor 
$F : \cat{K} \to \cat{E}$, and the situation is that of Definition 
\ref{dfn:3237}.

Assume that the abelian category $\cat{M}$ has enough projectives. 
Define $\cat{P}$ to be the full subcategory of $\cat{K}$ on the bounded above 
complexes of projective objects. 
We will prove later (in Corollary \ref{cor:1923}  and Theorem \ref{thm:1540}) 
that properties (a) and (b) of Theorem \ref{thm:1500} hold 
in this situation. Therefore we have a left derived functor
\begin{equation} \label{eqn:5033}
(\mrm{L}^- F, \eta^-) : \bcat{D}^{-}(\cat{M}) \to \bcat{D}^{-}(\cat{N}) . 
\end{equation}
For any $P \in \cat{P}$ the morphism 
\begin{equation} \label{eqn:2351}
\eta^-_P : \mrm{L}^- F(P) \to F(P) 
\end{equation}
in $\cat{E}$ is an isomorphism. 

Now assume the functor $F^0$ is right exact. Then $F^0$ has the classical 
left derived functors 
$\mrm{L}_q F^0 : \cat{M} \to \cat{N}$
for all $q \geq 0$, and $\mrm{L}_0 F^0 \cong F^0$. Formula (\ref{eqn:2351}) 
shows that for every $M \in \cat{M}$ there is an isomorphism 
$\mrm{L}_q F^0(M) \cong \opn{H}^{-q}(\mrm{L}^- F(M))$
in $\cat{N}$. As $M$ moves there is an isomorphism  
$\mrm{L}_q F^0 \cong \opn{H}^{-q} \circ \, \mrm{L}^- F$
of additive functors $\cat{M} \to \cat{N}$.

In this example, as in Example \ref{exer:1355}, we were careful to use 
distinct notations for the functor 
$F^0$ between abelian categories, and the induced triangulated functor 
$F = \opn{Q}^-_{\cat{N}} \circ \, \bcat{K}^{-}(F^0)$.
Later we will use the same notation for $F^0$ and $F$. 
\end{exa}

\begin{rem} \label{rem:4695}
The reader might ask the following question at this point: 
Suppose we are given additive functors
\begin{equation} \label{eqn:4705}
\cat{L} \xar{F} \cat{M} \xar{G} \cat{N}
\end{equation}
between abelian categories, and the categories 
$\cat{M}$ and $\cat{N}$ have enough injectives, so that the right derived 
functors 
\begin{equation} \label{eqn:4706}
\UseTips \xymatrix @C=8ex @R=8ex { 
\dcat{D}^+(\cat{L})
\ar[r]^{\mrm{R}^+ F}
\ar@(ur,ul)[rr]^{\mrm{R}^+ (G \circ F)}
&
\dcat{D}^+(\cat{M})
\ar[r]^{\mrm{R}^+ G}
&
\dcat{D}^+(\cat{N})
}
\end{equation}
exist (see Example \ref{exa:1305} above). Does this diagram commute (up to 
an isomorphism of triangulated functors)? Similar questions can be asked about 
left derived functors, etc. 

Here is an answer.  Assume that $F$ sends injectives in 
$\cat{L}$ to {\em right $G$-acyclic objects} of $\cat{M}$. See Definition
\ref{dfn:3887} (that talks about categories of graded modules, but the ideas 
are the same). By Lemma \ref{lem:3885} we know that for each bounded 
below complex of right $G$-acyclic objects $J$, the canonical morphism 
$G(J) \to \mrm{R}^+ G(J)$ in $\dcat{D}^+(\cat{M})$ is an isomorphism. 
Therefore, for every bounded below complex of injectives
$I \in \dcat{C}(\cat{L})$,
letting $J := F(I)$,  we have $J \cong \mrm{R}^+ F(I)$, and hence 
\[ \mrm{R}^+ (G \circ F)(I) \cong (G \circ F)(I) = G(F(I)) = G(J) \cong 
\mrm{R}^+ G(J) \cong \mrm{R}^+ G (\mrm{R}^+ F(I)) . \]
Thus there is an isomorphism of triangulated functors  
\begin{equation} \label{eqn:4707}
\mrm{R}^+ (G \circ F) \cong \mrm{R}^+ G \circ \mrm{R}^+  F : 
\dcat{D}^+(\cat{L}) \to \dcat{D}^+(\cat{N}) .
\end{equation}

The original application of the isomorphism (\ref{eqn:4707}),
in \cite[Proposition II.5.1]{RD}, was this: given morphisms of schemes 
$X \xar{f} Y \xar{g} Z$, the derived pushforward functors satisfy
\begin{equation} \label{eqn:5034}
\mrm{R}^+ (g_* \circ f_*) \cong \mrm{R}^+ g_* \circ \mrm{R}^+ f_*
\end{equation}
as functors $\dcat{D}^+(X) \to \dcat{D}^+(Z)$. 
This is true because the functor $f_*$ sends injective $\OO_X$-modules to 
flasque $\OO_Y$-modules. 

This pleasant phenomenon seems to fail in some important situations 
that we study in our book. For instance, take a noncommutative ring $A$ and 
$A$-bimodules $S, T$ that are tilting complexes but not invertible 
bimodules (see Subsection \ref{subsec:tilting-rings}).  
Consider the abelian categories 
$\cat{L} = \cat{M} = \cat{N} := \dcat{M}(A)$
and the functors $F := \opn{Hom}_A(T, -)$ and $G := \opn{Hom}_A(S, -)$. 
It is quite likely that the functors $\mrm{R}^+G \circ \mrm{R}^+F$ and 
$\mrm{R}^+(G \circ F)$ are not isomorphic. 
A concrete instance in which this can be tested is the ring $A$ 
from Example \ref{exa:4870}, and the bimodules $S = T := A^*$ there.
\end{rem}

\begin{rem} \label{rem:4705}
The composition of derived functors explained in the previous remark is also 
related to {\em spectral sequences}.
\index{Spectral sequence}
Given the input (\ref{eqn:4705}), in classical homological algebra it was 
impossible to ask for something like the commutative diagram (\ref{eqn:4706}).
The best approximation was to have a {\em convergent spectral sequence}
\begin{equation} \label{eqn:4695}
E^{i, j}_k = \mrm{R}^j G \circ \mrm{R}^i F \Rightarrow
\mrm{R}^n (G \circ F) . 
\end{equation}
Indeed, under the assumption in the previous remark (namely that
$F$ sends injective objects to right $G$-acyclic objects), such a convergent 
spectral sequence does exist; see   
\cite[Chapters XI-XII]{Mac1}, \cite[Section 11]{Rot} or 
\cite[Section 5.7]{We}. In the geometric example above, for which the 
isomorphism (\ref{eqn:5034}) holds, this is the {\em Leray spectral 
sequence}.  

There are several kinds of spectral sequences, but basically they are all of the 
same nature, i.e.\ approximations of the composition of derived functors. 

Most traditional applications of spectral sequences are performed more 
effectively by filtrations and truncations of complexes, as can be seen in 
Section \ref{sec:adj-equ-cohdim} of this book.
Still, there are a few instances in which spectral sequences 
cannot be left out, such as the {\em niveau spectral sequence} in algebraic 
geometry (see \cite[Chapter IV]{RD}), or its noncommutative relative, the {\em 
double Ext spectral sequence} for Auslander dualizing complexes 
(see Remark \ref{rem:3821} and the paper \cite{YeZh1}).  
\end{rem}



\mysubsection{Contravariant Triangulated Derived Functors} 
\label{subsec:cntr-tri-der-funs}

In this subsection there is a fixed nonzero commutative base ring $\K$, and 
all categories and functors here are $\K$-linear. The next setup is assumed. 

\begin{setup} \label{set:3225}
We are given a DG ring $A$, an abelian category
$\cat{M}$, and a full additive subcategory
$\cat{K} \sub \dcat{K}(A, \cat{M})$ s.t.\ 
$\cat{K}^{\mrm{op}} \sub \dcat{K}(A, \cat{M})^{\mrm{op}}$
is triangulated. We write 
$\cat{S} := \cat{K} \cap \, \dcat{S}(A, \cat{M})$, 
the set of quasi-isomorphisms in $\cat{K}$, and 
$\cat{D}^{\mrm{op}} := (\cat{K}^{\mrm{op}})_{\cat{S}^{\mrm{op}}}$,
the derived category of $\cat{K}^{\mrm{op}}$. 
\end{setup}

As explained in Subsections \ref{subsec:opp-hom-triang} and 
\ref{subsec:opp-dercat-triang}, the categories 
$\cat{K}^{\mrm{op}}$ and $\cat{D}^{\mrm{op}}$ 
have canonical triangulated structures, and the localization functor 
$\opn{Q}^{\mrm{op}} : \cat{K}^{\mrm{op}} \to \cat{D}^{\mrm{op}}$
is triangulated. 

\begin{dfn} \label{dfn:2356}
\index{Derived functor! contravariant triangulated right}
\index{Triangulated functor! contravariant right derived}
Let $\cat{E}$ be a triangulated category, and let 
$F : \cat{K}^{\mrm{op}} \to \cat{E}$
be a triangulated functor.  
A {\em triangulated right derived functor} of $F$ is a triangulated 
right derived functor 
$(\mrm{R} F, \eta^{\mrm{R}}) : \cat{D}^{\mrm{op}} \to \cat{E}$
of $F$ with respect to $\cat{S}^{\mrm{op}}$, in the sense 
of Definition \ref{dfn:131}. 
\end{dfn}

Likewise:

\begin{dfn} \label{dfn:2357} 
\index{Derived functor! contravariant triangulated left}
\index{Triangulated functor! contravariant left derived}
Let $\cat{E}$ be a triangulated category, and let 
$F : \cat{K}^{\mrm{op}} \to \cat{E}$ 
be a triangulated functor. A {\em triangulated left derived functor} of $F$ is 
a triangulated left derived functor 
$(\mrm{L} F, \eta^{\mrm{L}}) : \cat{D}^{\mrm{op}} \to \cat{E}$
of $F$ with respect to $\cat{S}^{\mrm{op}}$, in the sense 
of Definition \ref{dfn:1360}. 
\end{dfn}

The uniqueness results (Propositions \ref{prop:1350} and \ref{prop:1360})
and existence results (Theorems \ref{thm:1470} and \ref{thm:1500}) 
apply in the contravariant case as well. We shall return to the existence 
question the end of Section \ref{sec:resol}, when we talk about {\em 
resolving subcategories}; this will make matters more concrete. 
For now we give one example. 

\begin{exa} \label{exa:3055}
Let $\cat{M}$ and $\cat{N}$ be abelian categories, and let 
$F^0 : \cat{M} \to \cat{N}$
be a contravariant additive functor. This is the same as an additive functor
$F^0 \circ \opn{Op} : \cat{M}^{\mrm{op}} \to \cat{N}$.
As in Theorem \ref{thm:3015}, $F^0$ gives 
rise to a contravariant DG functor 
\[ \dcat{C}^-(F^0) : \dcat{C}^-(\cat{M}) \to \dcat{C}^+(\cat{N}) . \]
This means, according to Proposition \ref{prop:2500}, that we have a  DG 
functor 
\[ \dcat{C}^-(F^0) \circ \opn{Op} : \dcat{C}^-(\cat{M})^{\mrm{op}} \to 
\dcat{C}^+(\cat{N}) . \]
Passing to homotopy categories, and using Theorem \ref{thm:2995}, we obtain a 
triangulated functor 
\[ \opn{Ho} \bigl( \dcat{C}^-(F^0) \circ \opn{Op} \bigr) :
\dcat{K}^-(\cat{M})^{\mrm{op}} \to \dcat{K}^+(\cat{N}) \]
(with a translation isomorphism $\bar{\tau}$ that we ignore 
now). Composing this with the localization functor
$\opn{Q}_{\cat{N}}$ we have a triangulated functor 
\begin{equation} \label{eqn:5035}
F := \opn{Q}_{\cat{N}} \circ \, 
\opn{Ho} \bigl( \dcat{C}^-(F^0) \circ \opn{Op} \bigr) : 
\dcat{K}^-(\cat{M})^{\mrm{op}} \to 
\dcat{D}^+(\cat{N}) .
\end{equation}
Similarly there is a triangulated functor 
\begin{equation} \label{eqn:5036}
F^{\mrm{flip}} := \opn{Q}_{\cat{N}} \circ \, 
\opn{Ho} \bigl( \dcat{C}^+(F^0 \circ \opn{Op}) \bigr) :
\dcat{K}^+(\cat{M}^{\mrm{op}}) \to \dcat{D}^+(\cat{N}) .
\end{equation}
Observe that this triangulated functor is of the same sort as the one we had in 
(\ref{eqn:5030}). The functors (\ref{eqn:5035}) and (\ref{eqn:5036}) 
fit into the the following commutative diagram of triangulated functors 
\[ \UseTips \xymatrix @C=10ex @R=6ex {
\dcat{K}^-(\cat{M})^{\mrm{op}}
\ar[r]^{F}
\ar[d]_{\ol{\opn{Flip}}}^{\cong}
&
\dcat{D}^+(\cat{N})
\\
\dcat{K}^+(\cat{M}^{\mrm{op}}) 
\ar[ur]_{F^{\mrm{flip}}}
} \]
where $\ol{\opn{Flip}}$ is the isomorphism of triangulated categories from 
(\ref{eqn:2995}). 

The triangulated right derived functor 
\begin{equation} \label{eqn:5037}
\mrm{R}^- F : \dcat{D}^-(\cat{M})^{\mrm{op}} \to \dcat{D}^+(\cat{N})
\end{equation}
of the functor $F$ from (\ref{eqn:5035}) would sit inside this diagram of 
triangulated functors, which is 
commutative up to a unique isomorphism: 
\[ \UseTips \xymatrix @C=10ex @R=6ex {
\dcat{D}^-(\cat{M})^{\mrm{op}}
\ar[r]^{\mrm{R}^-F}
\ar[d]_{\ol{\opn{Flip}}}^{\cong}
&
\dcat{D}^+(\cat{N})
\\
\dcat{D}^+(\cat{M}^{\mrm{op}}) 
\ar[ur]_{\mrm{R}^+ F^{\mrm{flip}}}
} \]

Like in Example \ref{exa:1305}, to show existence of $\mrm{R}^+ F^{\mrm{flip}}$ 
using Theorem \ref{thm:1470},
and to calculate it, we would need a full triangulated subcategory 
$\cat{J} \sub \dcat{K}^+(\cat{M}^{\mrm{op}})$, 
such that every $M \in \dcat{K}^+(\cat{M}^{\mrm{op}})$
admits a quasi-isomorphism 
$M \to I$ with $I \in \cat{J}$; and 
every quasi-isomorphism $\psi$ in $\cat{J}$ goes to an isomorphism 
$F^{\mrm{flip}}(\psi)$ in $\dcat{D}^+(\cat{N})$.
This translates, by the flip, to a full triangulated subcategory
$\cat{P} \sub \dcat{K}^-(\cat{M})$ such that 
$\cat{J} = \ol{\opn{Flip}}(\cat{P})$. 
The conditions on $\cat{P}$ are these: 
each $M \in \dcat{K}^-(\cat{M})$
admits a quasi-isomorphism 
$P \to M$ with $P \in \cat{P}$; and each quasi-isomorphism $\phi$ 
in $\cat{P}$ goes to an isomorphism
$F(\phi)$ in $\dcat{D}^+(\cat{N})$.
For instance, if $\cat{M}$ has enough projectives, then according to Corollary 
\ref{cor:1900}, the category $\cat{P}$ of bounded above complexes of 
projectives satisfies these conditions. 

To calculate the derived functor $\mrm{R}^- F$ in equation (\ref{eqn:5037}) we 
choose a system of left $\cat{P}$-resolutions 
in $\dcat{K}^-(\cat{M})$, as in Definition \ref{dfn:2095}. 
Thus for every $M$ we have a chosen quasi-isomorphism 
$\rho_M : P(M) \to M$ in $\dcat{K}^{-}(\cat{M})$. We take
$\mrm{R}^- F(M) := F(P(M)) \in \dcat{D}^+(\cat{N})$.
The morphism
$\eta^{-}_M : F(M) \to \mrm{R}^- F(M)$
is 
$\eta^{-}_M := F(\rho_M) : F(M) \to F(P(M))$.

The story is very similar for the triangulated left derived functor 
\[ \mrm{L}^+ F : \dcat{D}^+(\cat{M})^{\mrm{op}} \to \dcat{D}^-(\cat{N})  \]
of 
\[ F := \opn{Q}_{\cat{N}} \circ \, 
\opn{Ho} \bigl( \dcat{C}^+(F^0) \circ \opn{Op} \bigr) : 
\dcat{K}^+(\cat{M})^{\mrm{op}} \to 
\dcat{D}^-(\cat{N}) , \]
so we won't tell it. 
\end{exa}

In later parts of the book we shall use the same symbol $F$ to denote also the 
functors $F^0$, $F^0 \circ \opn{Op}$ and  
$\dcat{C}^{\star}(F^0) \circ \opn{Op}$
from the example above. This will simplify our notation. 

Though indirectly defined (in Definition \ref{dfn:2995}), the  translation 
functor $\opn{T}^{\mrm{op}}$
of $\dcat{K}(A, \cat{M})^{\mrm{op}}$
is known to have some useful properties. Note that the objects of 
$\dcat{K}(A, \cat{M})^{\mrm{op}}$ are the same as those of 
$\dcat{K}(A, \cat{M})$. 

\begin{prop} \label{prop:2307}
For every object $M \in \dcat{K}(A, \cat{M})^{\mrm{op}}$ there are canonical 
isomorphisms 
$\opn{T}^{\mrm{op}}(M)^i \cong M^{i - 1}$
and 
$\opn{H}^i(\opn{T}^{\mrm{op}}(M)) \cong \opn{H}^{i - 1}(M)$
in $\cat{M}$. 
\end{prop}

\begin{proof}
This is immediate from Theorem \ref{thm:2495} and Definition \ref{dfn:2995}.
\end{proof}

Here are several useful properties of the triangulated structure on 
$\dcat{D}(A, \cat{M})^{\mrm{op}}$.

\begin{prop} \label{prop:3060}
Suppose 
\[ 0 \to L \xar{\phi} M \xar{\psi} N \to 0 \]
is an exact sequence in 
$\dcat{C}_{\mrm{str}}(A, \cat{M})^{\mrm{op}}$. 
Then there is a distinguished triangle 
\[ L \xar{\opn{Q}^{\mrm{op}}(\phi)} 
M \xar{\opn{Q}^{\mrm{op}}(\psi)} 
N \xar{\th} \opn{T}^{\mrm{op}}(L) \]
in $\dcat{D}(A, \cat{M})^{\mrm{op}}$.
\end{prop}

\begin{proof}
Apply the functor $\opn{Flip}$ to the given exact sequence. By Theorem 
\ref{thm:2495}(3) we obtain a short exact sequence 
\[ 0 \to \opn{Flip}(L) \xar{\opn{Flip}(\phi)} \opn{Flip}(M) 
\xar{\opn{Flip}(\psi)} \opn{Flip}(N) \to 0 \]
in 
$\dcat{C}_{\mrm{str}}(A^{\mrm{op}}, \cat{M}^{\mrm{op}})$. 
According to Proposition \ref{prop:2165} there is a  distinguished triangle
\[ \opn{Flip}(L) \xar{\opn{Q}^{\mrm{flip}}(\opn{Flip}(\phi))}
\opn{Flip}(M) \xar{\opn{Q}^{\mrm{flip}}(\opn{Flip}(\psi))}
\opn{Flip}(N) \xar{\th'}
\opn{T}^{\mrm{flip}}(\opn{Flip}(L)) \]
in 
$\dcat{D}(A^{\mrm{op}}, \cat{M}^{\mrm{op}})$.
Take the morphism 
\[ \th := \opn{Flip}^{-1}(\th') : N \to 
\opn{Flip}^{-1}(\opn{T}^{\mrm{flip}}(\opn{Flip}(L))) = 
\opn{T}^{\mrm{op}}(L) \]
in $\dcat{D}(A, \cat{M})^{\mrm{op}}$.
Since 
$\opn{Flip}^{-1}(\opn{Q}^{\mrm{flip}}(\opn{Flip}(\phi))) = 
\opn{Q}^{\mrm{op}}(\phi)$,
and likewise for $\psi$, we have the desired distinguished triangle. 
\end{proof}

Recall from Proposition \ref{prop:3170}
that if $A$ is a nonpositive DG ring, then for any 
$M \in \dcat{C}(A, \cat{M})$ and integer $i$ the smart truncations 
$\opn{smt}^{\geq i}(M), \, \opn{smt}^{\leq i}(M) \in \dcat{C}(A, \cat{M})$
exist. 

\begin{prop} \label{prop:3061}
Assume $A$ is a nonpositive DG ring. 
For every object $M \in \dcat{C}(A, \cat{M})$
there is a distinguished triangle 
\[ \opn{smt}^{\geq i + 1}(M) \to M \to 
\opn{smt}^{\leq i}(M) \to 
\opn{T}^{\mrm{op}}(\opn{smt}^{\geq i + 1}(M)) \]
in $\dcat{D}(A, \cat{M})^{\mrm{op}}$. The truncations are performed in 
$\dcat{C}(A, \cat{M})$ as above. 
\end{prop}

\begin{proof}
Let 
$N := \opn{Flip}(M) \in \dcat{D}(A^{\mrm{op}}, \cat{M}^{\mrm{op}})$. 
By Proposition \ref{prop:2320} there is a distinguished triangle
\[ \opn{smt}^{\leq i}(N) \xar{} N \xar{} \opn{smt}^{\geq i + 1}(N) 
\xar{} \opn{T}^{\mrm{flip}}(\opn{smt}^{\leq i}(N)) \]
in $\dcat{D}(A^{\mrm{op}}, \cat{M}^{\mrm{op}})$,
where now the truncations are done in 
$\dcat{C}(A^{\mrm{op}}, \cat{M}^{\mrm{op}})$. 
The formulas defining the flip functor in the proof of Theorem \ref{thm:2495}
show that \lb
$\opn{Flip}^{-1}(\opn{smt}^{\leq i}(N)) = \opn{smt}^{\geq i + 1}(M)$
and 
$\opn{Flip}^{-1}(\opn{smt}^{\geq i + 1}(N)) = \opn{smt}^{\leq i}(M)$.
\end{proof}



